\chapter{Finite-Group MPU Representations and Scalar Gauge Data}%
\label{ch:mpu_gauging}

\section{Finite-group unitary averages}

\begin{definition}[Finite-group unitary average]
    \label{def:mpug_finite_group_unitary_average}
    \lean{TNLean.Algebra.finiteGroupUnitaryAverage,
        TNLean.Algebra.unitaryMatrixRepresentation}
    \leanok
    \discussion{7458}
    Let $G$ be a finite group and let $\rho\colon G\to\mathrm U(n)$ be a
    unitary representation. Define
    \begin{align}
        P_\rho & = \frac{1}{|G|}\sum_{g\in G}\rho(g).
    \end{align}
    For a fixed lattice site $j$, the specialization
    $\rho(g)=\mathcal G_g^j$ gives the local Gauss-law average $P_j$ in
    \cite[lines~459--461]{FrancoRubio2025MPUGauging}.
\end{definition}

\begin{theorem}[Finite-group unitary average is an orthogonal projection]
    \label{thm:mpug_finite_group_unitary_average_projection}
    \lean{TNLean.Algebra.isStarProjection_finiteGroupUnitaryAverage,
        TNLean.Algebra.isOrthogonalProjection_finiteGroupUnitaryAverage}
    \leanok
    \discussion{7458}
    \uses{def:mpug_finite_group_unitary_average}
    For every finite-dimensional unitary representation $\rho$ of a finite
    group,
    \begin{align}
        P_\rho^\dagger & =P_\rho, \\
        P_\rho^2       & =P_\rho.
    \end{align}
    Hence $P_\rho$ is an orthogonal projection.
\end{theorem}

\begin{proof}\leanok
    \uses{def:mpug_finite_group_unitary_average}
    Regard each matrix $\rho(g)$ as the corresponding linear map and let
    $e_G=|G|^{-1}\sum_{g\in G}g$ be the normalized group-algebra average. The
    finite-group representation-averaging theorem used here is the projection
    identity whose concrete group-algebra calculation is
    \begin{align}
        h e_G & =e_G \quad (h\in G),                  &
        e_G^2 & =\frac{1}{|G|}\sum_{h\in G}h e_G=e_G.
    \end{align}
    Applying the linear representation and transporting back to matrices gives
    $P_\rho^2=P_\rho$.
    Unitarity gives $\rho(g)^\dagger=\rho(g^{-1})$. Since inversion permutes
    $G$,
    \begin{align}
        P_\rho^\dagger
         & = \frac{1}{|G|}\sum_{g\in G}\rho(g)^\dagger
        = \frac{1}{|G|}\sum_{g\in G}\rho(g^{-1})
        = P_\rho.
    \end{align}
\end{proof}

\section{Finite-group representations by matrix product unitaries}
\label{sec:mpug_group_representations}

Let $G$ be a group, let $d\in\N$, and let
$D\colon G\to\N_{>0}$ assign a virtual bond dimension to each group element.
For every $g\in G$, let $\mathcal U_g$ be an MPO tensor with physical input and
output dimension $d$ and virtual bond dimension $D_g$. Write $U_g^{(N)}$ for
the periodic operator obtained by closing a chain of $N$ copies of
$\mathcal U_g$. Throughout this section the chain length is positive:
$N\geq1$. In particular, no identity at length zero is asserted.

\begin{definition}[Matrix product unitarity on nonempty chains]
    \label{def:mpug_positive_mpu}
    \lean{MPOTensor.IsMPUPos}
    \leanok
    An MPO tensor $\mathcal U$ is a matrix product unitary on nonempty chains
    when $U^{(N)}$ is unitary for every $N\geq1$. This is the system-size
    convention in \cite[line~207]{FrancoRubio2025MPUGauging}.
\end{definition}

\begin{theorem}[Comparison with the $N>1$ MPU convention]
    \label{thm:mpug_positive_mpu_to_mpu}
    \lean{MPOTensor.IsMPUPos.isMPU}
    \leanok
    \uses{def:mpug_positive_mpu, def:mpu}
    Every matrix product unitary on nonempty chains is a matrix product unitary
    in the sense of Definition~\ref{def:mpu}.
\end{theorem}

\begin{proof}\leanok
    \uses{def:mpug_positive_mpu, def:mpu}
    The nonempty-chain condition applies in particular to every $N>1$.
\end{proof}

\begin{theorem}[Physical blocking preserves nonempty-chain unitarity]
    \label{thm:mpug_positive_mpu_blocking}
    \lean{MPOTensor.IsMPUPos.blockTensor}
    \leanok
    \uses{def:mpug_positive_mpu, def:mpdo_physical_blocking}
    If $\mathcal U$ is a matrix product unitary on nonempty chains and $L>0$,
    then $\mathcal U^{[L]}$ is also a matrix product unitary on nonempty chains.
\end{theorem}

\begin{proof}\leanok
    \uses{thm:mpu_reindex_unitary, thm:mpdo_closed_chain_physical_blocking}
    At chain length $N\geq1$, the canonical reindexing of configurations gives
    \begin{align}
        O_N(\mathcal U^{[L]}) & \cong U^{(NL)}.
    \end{align}
    Since $N,L>0$, the length $NL$ is positive. Thus $U^{(NL)}$ is unitary,
    and simultaneous reindexing preserves unitarity.
\end{proof}

\begin{definition}[Group representation by simple injective MPUs]
    \label{def:mpug_group_representation}
    \lean{MPOTensor.GroupFamily,
        MPOTensor.GroupFamily.IsRepresentation}
    \leanok
    \uses{def:injective, def:mpu_simple, def:mpug_positive_mpu}
    A group representation by simple injective MPUs consists of the
    dependent-bond family $\{\mathcal U_g\}_{g\in G}$ above such that every
    $\mathcal U_g$ is injective and simple, every $\mathcal U_g$ is a matrix
    product unitary on nonempty chains, and
    \begin{align}
        U_e^{(N)}          & =\Id,
        \label{eq:mpug_rep_identity}       \\
        U_g^{(N)}U_h^{(N)} & =U_{gh}^{(N)}
        \label{eq:mpug_rep_multiplication}
    \end{align}
    for all $g,h\in G$ and all $N\geq1$.
    The multiplicative law is the representation-level input
    in~\cite[Subsection ``MPU group representations'']
    {FrancoRubio2025MPUGauging}.
    The identity law is its standard consequence and is recorded explicitly for
    positive lengths. It is also compatible with the later, stronger convention
    that $\mathcal U_e=\Id$ has bond dimension one
    in~\cite[Subsection ``Assumptions for the rest of the paper'']
    {FrancoRubio2025MPUGauging}.
    The dimensions $D_gD_h$ and $D_{gh}$ need not agree, so
    (\ref{eq:mpug_rep_multiplication}) is an identity of periodic operators,
    not an equality of local tensors.
\end{definition}

\begin{theorem}[Powers of represented operators]
    \label{thm:mpug_operator_power}
    \lean{MPOTensor.GroupFamily.IsRepresentation.operator_pow}
    \leanok
    \discussion{7445}
    \uses{def:mpug_group_representation}
    For every $g\in G$, every $n\in\N$, and every $N\geq1$,
    \begin{align}
        \bigl(U_g^{(N)}\bigr)^n & =U_{g^n}^{(N)}.
        \label{eq:mpug_operator_power}
    \end{align}
\end{theorem}

\begin{proof}\leanok
    \uses{def:mpug_group_representation}
    Induct on $n$. For $n=0$, the identity law gives
    $U_e^{(N)}=\Id$. For the induction step, multiply the induction hypothesis
    by $U_g^{(N)}$ and apply the representation multiplication law:
    \begin{align}
        \bigl(U_g^{(N)}\bigr)^{n+1}
         & =U_{g^n}^{(N)}U_g^{(N)}
        =U_{g^{n+1}}^{(N)}.
    \end{align}
\end{proof}

\begin{theorem}[Finite order of represented operators]
    \label{thm:mpug_operator_finite_order}
    \lean{MPOTensor.GroupFamily.IsRepresentation.operator_pow_orderOf_eq_one}
    \leanok
    \discussion{7445}
    \uses{def:mpug_group_representation}
    Suppose that $G$ is finite. For every $g\in G$ and every $N\geq1$, the
    represented operator has finite order dividing the order of $g$:
    \begin{align}
        \bigl(U_g^{(N)}\bigr)^{\operatorname{ord}(g)} & =\Id.
    \end{align}
    This is the finite-order assertion in
    \cite[line~1547]{FrancoRubio2025MPUGauging}.
    % Source anchor: Papers/2502.20257/main.tex, line 1547.
\end{theorem}

\begin{proof}\leanok
    \uses{def:mpug_group_representation, thm:mpug_operator_power}
    Apply (\ref{eq:mpug_operator_power}) with $n=\operatorname{ord}(g)$.
    The finite-order identity $g^{\operatorname{ord}(g)}=e$ and the
    representation identity law then give
    $\bigl(U_g^{(N)}\bigr)^{\operatorname{ord}(g)}=U_e^{(N)}=\Id$.
\end{proof}

% Tracking boundary: these entries cover only finite order of the periodic
% operators. Issue #7318 remains open for public-index vanishing, the
% rank-product input r l = d^2, and the final rank equality r = l = d.

\begin{theorem}[Product-tensor operator identity]
    \label{thm:mpug_product_tensor_operator}
    \lean{MPOTensor.GroupFamily.IsRepresentation.sameMPV₂Pos_mulTensor}
    \leanok
    \uses{def:mpug_group_representation, def:mpdo_operator_product}
    Let $\mathcal U_g\mathbin{\cdot}\mathcal U_h$ denote the product MPO tensor,
    whose virtual bond dimension is $D_gD_h$. For every $N\geq1$,
    \begin{align}
        O_N(\mathcal U_g\mathbin{\cdot}\mathcal U_h)
         & =U_g^{(N)}U_h^{(N)}
        =U_{gh}^{(N)}.
        \label{eq:mpug_product_tensor_positive_chain}
    \end{align}
    Thus the product tensor and the selected tensor $\mathcal U_{gh}$ generate
    the same nonempty-chain periodic operators, although their virtual bond
    dimensions may differ.
\end{theorem}

\begin{proof}\leanok
    \uses{def:mpug_group_representation, thm:mpdo_operator_product_mpo}
    The first equality in
    (\ref{eq:mpug_product_tensor_positive_chain}) is multiplicativity of the
    periodic operator under the product-tensor contraction. The second is
    (\ref{eq:mpug_rep_multiplication}).
\end{proof}

\begin{theorem}[The physical adjoint represents the inverse]
    \label{thm:mpug_physical_adjoint_inverse}
    \lean{MPOTensor.GroupFamily.IsRepresentation.mpo_physicalAdjointTensor_eq_inv}
    \leanok
    \uses{def:mpug_group_representation, def:mpdo_physical_adjoint}
    For every $g\in G$ and every $N\geq1$, the physical adjoint of
    $\mathcal U_g$ generates the operator representing $g^{-1}$:
    \begin{align}
        O_N(\mathcal U_g^\sharp)
         & =U_{g^{-1}}^{(N)}.
        \label{eq:mpug_physical_adjoint_inverse}
    \end{align}
    This is the first assertion in
    \cite[line~1552]{FrancoRubio2025MPUGauging}. It requires no finiteness
    assumption on $G$.
\end{theorem}

\begin{proof}\leanok
    \uses{def:mpug_group_representation, def:mpdo_physical_adjoint}
    The representation law and unitarity give
    \begin{align}
        U_{g^{-1}}^{(N)}U_g^{(N)}              & =U_e^{(N)}=\Id, \\
        U_g^{(N)}\bigl(U_g^{(N)}\bigr)^\dagger & =\Id.
    \end{align}
    Hence
    \begin{align}
        U_{g^{-1}}^{(N)}
         & =U_{g^{-1}}^{(N)}
        \bigl(U_g^{(N)}\bigl(U_g^{(N)}\bigr)^\dagger\bigr)
        =\bigl(U_{g^{-1}}^{(N)}U_g^{(N)}\bigr)
        \bigl(U_g^{(N)}\bigr)^\dagger
        =\bigl(U_g^{(N)}\bigr)^\dagger.
    \end{align}
    Closing the physical-adjoint tensor gives the conjugate transpose, which
    proves (\ref{eq:mpug_physical_adjoint_inverse}).
\end{proof}

\begin{theorem}[Positive-length matrix product vector consequence]
    \label{thm:mpug_physical_adjoint_inverse_mpv}
    \lean{MPOTensor.GroupFamily.IsRepresentation.sameMPV₂Pos_physicalAdjointTensor_inv}
    \leanok
    \uses{def:mpug_group_representation, def:mpdo_physical_adjoint,
        def:mpo_to_mps, def:same_mpv2_pos}
    View $\mathcal U_g^\sharp$ and $\mathcal U_{g^{-1}}$ as MPS tensors by
    pairing their two physical indices. They generate the same matrix product
    vectors at every positive length:
    \begin{align}
        \mc{V}^+\bigl((\mathcal U_g^\sharp)^{\mathrm{MPS}}\bigr)
         & =\mc{V}^+\bigl(\mathcal U_{g^{-1}}^{\mathrm{MPS}}\bigr).
        \label{eq:mpug_physical_adjoint_inverse_mpv}
    \end{align}
\end{theorem}

\begin{proof}\leanok
    \uses{thm:mpug_physical_adjoint_inverse,
        lem:mpdo_three_first_site_contraction_coefficients}
    For a doubled physical configuration $\rho$, decode each site as a pair of
    ket and bra indices $(\sigma_n,\tau_n)$. The corresponding matrix product
    vector coefficient is the $(\sigma,\tau)$ matrix entry of the periodic
    operator. Equality
    (\ref{eq:mpug_physical_adjoint_inverse}) therefore gives equality of every
    coefficient in (\ref{eq:mpug_physical_adjoint_inverse_mpv}).
\end{proof}

\begin{definition}[Exchange of product bond factors]
    \label{def:mpug_product_bond_swap}
    \lean{MPOTensor.productBondSwapEquiv}
    \leanok
    For bond dimensions $D_M$ and $D_N$, identify a product bond index with
    its pair of factor indices. The exchange of product bond factors is the
    equivalence
    \begin{align}
        \{0,\ldots,D_M-1\}\times\{0,\ldots,D_N-1\}
                                                    & \longrightarrow
        \{0,\ldots,D_N-1\}\times\{0,\ldots,D_M-1\}, &
        (\alpha,\gamma)                             & \longmapsto (\gamma,\alpha).
        \label{eq:mpug_product_bond_swap}
    \end{align}
\end{definition}

\begin{theorem}[Physical adjunction reverses the operator product]
    \label{thm:mpug_physical_adjoint_product_reversal}
    \lean{MPOTensor.physicalAdjointTensor_mulTensor}
    \leanok
    \uses{def:mpug_product_bond_swap, def:mpdo_operator_product,
        def:mpdo_physical_adjoint}
    Let $M$ and $N$ be MPO tensors with bond dimensions $D_M$ and $D_N$.
    Write a left bond index of the product tensor as a pair
    $(\alpha,\gamma)$, with $\alpha$ in the bond space of $M$ and $\gamma$ in
    the bond space of $N$, and a right bond index as $(\beta,\delta)$
    likewise. Then, for every pair of physical indices $i,k$,
    \begin{align}
        \bigl((M\mathbin{\cdot}N)^\sharp\bigr)^{ik}
            _{(\alpha,\gamma),(\beta,\delta)}
         & =\bigl(N^\sharp\mathbin{\cdot}M^\sharp\bigr)^{ik}
        _{(\gamma,\alpha),(\delta,\beta)}.
        \label{eq:mpug_physical_adjoint_product_reversal}
    \end{align}
    The two bond spaces need not have the same dimension, so the exchange of
    the two factors of a product bond space is recorded here for unequal
    dimensions.
    % Source anchor: arXiv:2502.20257, line 1662 and prop:zetas,
    % lines 1662--1677.
    This is the inverse-product step used in the proof of the inversion
    identities for $\zeta$ and $\omega$
    in~\cite[lines~1662--1677]{FrancoRubio2025MPUGauging}, for the dagger
    operation described there, which conjugates every entry and exchanges the
    two legs lying on the same side of a tensor.
\end{theorem}

\begin{proof}\leanok
    \uses{def:mpug_product_bond_swap, def:mpdo_operator_product,
        def:mpdo_physical_adjoint}
    Expand both sides. The left side is the complex conjugate of the
    $\bigl((\alpha,\gamma),(\beta,\delta)\bigr)$ entry of
    $\sum_j M^{kj}\otimes N^{ji}$, namely
    $\sum_j\overline{M^{kj}_{\alpha\beta}}\,
        \overline{N^{ji}_{\gamma\delta}}$.
    The right side is the
    $\bigl((\gamma,\alpha),(\delta,\beta)\bigr)$ entry of
    $\sum_j (N^\sharp)^{ij}\otimes(M^\sharp)^{jk}$, namely
    $\sum_j\overline{N^{ji}_{\gamma\delta}}\,
        \overline{M^{kj}_{\alpha\beta}}$.
    The two sums agree term by term.
\end{proof}

Suppose now that the normalized flattening of every $\mathcal U_g$ is
left-canonical.  This is the canonical-form hypothesis used in
\cite[lines~1552--1557]{FrancoRubio2025MPUGauging}.

\begin{theorem}[Unitary gauges for the dagger and the group inverse]
    \label{thm:mpug_dagger_inverse_gauge}
    \lean{MPOTensor.GroupFamily.IsRepresentation.bondDim_inv,
        MPOTensor.GroupFamily.inverseFlattening,
        MPOTensor.GroupFamily.IsRepresentation.exists_daggerInverseGauge,
        MPOTensor.GroupFamily.IsRepresentation.daggerInverseGauge,
        MPOTensor.GroupFamily.IsRepresentation.physicalAdjointTensor_eq_daggerInverseGauge}
    \leanok
    \uses{def:mpug_group_representation, def:mpu_normalized_flattening,
        def:left_canonical_tensor, def:mpdo_physical_adjoint}
    Let $G$ be a group, and let
    $\{\mathcal U_g\}_{g\in G}$ be a representation by simple injective MPU
    tensors whose normalized flattenings are left-canonical.  Then
    \begin{align}
        D_g & =D_{g^{-1}}.
        \label{eq:mpug_dagger_inverse_bond_dimension}
    \end{align}
    After identifying these bond spaces, for every $g\in G$ there is a
    unitary matrix $T_g\in\operatorname{Mat}_{D_g}(\C)$ such that
    \begin{align}
        \mathcal U_g^\dagger
         & =T_g^\dagger\,\mathcal U_{g^{-1}}T_g.
        \label{eq:mpug_dagger_inverse_gauge}
    \end{align}
    Here $\mathcal U_g^\dagger$ is the physical-adjoint tensor.  This is
    the dagger--inverse gauge relation in
    \cite[lines~1552--1557]{FrancoRubio2025MPUGauging}.
\end{theorem}

\begin{proof}\leanok
    \uses{thm:mpug_physical_adjoint_inverse_mpv,
        thm:mps_equality_reduction_exists,
        thm:mps_rectangular_reduction_identities, thm:ft_single,
        thm:injective_irreducible,
        lem:left_canonical_irreducible_same_phase_unitary_gauge,
        thm:mpu_physical_adjoint_normalized_flattening}
    Apply the equality-case reduction theorem first to
    $\mathcal U_g^\dagger$ and $\mathcal U_{g^{-1}}$, and then with $g$ and
    $g^{-1}$ exchanged.  The two reductions give
    $D_{g^{-1}}\leq D_g$ and $D_g\leq D_{g^{-1}}$, proving
    (\ref{eq:mpug_dagger_inverse_bond_dimension}).  Their positive-length
    matrix product vectors agree; at length zero they agree because the bond
    dimensions have just been identified.  The single-block Fundamental
    Theorem therefore gives an invertible gauge.  The normalized flattenings
    of both tensors are left-canonical, while injectivity implies
    irreducibility, so the canonical-form unitary-gauge theorem replaces this
    gauge by a unitary one.  Reversing its orientation gives
    (\ref{eq:mpug_dagger_inverse_gauge}).
\end{proof}

For a complex matrix $A$, write
$A^*=(\overline{A_{ij}})_{i,j}$ for its entrywise complex conjugate.  This is
not the Hermitian adjoint $A^\dagger=(A^*)^{\mathsf T}$.

\begin{theorem}[Scalar relation for the dagger--inverse gauges]
    \label{thm:mpug_dagger_inverse_scalar}
    \lean{MPOTensor.GroupFamily.IsRepresentation.inverseDaggerInverseGauge,
        MPOTensor.GroupFamily.IsRepresentation.exists_daggerInverseScalar,
        MPOTensor.GroupFamily.IsRepresentation.daggerInverseScalar,
        MPOTensor.GroupFamily.IsRepresentation.daggerInverseGauge_mul_mapStar_inverse_eq_smul_one,
        MPOTensor.GroupFamily.IsRepresentation.daggerInverseScalar_mul_star_eq_one,
        MPOTensor.GroupFamily.IsRepresentation.norm_daggerInverseScalar,
        MPOTensor.GroupFamily.IsRepresentation.daggerInverseScalar_mul_inverse_eq_one,
        MPOTensor.GroupFamily.IsRepresentation.daggerInverseScalar_eq_one_or_neg_one_of_inv_eq}
    \leanok
    \uses{thm:mpug_dagger_inverse_gauge}
    The gauges in Theorem~\ref{thm:mpug_dagger_inverse_gauge} may be chosen
    together with scalars $\sigma_g\in\C$ satisfying
    \begin{align}
        T_gT_{g^{-1}}^* & =\sigma_g\mathbf 1.
        \label{eq:mpug_dagger_inverse_scalar}
    \end{align}
    Every such scalar lies in $U(1)$, and
    \begin{align}
        |\sigma_g|              & =1, &
        \sigma_g\sigma_{g^{-1}} & =1.
    \end{align}
    If $g^{-1}=g$, then $\sigma_g\in\{1,-1\}$.  These are
    the dagger--inverse scalar relation and its immediate consequences
    in \cite[lines~1557--1567]{FrancoRubio2025MPUGauging}.
\end{theorem}

\begin{proof}\leanok
    \uses{thm:mpug_dagger_inverse_gauge,
        thm:mpu_physical_adjoint_identities,
        lem:gauge_unique_up_to_scalar,
        thm:mpug_unitary_entrywise_scalar_phase,
        thm:mpug_unitary_entrywise_paired_scalars,
        thm:mpug_unitary_entrywise_self_sign}
    Apply physical adjunction to
    (\ref{eq:mpug_dagger_inverse_gauge}) for $g^{-1}$.  Its involutivity and
    unitarity give a second gauge from $\mathcal U_{g^{-1}}$ to
    $\mathcal U_g^\dagger$, whose matrix is $T_{g^{-1}}^*$.  Injective gauge
    uniqueness says that $T_{g^{-1}}^*$ is a nonzero scalar multiple of
    $T_g^\dagger$.  Multiplication on the left by $T_g$ gives
    (\ref{eq:mpug_dagger_inverse_scalar}) with the displayed orientation.
    The unit-modulus, inverse-pair, and self-inverse conclusions now follow
    from the three unitary matrix identities below.
\end{proof}

\begin{definition}[Common positive physical blocking]
    \label{def:mpug_common_blocking}
    \lean{MPOTensor.GroupFamily.block}
    \leanok
    \uses{def:mpug_group_representation, def:mpdo_physical_blocking}
    For a positive integer $L$, the common length-$L$ physical blocking is the
    family
    \begin{align}
        \widetilde{\mathcal U}_g & =\mathcal U_g^{[L]}.
        \label{eq:mpug_common_blocked_family}
    \end{align}
    Each blocked tensor has physical dimension $d^L$ and the same virtual bond
    dimension $D_g$ as $\mathcal U_g$. The physical alphabet is replaced
    simultaneously for every group element by words of length $L$.
\end{definition}

\begin{theorem}[Common blocking preserves the representation laws]
    \label{thm:mpug_common_blocking_representation}
    \lean{MPOTensor.GroupFamily.IsRepresentation.block,
        MPOTensor.mpo_blockTensor_eq_reindex,
        MPOTensor.blockTensor_mulTensor}
    \leanok
    \uses{def:mpug_common_blocking}
    Let $L>0$. The commonly blocked family
    $\{\mathcal U_g^{[L]}\}_{g\in G}$ is again a group representation by simple
    injective MPUs. Under the canonical identification of blocked
    configurations with words of length $NL$,
    \begin{align}
        O_N(\mathcal U_g^{[L]})
         & \cong U_g^{(NL)},
        \label{eq:mpug_blocked_operator} \\
        O_N(\mathcal U_e^{[L]})
         & =\Id,
        \label{eq:mpug_blocked_identity} \\
        O_N(\mathcal U_g^{[L]})O_N(\mathcal U_h^{[L]})
         & =O_N(\mathcal U_{gh}^{[L]})
        \label{eq:mpug_blocked_multiplication}
    \end{align}
    for every $N\geq1$. Physical blocking also commutes with tensor products:
    \begin{align}
        (\mathcal U_g\mathbin{\cdot}\mathcal U_h)^{[L]}
         & =\mathcal U_g^{[L]}\mathbin{\cdot}\mathcal U_h^{[L]}.
        \label{eq:mpug_block_product_tensor}
    \end{align}
\end{theorem}

\begin{proof}\leanok
    \uses{thm:mpug_positive_mpu_blocking,
        thm:mpdo_blocking_doubled_index,
        thm:mpu_bounded_simple_blocking,
        thm:mpdo_closed_chain_physical_blocking,
        thm:mpdo_blocking_product}
    Physical blocking preserves nonempty-chain unitarity, injectivity, and
    simplicity. By (\ref{eq:mpug_blocked_operator}), the blocked periodic operator is
    the original operator at length $NL$ up to the canonical simultaneous reindexing.
    Since $N,L>0$,
    the length $NL$ is positive, so (\ref{eq:mpug_rep_identity}) and
    (\ref{eq:mpug_rep_multiplication}) apply. Simultaneous reindexing preserves
    identities and products, which gives (\ref{eq:mpug_blocked_identity}) and
    (\ref{eq:mpug_blocked_multiplication}). Expanding one blocked letter gives
    (\ref{eq:mpug_block_product_tensor}).
\end{proof}

Assume from now until the end of this section that $G$ is finite and $d>0$.

\begin{definition}[Raw MPU representation]
    \label{def:mpug_raw_group_representation}
    \lean{MPOTensor.GroupFamily.IsRawRepresentation}
    \leanok
    \uses{def:injective, def:mpug_positive_mpu,
        def:mpug_group_representation}
    A raw MPU representation is a dependent-bond family
    $\{\mathcal U_g\}_{g\in G}$ of injective matrix product unitaries on
    nonempty chains satisfying (\ref{eq:mpug_rep_identity}) and
    (\ref{eq:mpug_rep_multiplication}) for all $g,h\in G$ and all $N\geq1$.
    No simplicity hypothesis is imposed.
\end{definition}

\begin{definition}[Common simplicity length]
    \label{def:mpug_common_simple_length}
    \lean{MPOTensor.GroupFamily.commonSimpleLength}
    \leanok
    \uses{def:mpug_raw_group_representation}
    For a finite dependent-bond family, define
    \begin{align}
        L_* & =\sum_{g\in G}D_g^4.
        \label{eq:mpug_common_simple_length}
    \end{align}
\end{definition}

\begin{theorem}[A finite MPU family has a common simple blocking]
    \label{thm:mpug_finite_common_simple_blocking}
    \lean{MPOTensor.GroupFamily.commonSimpleLength_pos,
        MPOTensor.GroupFamily.IsRawRepresentation.block_common,
        MPOTensor.GroupFamily.IsRawRepresentation.block_common_then_block}
    \leanok
    \uses{def:mpug_raw_group_representation,
        def:mpug_common_simple_length, def:mpug_group_representation,
        def:mpu_simple, def:mpdo_physical_blocking}
    Let $\{\mathcal U_g\}_{g\in G}$ be a raw MPU representation. Then $L_*>0$,
    and every $\mathcal U_g^{[L_*]}$ is simple. For every $K>0$, each further
    block $(\mathcal U_g^{[L_*]})^{[K]}$ is simple. Consequently the blocked
    family at length $L_*$ satisfies
    Definition~\ref{def:mpug_group_representation}, and every further common
    positive blocking does as well.
\end{theorem}

\begin{proof}\leanok
    \uses{thm:mpug_positive_mpu_blocking,
        thm:mpug_positive_mpu_to_mpu,
        thm:mpu_common_pow_four_simple_blocking,
        thm:mpu_all_later_simple_blockings,
        thm:mpdo_blocking_doubled_index,
        thm:mpdo_closed_chain_physical_blocking,
        thm:mpug_common_blocking_representation}
    Positivity of the bond dimensions gives $L_*>0$. For each $g$, the direct
    block of length $D_g^4$ is simple, and $D_g^4\leq L_*$ by
    (\ref{eq:mpug_common_simple_length}). Simplicity holds at every later direct
    blocking, hence at the common length $L_*$. Nonempty-chain unitarity and
    injectivity are preserved by physical blocking. Reindexing the closed chain
    at length $NL_*$ carries the identity and multiplication laws of the raw
    representation to the blocked family. The length-$L_*$ family is therefore
    a group representation by simple injective MPUs. Applying
    Theorem~\ref{thm:mpug_common_blocking_representation} to this established
    representation proves the assertion for every further positive blocking.
\end{proof}

\begin{definition}[Selected reductions for the pairs of a representation]
    \label{def:mpug_pair_reductions}
    \lean{MPOTensor.GroupFamily.productTensor,
        MPOTensor.GroupFamily.targetTensor,
        MPOTensor.GroupFamily.ReductionFamily,
        MPOTensor.GroupFamily.ReductionFamily.nilpotencyLength,
        MPOTensor.GroupFamily.ReductionFamily.HasExteriorBufferLength}
    \leanok
    \uses{def:mpug_group_representation, def:mpdo_operator_product,
        def:mpo_to_mps, def:mps_rectangular_reduction,
        def:mps_reduction_residual_nilpotency_length,
        def:mps_reduction_exterior_buffer_length}
    For $g,h\in G$, write
    \begin{align}
        B_{g,h} & =\mathcal U_g\mathbin{\cdot}\mathcal U_h,
                & A_{g,h}                                   & =\mathcal U_{gh}
        \label{eq:mpug_pair_tensors}
    \end{align}
    for the product tensor and the target tensor, both regarded as MPS
    tensors whose physical letters are the ket--bra pairs $(i,j)$, of bond
    dimensions $D_gD_h$ and $D_{gh}$.  A \emph{family of selected reductions}
    consists of matrices
    \begin{align}
        V_{g,h} & \in\C^{D_{gh}\times D_gD_h},
                & W_{g,h}                      & \in\C^{D_gD_h\times D_{gh}}
        \label{eq:mpug_pair_reduction_matrices}
    \end{align}
    such that $(V_{g,h},W_{g,h})$ is a reduction from $B_{g,h}$ to $A_{g,h}$
    for every pair: $V_{g,h}W_{g,h}=\Id$ and
    $V_{g,h}B_{g,h}^{\mathbf a}W_{g,h}=A_{g,h}^{\mathbf a}$ for every word
    $\mathbf a$.  The identity on words is the interior fusion equation in
    \cite[equation labelled \emph{eq:fusion\_2}, lines 1459--1497]
    {FrancoRubio2025MPUGauging}, with $F^<_{g,h}=V_{g,h}$ and
    $F^>_{g,h}=W_{g,h}$.  Write
    $N_0(g,h)=\ell(B_{g,h},A_{g,h};V_{g,h},W_{g,h})$ for the residual
    nilpotency length of the selected reduction of $(g,h)$.  The family has
    \emph{exterior buffer length} $m$ if $m$ is an exterior buffer length for
    every selected reduction:
    \begin{align}
        B_{g,h}^{\mathbf p}W_{g,h}A_{g,h}^{\mathbf c}V_{g,h}B_{g,h}^{\mathbf q}
         & =B_{g,h}^{\mathbf p\mathbf c\mathbf q}
        \label{eq:mpug_pair_exterior_identity}
    \end{align}
    for all $g,h\in G$, every nonempty word $\mathbf c$, and all words with
    $|\mathbf p|,|\mathbf q|\geq m$.  This is the exterior fusion equation in
    \cite[equation labelled \emph{eq:fusion\_1}, lines 1406--1457, with the
        quantifiers on line 1498]{FrancoRubio2025MPUGauging}, with $m\geq\ell$.
\end{definition}

\begin{theorem}[Existence of selected reductions]
    \label{thm:mpug_pair_reductions_exist}
    \lean{MPOTensor.GroupFamily.IsRepresentation.nonempty_reductionFamily}
    \leanok
    \uses{def:mpug_group_representation, def:mpug_pair_reductions}
    Every MPU representation by simple injective tensors admits a family of
    selected reductions.
\end{theorem}

\begin{proof}
    \leanok
    \uses{thm:mpug_product_tensor_operator,
        thm:mps_equality_reduction_exists}
    By (\ref{eq:mpug_product_tensor_positive_chain}), $B_{g,h}$ and $A_{g,h}$
    have the same positive-length matrix product vectors, and $A_{g,h}$ is
    injective with $D_{gh}>0$.  Theorem~\ref{thm:mps_equality_reduction_exists}
    gives a reduction for each pair; choose one for every pair.
\end{proof}

\begin{lemma}[Exterior buffers from the nilpotency lengths]
    \label{lem:mpug_pair_reductions_exterior}
    \lean{MPOTensor.GroupFamily.ReductionFamily.isReductionExteriorBufferLength_nilpotencyLength,
        MPOTensor.GroupFamily.ReductionFamily.hasExteriorBufferLength_of_le}
    \leanok
    \uses{def:mpug_pair_reductions}
    For a family of selected reductions of an MPU representation by simple
    injective tensors, $N_0(g,h)-1$ is an exterior buffer length for the
    selected reduction of $(g,h)$.  Consequently, if $N_0(g,h)\leq m+1$ for
    every pair, the family has exterior buffer length $m$.
\end{lemma}

\begin{proof}
    \leanok
    \uses{thm:mps_reduction_blocking, thm:mpug_product_tensor_operator}
    Equality of the positive-length matrix product vectors of $B_{g,h}$ and
    $A_{g,h}$ makes $N_0(g,h)$ a residual-word nilpotency bound, so
    Theorem~\ref{thm:mps_reduction_blocking} applies; enlarging the buffer
    preserves (\ref{eq:mpug_pair_exterior_identity}).
\end{proof}

\begin{definition}[Blocked selected reductions and the common buffer length]
    \label{def:mpug_pair_reductions_block}
    \lean{MPOTensor.GroupFamily.productTensor_block,
        MPOTensor.GroupFamily.targetTensor_block,
        MPOTensor.GroupFamily.ReductionFamily.block,
        MPOTensor.GroupFamily.ReductionFamily.commonBufferLength,
        MPOTensor.GroupFamily.ReductionFamily.commonBufferLength_pos,
        MPOTensor.GroupFamily.ReductionFamily.nilpotencyLength_le_commonBufferLength_add_one}
    \leanok
    \uses{def:mpug_pair_reductions, def:mpug_common_blocking,
        def:common_buffer_length}
    Let $L>0$.  Under the canonical pairing of blocked ket and bra words,
    \begin{align}
        \mathcal U_g^{[L]}\mathbin{\cdot}\mathcal U_h^{[L]}
         & \cong B_{g,h}^{[L]},
         & \mathcal U_{gh}^{[L]}
         & \cong A_{g,h}^{[L]},
        \label{eq:mpug_pair_tensors_block}
    \end{align}
    so a family of selected reductions $(V_{g,h},W_{g,h})$ of
    $\{\mathcal U_g\}_{g\in G}$ is, with the same matrices, a family of
    selected reductions of the blocked representation
    $\{\mathcal U_g^{[L]}\}_{g\in G}$.  The \emph{common buffer length} of
    the family is the common buffer length (\ref{eq:common_buffer_length})
    of the finite family $(N_0(g,h))_{(g,h)\in G\times G}$; it is positive
    and satisfies $N_0(g,h)\leq L+1$ for every pair.
\end{definition}

\begin{theorem}[Common blocking for the one-site fusion identities]
    \label{thm:mpug_common_blocking_one_site}
    \lean{MPOTensor.GroupFamily.ReductionFamily.block_hasExteriorBufferLength_one,
        MPOTensor.GroupFamily.ReductionFamily.block_commonBufferLength_hasExteriorBufferLength_one,
        MPOTensor.GroupFamily.IsRepresentation.exists_block_reductionFamily_one}
    \leanok
    \discussion{7396}
    \uses{def:mpug_pair_reductions_block, def:mpug_group_representation}
    Let $\{\mathcal U_g\}_{g\in G}$ be an MPU representation by simple
    injective tensors with a family of selected reductions, and let $L>0$
    satisfy $N_0(g,h)\leq L+1$ for every pair.  Then the blocked family of
    selected reductions of $\{\mathcal U_g^{[L]}\}_{g\in G}$ has exterior
    buffer length $1$: for all $g,h\in G$ and all nonempty blocked words
    $\mathbf p,\mathbf c,\mathbf q$,
    \begin{align}
        (B_{g,h}^{[L]})^{\mathbf p}W_{g,h}(A_{g,h}^{[L]})^{\mathbf c}
        V_{g,h}(B_{g,h}^{[L]})^{\mathbf q}
         & =(B_{g,h}^{[L]})^{\mathbf p\mathbf c\mathbf q}.
        \label{eq:mpug_common_blocking_one_site}
    \end{align}
    The common buffer length is such an $L$.  Consequently, every MPU
    representation of a finite group by simple injective tensors admits a
    positive $L$ such that $\{\mathcal U_g^{[L]}\}_{g\in G}$ is again an MPU
    representation by simple injective tensors and has a family of selected
    reductions with exterior buffer length $1$.

    This is the assertion of \cite{FrancoRubio2025MPUGauging},
    \texttt{main.tex} line 1498, and of the footnote to Theorem~1 of
    \cite{GarreRubioSchuch2024DomainWall}, that the nilpotency length can be
    reduced to $1$ by blocking, read as the one-site exterior identity
    (\ref{eq:mpug_common_blocking_one_site}).  Blocking need not make a
    blocked residual letter vanish, and the blocked nilpotency length
    $\ell(B_{g,h}^{[L]},A_{g,h}^{[L]};V_{g,h},W_{g,h})$ of Definition~8 in
    \cite{Molnar2018SemiInjective} need not be $1$
    \cite{gap:mgsc18_nilpotency_length_one_terminology}.
\end{theorem}

\begin{proof}
    \leanok
    \uses{lem:mpug_pair_reductions_exterior, thm:mps_reduction_blocking,
        thm:mpdo_blocking_doubled_index, thm:mpdo_blocking_product,
        thm:mpug_common_blocking_representation,
        thm:mpug_pair_reductions_exist}
    By Lemma~\ref{lem:mpug_pair_reductions_exterior}, $N_0(g,h)-1$ is an
    exterior buffer length for the pair $(g,h)$, and $N_0(g,h)-1\leq L$.
    Theorem~\ref{thm:mps_reduction_blocking} with $k=1$ makes $1$ an
    exterior buffer length for $(V_{g,h},W_{g,h})$ as a reduction from
    $B_{g,h}^{[L]}$ to $A_{g,h}^{[L]}$, and the pairing
    (\ref{eq:mpug_pair_tensors_block}), obtained from the commutation of
    blocking with the doubled-index view and with the product tensor,
    transports it to the blocked representation.  For the last assertion,
    choose a family of selected reductions by
    Theorem~\ref{thm:mpug_pair_reductions_exist}, take $L$ to be its common
    buffer length, and apply
    Theorem~\ref{thm:mpug_common_blocking_representation}.
\end{proof}

\begin{theorem}[Zero specified source-index value determines the source ranks]
    \label{thm:mpug_zero_specified_source_index_ranks}
    \lean{MPOTensor.sourceRanks_eq_physDim_of_sourceIndexValue_eq_zero}
    \leanok
    \uses{def:mpu_specified_source_index_value,
        def:mpu_right_rank, def:mpu_left_rank}
    Let $\mathcal U$ be a specified MPO tensor of positive physical dimension
    $d$, with positive right and left source-cut ranks $r$ and $\ell$. If
    \begin{align}
        r\ell & =d^2
    \end{align}
    and $\operatorname{ind}_{\mathrm{src}}(\mathcal U)=0$, then
    \begin{align}
        r    & =d, \\
        \ell & =d.
    \end{align}
    This is the arithmetic conclusion in
    \cite[line~1547]{FrancoRubio2025MPUGauging}, under the displayed
    hypotheses.
    % Scope: this entry does not assert finite-order vanishing, the rank-product
    % theorem, blocking independence, or the public MPU index.
    % Source anchor: Papers/2502.20257/main.tex, line 1547.
\end{theorem}

\begin{proof}\leanok
    \uses{thm:mpu_specified_source_index_identities}
    The source-index identity gives
    $0=\log_2(r/d)$. Since $r/d>0$, this implies $r/d=1$, hence $r=d$.
    Substitution into $r\ell=d^2$ and cancellation of the positive factor $d$
    give $\ell=d$.
\end{proof}

\begin{theorem}[Scalar phase from a unitary relation]
    \label{thm:mpug_unitary_entrywise_scalar_phase}
    \lean{Matrix.scalar_mul_star_eq_one_of_mul_map_star_eq_smul_one}
    \leanok
    Let $I$ be a nonempty finite set, let $T,S\in\operatorname{Mat}_I(\C)$ be
    unitary, and let $\sigma\in\C$. If
    \begin{align}
        TS^* & =\sigma\mathbf 1,
        \label{eq:mpug_unitary_entrywise_scalar_relation}
    \end{align}
    then
    \begin{align}
        \sigma\overline\sigma & =1.
        \label{eq:mpug_unitary_entrywise_scalar_phase}
    \end{align}
    This is the scalar phase conclusion following
    \cite[lines~1557--1562]{FrancoRubio2025MPUGauging} once the matrix relation
    has been supplied.
\end{theorem}

\begin{proof}\leanok
    The matrix $S^*$ is unitary, so $TS^*=\sigma\mathbf 1$ is unitary. Hence
    \begin{align}
        (\sigma\mathbf 1)(\sigma\mathbf 1)^\dagger
         & =\sigma\overline\sigma\,\mathbf 1=\mathbf 1.
    \end{align}
    Evaluating a diagonal entry gives
    (\ref{eq:mpug_unitary_entrywise_scalar_phase}).
\end{proof}

\begin{theorem}[Reciprocal scalars in the paired relations]
    \label{thm:mpug_unitary_entrywise_paired_scalars}
    \lean{Matrix.paired_scalars_mul_eq_one_of_mul_map_star_eq_smul_one}
    \leanok
    Let $I$ be a nonempty finite set, let $T,S\in\operatorname{Mat}_I(\C)$ be
    unitary, and let $\sigma,\tau\in\C$. If
    \begin{align}
        TS^* & =\sigma\mathbf 1, \\
        ST^* & =\tau\mathbf 1,
        \label{eq:mpug_unitary_entrywise_paired_relations}
    \end{align}
    then
    \begin{align}
        \sigma\tau & =1.
        \label{eq:mpug_unitary_entrywise_paired_scalars}
    \end{align}
    This is the conditional matrix-algebra content of
    $\sigma_g\sigma_{g^{-1}}=1$ in
    \cite[lines~1559--1563]{FrancoRubio2025MPUGauging}.
\end{theorem}

\begin{proof}\leanok
    \uses{thm:mpug_unitary_entrywise_scalar_phase}
    Multiplying the first relation by $T^{-1}$ on the left and by $T$ on the
    right gives $S^*T=\sigma\mathbf 1$. Entrywise conjugation of the second
    relation gives $S^*T=\overline\tau\,\mathbf 1$. A diagonal entry therefore
    gives $\sigma=\overline\tau$. Applying
    Theorem~\ref{thm:mpug_unitary_entrywise_scalar_phase} to the second relation
    yields $\tau\overline\tau=1$, and hence $\sigma\tau=1$.
\end{proof}

\begin{theorem}[Sign dichotomy in the self relation]
    \label{thm:mpug_unitary_entrywise_self_sign}
    \lean{Matrix.scalar_eq_one_or_neg_one_of_mul_map_star_self_eq_smul_one}
    \leanok
    Let $I$ be a nonempty finite set, let $T\in\operatorname{Mat}_I(\C)$ be
    unitary, and let $\sigma\in\C$. If
    \begin{align}
        TT^* & =\sigma\mathbf 1,
        \label{eq:mpug_unitary_entrywise_self_relation}
    \end{align}
    then
    \begin{align}
        \sigma & \in\{1,-1\}.
        \label{eq:mpug_unitary_entrywise_self_sign}
    \end{align}
    This is the sign step $\Id=xx^\dagger=e^{2i\phi}\Id$ in
    \cite[lines~1036--1038]{Cirac2017MPU}, and the involution conclusion in
    \cite[lines~1563--1567]{FrancoRubio2025MPUGauging}, conditional on the
    self relation.
\end{theorem}

\begin{proof}\leanok
    \uses{thm:mpug_unitary_entrywise_paired_scalars}
    Apply Theorem~\ref{thm:mpug_unitary_entrywise_paired_scalars} with $S=T$
    and $\tau=\sigma$. Then $\sigma^2=1$, so
    $\sigma=1$ or $\sigma=-1$.
\end{proof}

These three matrix identities supply the final algebraic steps in
Theorem~\ref{thm:mpug_dagger_inverse_scalar} after the representation-level
gauges and scalar relation have been constructed.

\section{A nondegenerate alternating-network comparison}
\label{sec:mpug_alternating_network}

\begin{definition}[Connected open alternating network]
    \label{def:mpug_open_alternating_network}
    \lean{MPSTensor.openAlternatingNetwork}
    \leanok
    Let $n\geq 0$, and let $D_0,\ldots,D_n$ be horizontal bond dimensions.  At each
    $0\leq i<n$, let $A_i$ have square bond dimension $D_i$, and let $B_i$ have
    left and right bond dimensions $D_i$ and $D_{i+1}$.  For fixed
    physical indices $\sigma=(\sigma_i)_i$ and $\tau=(\tau_i)_i$, define the
    connected network with its horizontal ends open by
    \begin{align}
        \mathcal N_{A,B}(\sigma,\tau)_{a,b}
         & :=\sum_{\substack{x_0,\ldots,x_n \\
                x_0=a,\,x_n=b}}
        \ \sum_{y_0,\ldots,y_{n-1}}
        \prod_{i=0}^{n-1}
        (A_i^{\sigma_i})_{x_i,y_i}
        (B_i^{\tau_i})_{y_i,x_{i+1}}.
        \notag
    \end{align}
    The indices $y_i$ join $A_i$ to $B_i$, while $x_{i+1}$ joins $B_i$ to
    $A_{i+1}$ when $i+1<n$; $x_0$ and $x_n$ are the two open end indices.
    This is the coordinate form of the connected open diagram in
    \cite[lines~2299--2361]{FrancoRubio2025MPUGauging}.
\end{definition}

\begin{lemma}[One-sided-inverse contraction of an alternating network]
    \label{lem:mpug_inverse_contraction_open_alternating_network}
    \lean{MPSTensor.inverseContraction_openAlternatingNetwork}
    \leanok
    \uses{def:injective, def:decomposition_map,
        def:mpug_open_alternating_network}
    Suppose that every $A_i$ is injective.  For the matrix unit $E_{z_i,q_i}$,
    let $c_i^{\sigma_i}(z_i,q_i)$ be its coefficients under a one-sided
    inverse of $A_i$, so that
    \begin{align}
        \sum_{\sigma_i}c_i^{\sigma_i}(z_i,q_i)A_i^{\sigma_i}
         & =E_{z_i,q_i}.
        \notag
    \end{align}
    Then, for every horizontal path $z=(z_i)_{i=0}^{n}$, every intermediate
    path $q=(q_i)_{i=0}^{n-1}$, and every $\tau$, one has
    \begin{align}
        \sum_{\sigma_0,\ldots,\sigma_{n-1}}
        \left(\prod_{i=0}^{n-1}c_i^{\sigma_i}(z_i,q_i)\right)
        \mathcal N_{A,B}(\sigma,\tau)_{z_0,z_n}
         & =\prod_{i=0}^{n-1}(B_i^{\tau_i})_{q_i,z_{i+1}}.
        \notag
    \end{align}
    This is the simultaneous inverse-tensor contraction in
    \cite[lines~2365--2392]{FrancoRubio2025MPUGauging}.
\end{lemma}

\begin{proof}\leanok
    Expand the connected network and interchange the finite sums.  At site
    $i$, the sum over $\sigma_i$ is the $(x_i,y_i)$ entry of
    $E_{z_i,q_i}$.  Its two Kronecker deltas force $x_i=z_i$ and $y_i=q_i$.
    Thus every internal sum collapses, leaving precisely the displayed product
    of entries of the $B_i$.
\end{proof}

\begin{theorem}[Nondegenerate corrected alternating-network proportionality]
    \label{thm:mpug_alternating_network_nondegenerate}
    \lean{MPSTensor.exists_scalars_of_openAlternatingNetwork_eq}
    \leanok
    \uses{def:injective, def:mpug_open_alternating_network}
    Let $n\geq 0$, and let $D_0,\ldots,D_n$ be horizontal bond dimensions.
    For each $0\leq i<n$, let $A_i$ be an injective MPS tensor with square bond
    dimension $D_i$, and let $B_i,B'_i$ be tensors whose left and right
    horizontal bond dimensions are $D_i$ and $D_{i+1}$.  Suppose that every
    $B_i$ is nonzero and that the connected open networks agree:
    \begin{align}
        A_0B_0A_1B_1\cdots A_{n-1}B_{n-1}
         & =A_0B'_0A_1B'_1\cdots A_{n-1}B'_{n-1}.
        \label{eq:mpug_alternating_network_equality}
    \end{align}
    Here the equality retains the two horizontal end legs and all physical
    legs.  Then there are scalars $\beta_i\in\C$, for $0\leq i<n$, such that
    \begin{align}
        B_i                      & =\beta_iB'_i,
        \label{eq:mpug_alternating_network_proportionality} \\
        \prod_{i=0}^{n-1}\beta_i & =1.
        \label{eq:mpug_alternating_network_scalar_product}
    \end{align}

    This is the nondegenerate corrected form of Lemma~\texttt{lemma:1} in
    \cite[lines~2296--2395]{FrancoRubio2025MPUGauging}.  The printed wording
    omits nonvanishing.  The paper's truncated symmetry $U^L_g$ is unitary at
    line~2099;
    lines~2444--2466 identify the side used at line~2567 with
    $(U^L_{g^{-1}})^\dagger$, while lines~3033--3053 identify the side used at
    line~3157 with $\sigma_g(U^L_{g^{-1}})^\dagger$ for
    $\sigma_g\in\{1,-1\}$.  Both displayed networks are therefore nonzero.  A
    vanishing intervening factor would make its entire network zero, so taking
    the factors on these sides as the $B_i$ supplies $B_i\ne 0$ at both
    applications.  This local nondegeneracy convention is recorded in
    \cite{gap:fbc25_alternating_network_nondegeneracy}.  The unrestricted
    printed statement is not asserted here.
\end{theorem}

\begin{proof}\leanok
    \uses{def:decomposition_map,
        lem:mpug_inverse_contraction_open_alternating_network,
        thm:periodic_product_tensor_phase_extraction}
    Apply Lemma~\ref{lem:mpug_inverse_contraction_open_alternating_network} to
    both sides of (\ref{eq:mpug_alternating_network_equality}), choosing the
    matrix units to prescribe arbitrary entries of all $B_i$.  The resulting
    connected horizontal contractions retain precisely those entries.  Since
    the entries were arbitrary, this gives the pure-tensor identity
    \begin{align}
        \bigotimes_{i=0}^{n-1}B_i
         & =\bigotimes_{i=0}^{n-1}B'_i.
        \label{eq:mpug_alternating_network_pure_tensor_equality}
    \end{align}
    Apply Theorem~\ref{thm:periodic_product_tensor_phase_extraction} to
    (\ref{eq:mpug_alternating_network_pure_tensor_equality}), with global
    scalar $1$ and the nonzero reference entries of the $B_i$.  It
    gives the factorwise proportionalities and
    (\ref{eq:mpug_alternating_network_scalar_product}); reversing the
    resulting nonzero scalars gives the orientation
    (\ref{eq:mpug_alternating_network_proportionality}).
\end{proof}

\section{Scalar three-cocycles and fusion-tensor gauge freedom}
\label{sec:mpug_scalar_gauge}

Let $G$ be a group and write $\C^\times$ for the multiplicative group of
nonzero complex numbers. The fusion tensors have the scalar gauge freedom
\begin{align}
    F^>_{g,h} & \longmapsto \beta_{g,h}F^>_{g,h},      \\
    F^<_{g,h} & \longmapsto \beta_{g,h}^{-1}F^<_{g,h}.
    \label{eq:mpug_fusion_gauge}
\end{align}
% Source anchor: arXiv:2502.20257, eq:scalar_fus_ten.
This is the fusion-tensor gauge freedom of~\cite{FrancoRubio2025MPUGauging}.
Apart from the matrix results below, the results in this section formalize
only the scalar functions $\beta$, $\omega$, $\gamma$, $L$, $\sigma$, and
$\Xi$; they construct neither fusion tensors nor action tensors. The
representation-level data were specified in
Definition~\ref{def:mpug_group_representation}.

\begin{theorem}[Scalar factors of an identity Kronecker product]
    \label{thm:mpug_kronecker_identity_factors}
    \lean{Matrix.exists_eq_smul_one_of_kronecker_eq_one}
    \leanok
    Let $M$ and $N$ be nonempty sets with decidable equality, and let
    $A\in\operatorname{Mat}_M(\C)$ and $B\in\operatorname{Mat}_N(\C)$. If
    $A\otimes B=\mathbf 1$, then there is a nonzero $c\in\C$ such that
    \begin{align}
        A & =c\mathbf 1, & B & =c^{-1}\mathbf 1.
        \label{eq:mpug_kronecker_identity_factors}
    \end{align}
\end{theorem}

\begin{proof}\leanok
    Choose $i_0\in M$ and $j_0\in N$. The
    $((i_0,j_0),(i_0,j_0))$ entry gives
    $A_{i_0i_0}B_{j_0j_0}=1$. Thus $c=A_{i_0i_0}$ is nonzero and
    $B_{j_0j_0}=c^{-1}$. Comparing every
    $((i,j_0),(i',j_0))$ entry gives $A_{ii'}=c\delta_{ii'}$; comparing every
    $((i_0,j),(i_0,j'))$ entry gives $B_{jj'}=c^{-1}\delta_{jj'}$.
\end{proof}

\begin{theorem}[Unitary factors of a unitary Kronecker product]
    \label{thm:mpug_unitary_kronecker_factors}
    \lean{Matrix.exists_pos_real_rescaling_mem_unitaryGroup_of_kronecker_mem_unitaryGroup,
        Matrix.exists_pos_real_smul_unitary_factors_of_kronecker_mem_unitaryGroup}
    \leanok
    Let $M$ and $N$ be nonempty finite sets. If
    $K_1\in\operatorname{Mat}_M(\C)$ and
    $K_2\in\operatorname{Mat}_N(\C)$ satisfy that $K_1\otimes K_2$ is
    unitary, then there are a positive real number $\delta$ and unitary
    matrices $W_1,W_2$ such that
    \begin{align}
        K_1 & =\delta W_1, & K_2 & =\delta^{-1}W_2.
        \label{eq:mpug_unitary_kronecker_factors}
    \end{align}
    Equivalently, $\delta^{-1}K_1$ and $\delta K_2$ are unitary.
    % Source anchor: arXiv:2502.20257, Lemma lem:deco, lines 1052--1066.
    This is only the tensor-factor step in the decomposition result
    of~\cite{FrancoRubio2025MPUGauging}; it does not construct or compare the
    three-legged decompositions occurring there.
\end{theorem}

\begin{proof}\leanok
    \uses{thm:mpug_kronecker_identity_factors}
    Unitarity gives
    \begin{align}
        (K_1^\dagger K_1)\otimes(K_2^\dagger K_2)=\mathbf 1.
    \end{align}
    Theorem~\ref{thm:mpug_kronecker_identity_factors} supplies a nonzero
    $c\in\C$ such that
    $K_1^\dagger K_1=c\mathbf 1$ and
    $K_2^\dagger K_2=c^{-1}\mathbf 1$. Positivity of the first Gram matrix
    makes $c$ a positive real number. Taking $\delta=\sqrt c$ shows directly
    that $\delta^{-1}K_1$ and $\delta K_2$ have identity Gram matrices; these
    are $W_1$ and $W_2$ in
    (\ref{eq:mpug_unitary_kronecker_factors}).
\end{proof}

\begin{theorem}[Cross-product rigidity]
    \label{thm:mpug_cross_product_rigidity}
    \lean{Kraus.IsNormal.eq_smul_of_cross_mul_eq,
        Kraus.IsInjective.eq_smul_of_cross_mul_eq}
    \leanok
    \uses{def:injective, def:normal}
    Let $A_i$ and $B_i$ be two families of square matrices of the same size,
    indexed by the same finite alphabet, and suppose that, for every $i,j$,
    \begin{align}
        B_iA_j & =A_iB_j.
        \label{eq:mpug_cross_product_rigidity}
    \end{align}
    If $A$ is injective, or more generally normal, then there is a scalar
    $c\in\C$ with $B_i=cA_i$ for every $i$.
    % Source anchor: arXiv:2502.20257, prop:MPUsplus, lines 1255--1325.
    This is the commuting-tensor step in the proof of the source-factor
    contraction identities
    in~\cite[lines~1255--1325]{FrancoRubio2025MPUGauging}, where the tensor
    commuting with a normal matrix product unitary tensor is shown to be a
    multiple of it.
\end{theorem}

\begin{proof}\leanok
    \uses{def:decomposition_map, thm:normal_tensor_gram_rigidity}
    Let $N\geq1$ be a blocking length at which the length-$N$ words of $A$
    span the full matrix algebra, and decompose the identity as
    $\mathbf 1=\sum_\sigma c_\sigma A^{\sigma_1}\cdots A^{\sigma_N}$. The
    cross relation moves the marked letter successively through the word,
    \begin{align}
        B_iA^{\sigma_1}\cdots A^{\sigma_N}
         & =A_iA^{\sigma_1}\cdots A^{\sigma_{N-1}}B^{\sigma_N}.
        \notag
    \end{align}
    When $N=1$, the middle word on the right is absent. Thus $B_i=A_iC$ with
    $C=\sum_\sigma c_\sigma
        A^{\sigma_1}\cdots A^{\sigma_{N-1}}B^{\sigma_N}$.
    Substituting this identity into
    (\ref{eq:mpug_cross_product_rigidity}) gives
    $A_i(CA_j-A_jC)=0$ for all $i$ and $j$. Multiplying on the left by a
    length-$N$ word, whose last letter absorbs one of the $A_i$, and
    recombining the words into the identity gives $CA_j=A_jC$. A matrix
    commuting with every letter of a normal family is a multiple of the
    identity, so $C=c\mathbf 1$ and $B_i=cA_i$. The injective case is the
    case of blocking length one.
\end{proof}

\begin{theorem}[Comparison of factorizations from one-sided inverses]
    \label{thm:mpug_factorization_comparison}
    \lean{Matrix.factorization_comparison_of_one_sided_inverses}
    \leanok
    Let $A,B,C$ be finite index sets and let $\mathbb S$ be a semiring. Suppose
    that
    \begin{align}
        X,\widetilde X & \in\operatorname{Mat}_{A\times B}(\mathbb S), &
        Y,\widetilde Y & \in\operatorname{Mat}_{B\times C}(\mathbb S),
    \end{align}
    and $XY=\widetilde X\widetilde Y$. If
    $L\in\operatorname{Mat}_{B\times A}(\mathbb S)$ and
    $R\in\operatorname{Mat}_{C\times B}(\mathbb S)$ satisfy
    $L\widetilde X=\mathbf 1$ and $YR=\mathbf 1$, then the comparison matrix
    $K=LX$ obeys
    \begin{align}
        K & =\widetilde YR, & X & =\widetilde XK, & \widetilde Y & =KY.
        \label{eq:mpug_factorization_comparison}
    \end{align}
    % Source anchor: arXiv:2502.20257, Lemma lem:deco, lines 1061--1065.
    These are the one-sided-inverse identities used in the proof of
    \cite[Lemma at lines~1052--1066]{FrancoRubio2025MPUGauging}.
\end{theorem}

\begin{proof}\leanok
    Associativity and the two inverse identities give
    \begin{align}
        LX
         & =(LX)(YR)
        =L(XY)R
        =L(\widetilde X\widetilde Y)R
        =\widetilde YR.
    \end{align}
    Multiplying $XY=\widetilde X\widetilde Y$ on the right by $R$ gives
    $X=\widetilde XK$. Multiplying it on the left by $L$ gives
    $\widetilde Y=KY$.
\end{proof}

\begin{theorem}[Reciprocal unitary gauges from a Kronecker comparison]
    \label{thm:mpug_reciprocal_unitary_gauges}
    \lean{Matrix.exists_reciprocal_unitary_gauges_of_comparison}
    \leanok
    For $i=1,2$, suppose that two compatible matrix factorizations are related
    by square comparison matrices $K_i$ through
    \begin{align}
        X_i & =\widetilde X_iK_i, & \widetilde Y_i & =K_iY_i.
        \label{eq:mpug_comparison_matrices}
    \end{align}
    If $K_1\otimes K_2$ is unitary, then there are a positive real number
    $\delta$ and unitary matrices $W_1,W_2$ such that
    \begin{align}
        \widetilde X_1 & =\delta^{-1}X_1W_1,          &
        \widetilde Y_1 & =\delta W_1^\dagger Y_1,       \\
        \widetilde X_2 & =\delta X_2W_2,              &
        \widetilde Y_2 & =\delta^{-1}W_2^\dagger Y_2.
        \label{eq:mpug_reciprocal_unitary_gauges}
    \end{align}
    Thus the positive scalar factors in the orientation of
    \cite[Lemma at lines~1052--1066]{FrancoRubio2025MPUGauging} are
    $\delta_1=\delta^{-1}$ and $\delta_2=\delta$, and hence
    $\delta_1\delta_2=1$.
    % Source anchor: arXiv:2502.20257, Lemma lem:deco, lines 1052--1066.
\end{theorem}

\begin{proof}\leanok
    \uses{thm:mpug_unitary_kronecker_factors}
    Apply Theorem~\ref{thm:mpug_unitary_kronecker_factors} to
    $K_1\otimes K_2$. It gives $\delta>0$ and unitary matrices $U_1,U_2$ with
    $K_1=\delta U_1$ and $K_2=\delta^{-1}U_2$. Set
    $W_i=U_i^\dagger$. Substitution into
    (\ref{eq:mpug_comparison_matrices}) gives
    (\ref{eq:mpug_reciprocal_unitary_gauges}); positivity of $\delta$ gives
    $\delta^{-1}\delta=1$.
\end{proof}

The full decomposition theorem also requires the comparison matrices obtained
from the MPU source-gate identities to have unitary Kronecker product. Once
this tensor-specific statement is established,
Theorems~\ref{thm:mpug_factorization_comparison} and
\ref{thm:mpug_reciprocal_unitary_gauges} yield the claimed gauges.
% Formalization boundary: the MPU-specific source-gate transport and the final
% lem:deco wrapper remain outside this slice and depend on issue #7313.

\begin{definition}[Unitary factorization comparison]
    \label{def:mpug_unitary_factorization_comparison}
    \lean{Matrix.unitaryFactorizationComparison,
        Matrix.unitaryFactorizationComparison_coe}
    \leanok
    Let $I$ be a set and let $R_{a,b}$ be a family of unitary matrices of a
    fixed size, indexed by $(a,b)\in I^2$. The unitary comparison from $(a,b)$
    to the reference pair $(c,d)$ is
    \begin{align}
        \lambda_{a,b}^{c,d}
         & =(R_{c,d})^{-1}R_{a,b}
        =R_{c,d}^\dagger R_{a,b}.
        \label{eq:mpug_unitary_factorization_comparison}
    \end{align}
    % Source anchor: arXiv:2502.20257, Corollary after Proposition 3,
    % lines 2782--2821, especially the defining formula at line 2820.
    This is the orientation in the corollary after Proposition~3
    of~\cite[lines~2782--2821]{FrancoRubio2025MPUGauging}. In that corollary,
    $I=G$ and comparisons are used when $ab=cd$; this equality selects the
    factorizations and is not needed in the unitary calculation.
\end{definition}

\begin{theorem}[Identity, inverse, composition, and right-reference laws]
    \label{thm:mpug_unitary_factorization_comparison_laws}
    \lean{Matrix.unitaryFactorizationComparison_self,
        Matrix.unitaryFactorizationComparison_inv,
        Matrix.unitaryFactorizationComparison_trans,
        Matrix.unitaryFactorizationComparison_of_reference_eq_one,
        Matrix.unitaryFactorizationComparison_right_reference}
    \leanok
    \uses{def:mpug_unitary_factorization_comparison}
    For all indices $a,b,c,d,f,g\in I$,
    \begin{align}
        \lambda_{a,b}^{a,b} & =\mathbf 1.
    \end{align}
    Reversing the comparison gives
    \begin{align}
        \left(\lambda_{a,b}^{c,d}\right)^{-1}
         & =\lambda_{c,d}^{a,b}.
    \end{align}
    Comparisons compose as
    \begin{align}
        \lambda_{a,b}^{f,g}
         & =\lambda_{c,d}^{f,g}\lambda_{a,b}^{c,d}.
        \label{eq:mpug_unitary_factorization_comparison_laws}
    \end{align}
    If $R_{c,d}=\mathbf 1$, then
    $\lambda_{a,b}^{c,d}=R_{a,b}$. In particular, if $I=G$ is a group, the
    normalization $R_{e,x}=\mathbf 1$ gives
    \begin{align}
        R_{g,h} & =\lambda_{g,h}^{e,gh}.
        \label{eq:mpug_unitary_factorization_right_reference}
    \end{align}
    % Source anchor: arXiv:2502.20257, Corollary after Proposition 3,
    % line 2817 (right reference and composition).
    The last identity is the right-reference specialization
    $\lambda^R_{g,h}=\lambda_{g,h}^{e,gh}$ in the cited corollary.
\end{theorem}

\begin{proof}\leanok
    \uses{def:mpug_unitary_factorization_comparison}
    The identity, inverse, and right-reference formulas follow from
    $R_{a,b}^{-1}R_{a,b}=\mathbf 1$, inversion of a product, and
    $\mathbf 1^{-1}R_{a,b}=R_{a,b}$. For composition, insert the identity at
    the intermediate pair:
    \begin{align}
        (R_{f,g})^{-1}R_{a,b}
         & =(R_{f,g})^{-1}R_{c,d}(R_{c,d})^{-1}R_{a,b}.
    \end{align}
\end{proof}

\begin{definition}[Gauss action on the gauge legs]
    \label{def:mpug_gauss_leg_action}
    \lean{TNLean.Algebra.gaussLegAction,
        TNLean.Algebra.gaussLegAction_apply}
    \leanok
    \discussion{7503}
    Let $G$ be a group. For $g\in G$, define the permutation
    $T_g\colon G^2\to G^2$ by
    \begin{align}
        T_g(a,b) & =\left(ag^{-1},gb\right).
        \label{eq:mpug_gauss_leg_action}
    \end{align}
    These are the two gauge-leg labels in the local Gauss operator of
    \cite[lines~3380--3387]{FrancoRubio2025MPUGauging}.
\end{definition}

\begin{theorem}[Composition of the Gauss leg action]
    \label{thm:mpug_gauss_leg_action_mul}
    \lean{TNLean.Algebra.gaussLegAction_one,
        TNLean.Algebra.gaussLegAction_mul}
    \leanok
    \discussion{7503}
    \uses{def:mpug_gauss_leg_action}
    For all $g,h\in G$,
    \begin{align}
        T_e & =\id.
    \end{align}
    Moreover,
    \begin{align}
        T_g\circ T_h & =T_{gh}.
        \label{eq:mpug_gauss_leg_action_mul}
    \end{align}
\end{theorem}

\begin{proof}\leanok
    \uses{def:mpug_gauss_leg_action}
    The two coordinates give
    \begin{align}
        T_g(T_h(a,b))
         & =\left(ah^{-1}g^{-1},ghb\right)
        =\left(a(gh)^{-1},(gh)b\right).
    \end{align}
\end{proof}

\begin{definition}[Local Gauss operator]
    \label{def:mpug_gauss_operator}
    \lean{TNLean.Algebra.gaussOperator,
        TNLean.Algebra.gaussOperator_apply_target,
        TNLean.Algebra.gaussOperator_apply_of_ne}
    \leanok
    \discussion{7503}
    \uses{def:mpug_unitary_factorization_comparison,
        def:mpug_gauss_leg_action}
    Let $G$ be finite, let $R_{a,b}\in\mathrm U(n)$ for $a,b\in G$, and use
    the basis $\ket{i,a,b}$ of $\C^n\otimes\C[G]^{\otimes2}$. Define
    \begin{align}
        \mathcal G_g
         & =\sum_{a,b\in G}
        \lambda_{a,b}^{ag^{-1},gb}
        \otimes\ket{ag^{-1},gb}\!\bra{a,b}.
        \label{eq:mpug_gauss_operator}
    \end{align}
    Equivalently, the block from $(a,b)$ to $T_g(a,b)$ is
    $(R_{T_g(a,b)})^{-1}R_{a,b}$ and every other block in that column
    vanishes. This is the local Gauss operator in
    \cite[lines~3380--3387]{FrancoRubio2025MPUGauging}.
\end{definition}

\begin{theorem}[Gauss operators form a unitary representation]
    \label{thm:mpug_gauss_representation}
    \lean{TNLean.Algebra.gaussOperator_one,
        TNLean.Algebra.gaussOperator_mul,
        TNLean.Algebra.star_gaussOperator,
        TNLean.Algebra.gaussOperator_mem_unitaryGroup,
        TNLean.Algebra.gaussRepresentation,
        TNLean.Algebra.gaussRepresentation_coe}
    \leanok
    \discussion{7503}
    \uses{def:mpug_gauss_operator}
    The local Gauss operators satisfy
    \begin{align}
        \mathcal G_e & =\id.
    \end{align}
    Their product law is
    \begin{align}
        \mathcal G_g\mathcal G_h & =\mathcal G_{gh}.
        \label{eq:mpug_gauss_representation}
    \end{align}
    Finally,
    \begin{align}
        \mathcal G_g^\dagger & =\mathcal G_{g^{-1}}.
    \end{align}
    Hence $g\mapsto\mathcal G_g$ is a unitary representation. This is the
    representation part of the local Gauss-law proposition in
    \cite[lines~3380--3387]{FrancoRubio2025MPUGauging}; invariance of the
    gauged tensor is a separate assertion and is not included here.
\end{theorem}

\begin{proof}\leanok
    \uses{thm:mpug_unitary_factorization_comparison_laws,
        thm:mpug_gauss_leg_action_mul,
        def:mpug_gauss_operator}
    In a column labelled by $(a,b)$, multiplication has only the intermediate
    gauge label $T_h(a,b)$. The surviving matter block is
    \begin{align}
        \lambda_{T_h(a,b)}^{T_g(T_h(a,b))}
        \lambda_{a,b}^{T_h(a,b)}
         & =\lambda_{a,b}^{T_{gh}(a,b)},
    \end{align}
    by the comparison composition law and
    $T_gT_h=T_{gh}$. At $e$, every surviving comparison is the identity.
    Reversing a comparison gives
    $(\lambda_{a,b}^{c,d})^\dagger=\lambda_{c,d}^{a,b}$, so taking the adjoint
    gives $\mathcal G_g^\dagger=\mathcal G_{g^{-1}}$. The product with this
    adjoint is $\mathcal G_e=\id$, which proves unitarity.
\end{proof}

\begin{definition}[Unplaced local Gauss projector kernel]
    \label{def:mpug_gauss_projector}
    \lean{TNLean.Algebra.gaussProjector,
        TNLean.Algebra.gaussProjector_eq_average}
    \leanok
    \discussion{7578}
    \uses{def:mpug_finite_group_unitary_average,
        thm:mpug_gauss_representation}
    Let $G$ be finite and let $\mathcal G_g$ act on
    $\mathcal H=\C^n\otimes\C[G]^{\otimes2}$. Define the unplaced local
    projector kernel
    \begin{align}
        P_G(R) & =\frac{1}{|G|}\sum_{g\in G}\mathcal G_g.
        \label{eq:mpug_gauss_projector}
    \end{align}
    This is the local average introduced in
    \cite[lines~3494--3499]{FrancoRubio2025MPUGauging}, before a position on
    the chain is chosen. This definition neither places the kernel on
    $(j,j+1)$ nor asserts invariance of the gauged matrix product state.
\end{definition}

\begin{theorem}[The local Gauss average is a projection]
    \label{thm:mpug_gauss_projector_projection}
    \lean{TNLean.Algebra.isStarProjection_gaussProjector}
    \leanok
    \discussion{7578}
    \uses{def:mpug_gauss_projector}
    The kernel in~(\ref{eq:mpug_gauss_projector}) satisfies
    \begin{align}
        P_G(R)^2       & =P_G(R), \\
        P_G(R)^\dagger & =P_G(R).
    \end{align}
\end{theorem}

\begin{proof}\leanok
    \uses{thm:mpug_finite_group_unitary_average_projection,
        thm:mpug_gauss_representation}
    Identify the matrix average with the averaging operator of the unitary
    representation $g\mapsto\mathcal G_g$. The invariant-space averaging
    theorem makes this operator idempotent. For self-adjointness, use
    $\mathcal G_g^\dagger=\mathcal G_{g^{-1}}$ and reindex the finite sum by
    the bijection $g\mapsto g^{-1}$.
\end{proof}

\begin{definition}[Gauss projector on a periodic chain]
    \label{def:mpug_placed_gauss_projector}
    \lean{MPOTensor.gaussLocalCoordinateEquiv,
        MPOTensor.gaussLocalCoordinateEquiv_apply_zero,
        MPOTensor.gaussLocalCoordinateEquiv_apply_one,
        MPOTensor.gaussWindowProjector,
        MPOTensor.placedGaussProjector}
    \leanok
    \discussion{7605}
    \uses{def:mpug_gauss_projector}
    Let the matter degree of freedom at one site have dimension $d$.  The
    local kernel has row and column coordinates
    \begin{align}
        \bigl((i_0,i_1),(a,b)\bigr)
         & \in (\{0,\ldots,d-1\}^{\{0,1\}})\times G^2.
    \end{align}
    Identify these with two site coordinates by the bijection
    \begin{align}
        \kappa\bigl((i_0,i_1),(a,b)\bigr)
         & =\bigl((i_0,a),(i_1,b)\bigr),
    \end{align}
    followed at each site by a fixed enumeration of
    $\{0,\ldots,d-1\}\times G$.  If $\widetilde P_G(R)$ denotes the matrix
    obtained from $P_G(R)$ by this simultaneous row and column relabelling,
    define its placement on a periodic chain of length $N\geq2$ by
    \begin{align}
        P_j(R) & =\operatorname{Embed}_{j,j+1}\bigl(\widetilde P_G(R)\bigr).
        \label{eq:mpug_placed_gauss_projector}
    \end{align}
    Here the second site is $j+1$ modulo $N$.  This is the position-dependent
    projector $P_j$ in~\cite[lines~3494--3499]{FrancoRubio2025MPUGauging}.
    % Source anchor: arXiv:2502.20257, eq:gauge_projs.
\end{definition}

\begin{theorem}[The placed Gauss projector is a projection]
    \label{thm:mpug_placed_gauss_projector_projection}
    \lean{MPOTensor.isStarProjection_gaussWindowProjector,
        MPOTensor.isStarProjection_placedGaussProjector}
    \leanok
    \discussion{7605}
    \uses{def:mpug_placed_gauss_projector}
    For every $N\geq2$ and every position $j$, the operator
    in~(\ref{eq:mpug_placed_gauss_projector}) satisfies
    \begin{align}
        P_j(R)^2       & =P_j(R), \\
        P_j(R)^\dagger & =P_j(R).
    \end{align}
\end{theorem}

\begin{proof}\leanok
    \uses{thm:mpug_gauss_projector_projection,
        thm:mpdo_cyclic_local_embedding_reindex}
    Let $U_\kappa$ be the permutation unitary determined by the coordinate
    bijection in Definition~\ref{def:mpug_placed_gauss_projector}, with the
    orientation
    \begin{align}
        U_\kappa\ket{((i_0,i_1),(a,b))}
         & =\ket{\kappa((i_0,i_1),(a,b))}.
        \notag
    \end{align}
    Simultaneous row and column reindexing is therefore the conjugation
    \begin{align}
        \widetilde P_G(R)
         & =U_\kappa P_G(R)U_\kappa^{-1}.
        \label{eq:mpug_gauss_reindex_conjugation}
    \end{align}
    Since $U_\kappa^{-1}=U_\kappa^\dagger$, Theorem~\ref{thm:mpug_gauss_projector_projection}
    and~(\ref{eq:mpug_gauss_reindex_conjugation}) give
    \begin{align}
        \widetilde P_G(R)^2
         & =U_\kappa P_G(R)^2U_\kappa^{-1}
        =\widetilde P_G(R),
        \label{eq:mpug_gauss_reindexed_idempotent} \\
        \widetilde P_G(R)^\dagger
         & =U_\kappa P_G(R)^\dagger U_\kappa^{-1}
        =\widetilde P_G(R).
        \label{eq:mpug_gauss_reindexed_adjoint}
    \end{align}
    Let $V_j$ be the permutation unitary that sends a chain configuration to
    its restriction to the cyclic window $(j,j+1)$ followed by its restriction
    to the complement. By Theorem~\ref{thm:mpdo_cyclic_local_embedding_reindex},
    \begin{align}
        V_jP_j(R)V_j^{-1}
         & =\widetilde P_G(R)\otimes\Id_{\mathrm{comp}}.
        \label{eq:mpug_gauss_window_tensor_identity}
    \end{align}
    Using $V_j^{-1}=V_j^\dagger$ and
    (\ref{eq:mpug_gauss_reindexed_idempotent})--(\ref{eq:mpug_gauss_reindexed_adjoint})
    in~(\ref{eq:mpug_gauss_window_tensor_identity}), we obtain
    \begin{align}
        V_jP_j(R)^2V_j^{-1}
         & =\bigl(\widetilde P_G(R)\otimes\Id_{\mathrm{comp}}\bigr)^2
        =\widetilde P_G(R)^2\otimes\Id_{\mathrm{comp}}
        =V_jP_j(R)V_j^{-1},                                                 \\
        V_jP_j(R)^\dagger V_j^{-1}
         & =\bigl(\widetilde P_G(R)\otimes\Id_{\mathrm{comp}}\bigr)^\dagger
        =\widetilde P_G(R)^\dagger\otimes\Id_{\mathrm{comp}}
        =V_jP_j(R)V_j^{-1}.
    \end{align}
    Conjugating by $V_j^{-1}$ gives $P_j(R)^2=P_j(R)$ and
    $P_j(R)^\dagger=P_j(R)$.
\end{proof}

\begin{definition}[Common fixed subspace]
    \label{def:mpug_common_fixed_subspace}
    \lean{LinearMap.commonFixedSubmodule}
    \leanok
    Let $\mathcal H$ be a complex vector space and let $(P_j)_{j\in J}$ be any
    family of endomorphisms of $\mathcal H$. Define
    their common fixed subspace by
    \begin{align}
        \mathcal V(P)
         & =\bigcap_{j\in J}\ker(P_j-\Id_{\mathcal H}).
    \end{align}
    No commutativity of the operators is required. For the averaged Gauss
    projectors, this is the common $+1$ eigenspace called the gauge-invariant
    subspace in
    \cite[lines~4327--4335 and line~5198]{FrancoRubio2025MPUGauging}.
    Here $P_j$ is an averaged projector, not an individual Gauss operator
    $\mathcal G_g$.
    % Source anchors: arXiv:2502.20257, lines 4327--4335 and line 5198.
\end{definition}

\begin{theorem}[Membership in the common fixed subspace]
    \label{thm:mpug_mem_common_fixed_subspace}
    \lean{LinearMap.mem_commonFixedSubmodule_iff}
    \leanok
    \uses{def:mpug_common_fixed_subspace}
    A vector $\psi\in\mathcal H$ belongs to $\mathcal V(P)$ if and only if
    $P_j\psi=\psi$ for every $j\in J$.
\end{theorem}

\begin{proof}\leanok
    \uses{def:mpug_common_fixed_subspace}
    Membership in an indexed intersection is equivalent to membership in
    every factor, and membership in $\ker(P_j-\Id_{\mathcal H})$ is equivalent
    to $P_j\psi=\psi$.
\end{proof}

\begin{theorem}[Nonzero-vector lower bound]
    \label{thm:mpug_common_fixed_subspace_finrank}
    \lean{LinearMap.one_le_finrank_commonFixedSubmodule}
    \leanok
    \uses{def:mpug_common_fixed_subspace}
    If $\mathcal H$ is finite-dimensional, $\psi\in\mathcal V(P)$, and
    $\psi\neq0$, then $1\leq\dim_{\C}\mathcal V(P)$.
\end{theorem}

\begin{proof}\leanok
    The common fixed subspace is nonzero because it contains $\psi$, so its
    dimension is at least one.
\end{proof}

\begin{definition}[Prescribed defect maps]
    \label{def:mpug_defect_maps}
    \lean{TNLean.Algebra.DefectMaps}
    \leanok
    \discussion{7569}
    Let $\mathcal H_{\mathrm m}$ be the matter space of two sites. Prescribed
    defect maps assign to every ordered pair $(a,b)\in G^2$ a defect subspace
    \begin{align}
        \mathcal K_{a,b} & \subseteq\mathcal H_{\mathrm m}
    \end{align}
    together with a map $D_{a,b}$ of $\mathcal H_{\mathrm m}$ whose
    restriction to that subspace is an isometry: for every
    $\ket{\xi}\in\mathcal K_{a,b}$,
    \begin{align}
        D_{a,b}^\dagger D_{a,b}\ket{\xi}
         & =\ket{\xi}.
        \label{eq:mpug_defect_isometry}
    \end{align}
    When the fusion scalars are not block independent, the state-level
    construction of~\cite[lines~4215--4326]{FrancoRubio2025MPUGauging}
    prescribes the modified fusion operators only on such subspaces.
    % Source anchor: arXiv:2502.20257, lemma:modif and the sentence closing
    % its proof.
\end{definition}

\begin{definition}[Full unitary completion class]
    \label{def:mpug_completion_class}
    \lean{TNLean.Algebra.DefectMaps.IsCompletion,
        TNLean.Algebra.DefectMaps.completionClass,
        TNLean.Algebra.DefectMaps.mem_completionClass_iff}
    \leanok
    \discussion{7569}
    \uses{def:mpug_defect_maps}
    The full completion class of prescribed defect maps
    $D=(D_{a,b})_{a,b\in G}$ is
    \begin{align}
        \mathfrak C(D)
         & =\left\{R=(U_{a,b})_{a,b\in G}:
        U_{a,b}\in\mathrm U(\mathcal H_{\mathrm m}),\
        U_{a,b}|_{\mathcal K_{a,b}}=D_{a,b}|_{\mathcal K_{a,b}}\right\}.
        \label{eq:mpug_completion_class}
    \end{align}
    A completion is the entire family $(U_{a,b})_{a,b\in G}$, not one
    selected operator of that family. If $\mathcal K_{a,b}$ does not span the
    whole matter space, the extension of $D_{a,b}$ need not be unique, as
    observed at the end of the proof of the modified-fusion lemma
    of~\cite[lines~4215--4326]{FrancoRubio2025MPUGauging}.
\end{definition}

\begin{theorem}[Existence of full unitary completions]
    \label{thm:mpug_completion_class_nonempty}
    \lean{TNLean.Algebra.DefectMaps.completionClass_nonempty,
        TNLean.Algebra.DefectMaps.exists_completion}
    \leanok
    \discussion{7702}
    \uses{def:mpug_defect_maps, def:mpug_completion_class}
    Let $\mathcal H_{\mathrm m}$ be finite-dimensional. Every family of
    prescribed defect maps satisfying (\ref{eq:mpug_defect_isometry}) admits a
    simultaneous unitary extension:
    \begin{align}
        \mathfrak C(D) & \neq\varnothing.
        \label{eq:mpug_completion_class_nonempty}
    \end{align}
    Equivalently, there is a family
    $R=(U_{a,b})_{(a,b)\in G^2}$ of unitary operators such that, for every
    $a,b\in G$ and every $\ket{\xi}\in\mathcal K_{a,b}$,
    \begin{align}
        U_{a,b}\ket{\xi}
         & =D_{a,b}\ket{\xi}.
        \label{eq:mpug_completion_exists}
    \end{align}
    This is the finite-dimensional extension step implicit in the freedom of
    the modified fusion operators at
    \cite[lines~4215--4326]{FrancoRubio2025MPUGauging}. The subspaces
    $\mathcal K_{a,b}$ and their prescribed isometries are hypotheses here.
    % Source anchor: arXiv:2502.20257, lemma:modif and its proof,
    % lines 4215--4326. The theorem does not assert that the defect data arise
    % from locally orthogonal MPS blocks or action tensors; their construction
    % remains in issues #7349, #7350, and #7569.
\end{theorem}

\begin{proof}\leanok
    \uses{def:mpug_defect_maps, def:mpug_completion_class}
    For each pair $(a,b)$, regard the restriction of $D_{a,b}$ to
    $\mathcal K_{a,b}$ as a linear isometry. Extend it to a unitary operator
    on the whole finite-dimensional matter space. Choosing one such extension
    for every pair gives a family $R$ satisfying
    (\ref{eq:mpug_completion_exists}), hence a member of $\mathfrak C(D)$.
\end{proof}

\begin{theorem}[Transport identity for a completion]
    \label{thm:mpug_completion_transport}
    \lean{TNLean.Algebra.DefectMaps.IsCompletion.mulVec_eq_prescribed,
        TNLean.Algebra.DefectMaps.IsCompletion.mulVec_eq_mulVec,
        TNLean.Algebra.DefectMaps.IsCompletion.unitaryFactorizationComparison_mulVec,
        TNLean.Algebra.DefectMaps.IsCompletion.unitaryFactorizationComparison_gaussLegAction_mulVec}
    \leanok
    \discussion{7569}
    \uses{def:mpug_completion_class,
        def:mpug_unitary_factorization_comparison,
        def:mpug_gauss_leg_action}
    Let $R=(U_{a,b})\in\mathfrak C(D)$, let
    $\ket{\xi}\in\mathcal K_{a,b}$ and $\ket{\eta}\in\mathcal K_{c,d}$, and
    assume the exact covariance
    \begin{align}
        D_{a,b}\ket{\xi} & =D_{c,d}\ket{\eta}.
        \label{eq:mpug_prescribed_covariance}
    \end{align}
    Then
    \begin{align}
        U_{a,b}\ket{\xi}             & =D_{a,b}\ket{\xi},
        \label{eq:mpug_completion_agrees}                  \\
        U_{a,b}\ket{\xi}             & =U_{c,d}\ket{\eta},
        \label{eq:mpug_completion_covariance}              \\
        \lambda_{a,b}^{c,d}\ket{\xi} & =\ket{\eta}.
        \label{eq:mpug_completion_transport}
    \end{align}
    Taking $(c,d)=T_g(a,b)$ gives the transport identity used for the
    modified Gauss laws of~\cite[lines~4325--4335]{FrancoRubio2025MPUGauging}.
\end{theorem}

\begin{proof}\leanok
    \uses{def:mpug_completion_class,
        thm:mpug_unitary_factorization_comparison_laws}
    Identity~(\ref{eq:mpug_completion_agrees}) is the defining property of a
    completion on the defect subspace. Applying it at $(a,b)$ and at $(c,d)$
    and inserting~(\ref{eq:mpug_prescribed_covariance})
    gives~(\ref{eq:mpug_completion_covariance}), an equality of vectors in the
    whole matter space. Multiplying that equality by $U_{c,d}^\dagger$ and
    using $U_{c,d}^\dagger U_{c,d}=\Id$ gives
    \begin{align}
        \lambda_{a,b}^{c,d}\ket{\xi}
         & =U_{c,d}^\dagger U_{a,b}\ket{\xi}
        =U_{c,d}^\dagger U_{c,d}\ket{\eta}
        =\ket{\eta}.
    \end{align}
    The adjoint is applied only after the equality in the whole matter space
    is available; no step asserts that $U_{c,d}^\dagger$ restricts to
    $D_{c,d}^\dagger$ on a defect subspace.
\end{proof}

\begin{definition}[Label action at one window]
    \label{def:mpug_chain_label_action}
    \lean{MPOTensor.chainLabelAction,
        MPOTensor.extractWindow_chainLabelAction_zero,
        MPOTensor.extractWindow_chainLabelAction_one}
    \leanok
    \discussion{7569}
    \uses{def:mpug_gauss_leg_action}
    Let $\alpha=(a_1,\ldots,a_N)$ be a configuration of gauge labels on a
    periodic chain of length $N\geq2$. For a window $(j,j+1)$ and $g\in G$,
    the configuration $T_{j,g}\alpha$ carries the labels
    \begin{align}
        \left(a_jg^{-1},ga_{j+1}\right)
    \end{align}
    on that window and agrees with $\alpha$ at every other site.
\end{definition}

\begin{theorem}[Composition of the label action]
    \label{thm:mpug_chain_label_action_mul}
    \lean{MPOTensor.chainLabelAction_one,
        MPOTensor.chainLabelAction_chainLabelAction,
        MPOTensor.chainLabelActionEquiv,
        MPOTensor.chainLabelActionEquiv_apply}
    \leanok
    \discussion{7569}
    \uses{def:mpug_chain_label_action}
    For all $g,h\in G$,
    \begin{align}
        T_{j,e}        & =\id,      \\
        T_{j,g}T_{j,h} & =T_{j,gh}.
        \label{eq:mpug_chain_label_action_mul}
    \end{align}
    In particular $T_{j,g}$ is a bijection of gauge-label configurations with
    inverse $T_{j,g^{-1}}$.
\end{theorem}

\begin{proof}\leanok
    \uses{def:mpug_chain_label_action}
    Outside the window both sides leave $\alpha$ unchanged. On the window,
    replacing the labels twice keeps only the second replacement, and the two
    labels are
    \begin{align}
        \left(a_jh^{-1}g^{-1},gha_{j+1}\right)
         & =\left(a_j(gh)^{-1},(gh)a_{j+1}\right).
    \end{align}
    At $g=e$ the labels are unchanged, so $T_{j,e}=\id$, and the composition
    law with $h=g^{-1}$ gives the inverse.
\end{proof}

\begin{definition}[Placed Gauss operator of a completion]
    \label{def:mpug_placed_gauss_operator}
    \lean{MPOTensor.gaussWindowOperator,
        MPOTensor.gaussWindowOperator_apply,
        MPOTensor.placedGaussOperator}
    \leanok
    \discussion{7569}
    \uses{def:mpug_gauss_operator,
        def:mpug_placed_gauss_projector,
        def:mpug_completion_class}
    Let $R\in\mathfrak C(D)$ and let a chain site carry one matter index and
    one gauge label. Relabelling the row and column coordinates of
    $\mathcal G_g$ by the bijection $\kappa$ of
    Definition~\ref{def:mpug_placed_gauss_projector} and embedding the result
    on the periodic window $(j,j+1)$ defines
    \begin{align}
        \widetilde{\mathcal G}^{j,j+1}_g(R)
         & =\operatorname{Embed}_{j,j+1}\bigl(\widetilde{\mathcal G}_g(R)\bigr),
        \label{eq:mpug_placed_gauss_operator}
    \end{align}
    the modified Gauss law
    of~\cite[lines~4325--4335]{FrancoRubio2025MPUGauging}, whose block from
    $(a,b)$ to $T_g(a,b)$ is $U_{T_g(a,b)}^\dagger U_{a,b}$.
\end{definition}

\begin{theorem}[The placed projector is the average of the placed Gauss operators]
    \label{thm:mpug_placed_gauss_projector_average}
    \lean{MPOTensor.placedGaussProjector_eq_average}
    \leanok
    \discussion{7569}
    \uses{def:mpug_placed_gauss_operator}
    For every window $(j,j+1)$,
    \begin{align}
        \widetilde P_j(R)
         & =\frac{1}{|G|}\sum_{g\in G}\widetilde{\mathcal G}^{j,j+1}_g(R).
        \label{eq:mpug_placed_gauss_projector_average}
    \end{align}
\end{theorem}

\begin{proof}\leanok
    \uses{def:mpug_gauss_projector, def:mpug_placed_gauss_operator}
    Comparing matrix entries, the relabelling is entrywise and the embedding
    multiplies each entry by the indicator that the two chain configurations
    agree outside the window. Both operations commute with the normalized sum
    over $G$ defining the local projector kernel
    in~(\ref{eq:mpug_gauss_projector}).
\end{proof}

\begin{definition}[Gauged vector of matter defect states]
    \label{def:mpug_gauged_vector}
    \lean{MPOTensor.gaugedMatter,
        MPOTensor.gaugedLabel,
        MPOTensor.gaugedVector,
        MPOTensor.gaugedVector_apply}
    \leanok
    \discussion{7569}
    Every chain configuration splits into its matter part and its gauge-label
    part. Given matter defect states $\ket{\Psi_\alpha}$ indexed by the
    gauge-label configurations $\alpha$ of a chain of length $N$, the gauged
    vector is
    \begin{align}
        \ket{\Psi_{\mathrm{gauged}}}
         & =\sum_\alpha\ket{\Psi_\alpha}\otimes\ket{\alpha},
        \label{eq:mpug_gauged_vector}
    \end{align}
    whose component at a chain configuration is the component of
    $\ket{\Psi_\alpha}$ at its matter part, $\alpha$ being its gauge-label
    part. This is the gauged state of
    \cite[lines~4325--4335]{FrancoRubio2025MPUGauging}.
\end{definition}

\begin{definition}[Local defect support and exact defect covariance]
    \label{def:mpug_defect_state_hypotheses}
    \lean{MPOTensor.HasDefectSupport,
        MPOTensor.HasPrescribedDefectCovariance}
    \leanok
    \discussion{7569}
    \uses{def:mpug_defect_maps, def:mpug_chain_label_action,
        def:mpug_gauged_vector}
    The matter defect states have local defect support when, for every
    $\alpha$ and every window $(j,j+1)$ whose adjacent labels are
    $(a_j,a_{j+1})$, the matter support of $\ket{\Psi_\alpha}$ lies in
    \begin{align}
        \mathcal K_{a_j,a_{j+1}}\otimes\mathcal H_{\mathrm{rest}}.
    \end{align}
    They satisfy exact defect covariance when
    \begin{align}
        \left(D_{a_j,a_{j+1}}\otimes\id_{\mathrm{rest}}\right)
        \ket{\Psi_\alpha}
         & =\left(D_{T_g(a_j,a_{j+1})}\otimes\id_{\mathrm{rest}}\right)
        \ket{\Psi_{T_{j,g}\alpha}}
        \label{eq:mpug_exact_defect_covariance}
    \end{align}
    for every $\alpha$, every window $(j,j+1)$, and every $g\in G$.
\end{definition}

\begin{theorem}[Universal invariant vector]
    \label{thm:mpug_universal_invariant_vector}
    \lean{MPOTensor.embedLocalOperator_mulVec_replaceWindow,
        MPOTensor.placedGaussOperator_mulVec_gaugedVector,
        MPOTensor.placedGaussProjector_mulVec_gaugedVector}
    \leanok
    \discussion{7569}
    \uses{def:mpug_placed_gauss_operator,
        def:mpug_placed_gauss_projector,
        def:mpug_completion_class,
        def:mpug_gauged_vector,
        def:mpug_defect_state_hypotheses}
    Assume local defect support and exact defect
    covariance~(\ref{eq:mpug_exact_defect_covariance}). Here these properties
    are hypotheses rather than consequences of locally orthogonal MPS blocks
    \cite{gap:fbc25_state_level_gauging_covariance}. For every
    $R\in\mathfrak C(D)$, every window $(j,j+1)$, and every $g\in G$,
    \begin{align}
        \widetilde{\mathcal G}^{j,j+1}_g(R)\ket{\Psi_{\mathrm{gauged}}}
         & =\ket{\Psi_{\mathrm{gauged}}},
        \label{eq:mpug_universal_invariant_vector} \\
        \widetilde P_j(R)\ket{\Psi_{\mathrm{gauged}}}
         & =\ket{\Psi_{\mathrm{gauged}}}.
    \end{align}
\end{theorem}

\begin{proof}\leanok
    \uses{thm:mpug_completion_transport,
        thm:mpug_chain_label_action_mul,
        thm:mpug_placed_gauss_projector_average,
        def:mpug_gauss_operator}
    Fix an output chain configuration and let $\alpha$ be its gauge-label part,
    with $(a,b)=(a_j,a_{j+1})$ the two labels adjacent to the window. Among the
    columns contributing to this fixed output row of
    $\widetilde{\mathcal G}^{j,j+1}_g(R)$, only the gauge labels
    $T_{j,g^{-1}}\alpha$ survive. Thus, with
    $\beta=T_{j,g^{-1}}\alpha$, the component of
    $\widetilde{\mathcal G}^{j,j+1}_g(R)\ket{\Psi_{\mathrm{gauged}}}$ equals
    the component of
    \begin{align}
        \left(U_{a,b}^\dagger U_{ag,g^{-1}b}
        \otimes\id_{\mathrm{rest}}\right)\ket{\Psi_\beta},
    \end{align}
    since the labels of $\beta$ adjacent to the window are $(ag,g^{-1}b)$ and
    $T_g(ag,g^{-1}b)=(a,b)$, whence $T_{j,g}\beta=\alpha$ by the composition
    law~(\ref{eq:mpug_chain_label_action_mul}). Local defect support puts the
    matter slices of $\ket{\Psi_\beta}$ and $\ket{\Psi_\alpha}$ over that
    window in $\mathcal K_{ag,g^{-1}b}$ and $\mathcal K_{a,b}$, and exact
    covariance~(\ref{eq:mpug_exact_defect_covariance}) relates them, so the
    transport identity~(\ref{eq:mpug_completion_transport}) turns that vector
    into $\ket{\Psi_\alpha}$. Identifying the label configuration $\beta$ with
    $\alpha$ through the bijection $T_{j,g}$ is the reindexing of the sum
    over gauge-label configurations in~(\ref{eq:mpug_gauged_vector}). Averaging
    over $g$ and using~(\ref{eq:mpug_placed_gauss_projector_average}) gives the
    second identity.
\end{proof}

\begin{definition}[Gauge-invariant subspace of a completion]
    \label{def:mpug_gauge_invariant_subspace}
    \lean{MPOTensor.gaugeInvariantSubspace}
    \leanok
    \discussion{7569}
    \uses{def:mpug_common_fixed_subspace,
        def:mpug_completion_class,
        def:mpug_placed_gauss_projector}
    For $R\in\mathfrak C(D)$ and a chain of length $N\geq2$, the
    gauge-invariant subspace is the common fixed subspace
    \begin{align}
        \mathcal V_N(R)
         & =\left\{\ket{\phi}:
        \widetilde P_j(R)\ket{\phi}=\ket{\phi}
        \text{ for every }j\right\}.
        \label{eq:mpug_gauge_invariant_subspace}
    \end{align}
    No commutativity of the projectors on neighbouring windows is assumed;
    \cite[lines~4325--4335]{FrancoRubio2025MPUGauging} observes that the
    modified projectors need not commute.
\end{definition}

\begin{theorem}[The gauged vector spans at least one dimension]
    \label{thm:mpug_gauge_invariant_subspace_finrank}
    \lean{MPOTensor.gaugedVector_mem_gaugeInvariantSubspace,
        MPOTensor.one_le_finrank_gaugeInvariantSubspace}
    \leanok
    \discussion{7569}
    \uses{def:mpug_gauge_invariant_subspace,
        def:mpug_completion_class,
        def:mpug_gauged_vector,
        def:mpug_defect_state_hypotheses}
    Under local defect support and exact defect
    covariance~(\ref{eq:mpug_exact_defect_covariance}),
    \begin{align}
        \ket{\Psi_{\mathrm{gauged}}} & \in\mathcal V_N(R)
    \end{align}
    for every $R\in\mathfrak C(D)$. If moreover
    $\ket{\Psi_{\mathrm{gauged}}}\neq0$, then
    \begin{align}
        \dim_{\C}\mathcal V_N(R) & \geq1.
    \end{align}
    Neither the minimum nor the maximum of $\dim_{\C}\mathcal V_N(R)$ over
    $\mathfrak C(D)$ is asserted; the growth of this dimension with $N$ is the
    question raised in
    \cite[line~5198]{FrancoRubio2025MPUGauging}.
\end{theorem}

\begin{proof}\leanok
    \uses{thm:mpug_universal_invariant_vector,
        thm:mpug_mem_common_fixed_subspace,
        thm:mpug_common_fixed_subspace_finrank}
    Membership is the projector identity in
    Theorem~\ref{thm:mpug_universal_invariant_vector} at every window. A
    nonzero member of a subspace forces its dimension to be at least one.
\end{proof}

\begin{definition}[Scalar cochains and normalization]
    \label{def:mpug_scalar_cochains}
    \lean{TNLean.Algebra.ScalarThreeCochain,
        TNLean.Algebra.ScalarCocycle.IsNormalized,
        TNLean.Algebra.ScalarThreeCochain.IsNormalized}
    \leanok
    \uses{def:scalar_cocycle}
    Recall from Definition~\ref{def:scalar_cocycle} that a scalar $2$-cochain
    is a function $\beta\colon G^2\to\C^\times$. A scalar $3$-cochain is a
    function $\omega\colon G^3\to\C^\times$. Following the paper's later
    standing convention, it is normalized when, for all $g,h,k\in G$,
    \begin{align}
        \omega(e,h,k) & = 1, \\
        \omega(g,e,k) & = 1, \\
        \omega(g,h,e) & = 1.
        \label{eq:mpug_three_normalized}
    \end{align}
    The standing convention of~\cite{FrancoRubio2025MPUGauging} makes every
    fusion tensor involving the identity trivial, which is what forces all
    three identity axes. Its printed display repeats the third axis in place
    of the first two; the three-axis form is the local correction recorded in
    \cite{gap:fbc25_three_cocycle_normalization_typo}.
    % Source anchor: arXiv:2502.20257, standing convention at lines 1936--1937
    %   and its display eq:triv_omegas, whose printed form repeats
    %   omega(g,h,e) three times.
    % Local fix: docs/paper-gaps/fbc25_three_cocycle_normalization_typo.tex.
    The scalar $2$-cochain $\beta$ is normalized when
    $\beta(e,g)=\beta(g,e)=1$ for every $g\in G$.
\end{definition}

\begin{definition}[Scalar three-cocycle]
    \label{def:mpug_three_cocycle}
    \lean{TNLean.Algebra.ScalarThreeCochain.IsCocycle}
    \leanok
    \uses{def:mpug_scalar_cochains}
    A scalar $3$-cochain is a $3$-cocycle when, for all $g,h,k,l\in G$,
    \begin{align}
        \omega(gh,k,l)\omega(g,h,kl)
         & =\omega(g,h,k)\omega(g,hk,l)\omega(h,k,l).
        \label{eq:mpug_three_cocycle}
    \end{align}
    % Source anchor: arXiv:2502.20257, eq:3-cocycle.
    This is the scalar $3$-cocycle relation
    of~\cite{FrancoRubio2025MPUGauging}.
\end{definition}

\begin{definition}[Three-coboundary and fusion-tensor gauge action]
    \label{def:mpug_fusion_gauge}
    \lean{TNLean.Algebra.ScalarThreeCochain.coboundary,
        TNLean.Algebra.ScalarThreeCochain.fusionGauge}
    \leanok
    \uses{def:mpug_scalar_cochains}
    For a scalar $2$-cochain $\beta$, define
    \begin{align}
        (d\beta)(g,h,k)
         & =\frac{\beta(g,hk)\beta(h,k)}
        {\beta(g,h)\beta(gh,k)},
        \label{eq:mpug_coboundary}        \\
        (\beta\mathbin{\triangleright}\omega)(g,h,k)
         & =(d\beta)(g,h,k)\omega(g,h,k).
        \label{eq:mpug_gauge_action}
    \end{align}
    % Source anchor: arXiv:2502.20257, eq:omegagauge.
    Formula~(\ref{eq:mpug_gauge_action}) is the change of $\omega$ induced by
    (\ref{eq:mpug_fusion_gauge}), as in~\cite{FrancoRubio2025MPUGauging}.
\end{definition}

\begin{theorem}[Identity and composition of fusion gauges]
    \label{thm:mpug_gauge_action}
    \lean{TNLean.Algebra.ScalarThreeCochain.fusionGauge_one,
        TNLean.Algebra.ScalarThreeCochain.fusionGauge_comp}
    \leanok
    \uses{def:mpug_fusion_gauge}
    Let $1$ denote the constant scalar $2$-cochain and let
    $(\beta\gamma)(g,h)=\beta(g,h)\gamma(g,h)$. Then
    \begin{align}
        1\mathbin{\triangleright}\omega & =\omega,
        \label{eq:mpug_gauge_identity}                                                  \\
        \gamma\mathbin{\triangleright}
        (\beta\mathbin{\triangleright}\omega)
                                        & =(\beta\gamma)\mathbin{\triangleright}\omega.
        \label{eq:mpug_gauge_composition}
    \end{align}
\end{theorem}

\begin{proof}\leanok
    \uses{def:mpug_fusion_gauge}
    The constant cochain has $d1=1$. Commutativity of $\C^\times$ gives
    $d(\beta\gamma)=(d\beta)(d\gamma)$, and substitution into
    (\ref{eq:mpug_gauge_action}) gives
    (\ref{eq:mpug_gauge_identity}) and
    (\ref{eq:mpug_gauge_composition}).
\end{proof}

\begin{theorem}[A three-coboundary is a three-cocycle]
    \label{thm:mpug_coboundary_cocycle}
    \lean{TNLean.Algebra.ScalarThreeCochain.isCocycle_coboundary}
    \leanok
    \uses{def:mpug_three_cocycle, def:mpug_fusion_gauge}
    For every scalar $2$-cochain $\beta$, the cochain $d\beta$ satisfies
    (\ref{eq:mpug_three_cocycle}).
\end{theorem}

\begin{proof}\leanok
    \uses{def:mpug_three_cocycle, def:mpug_fusion_gauge}
    Expanding the two sides gives the same product after cancellation:
    \begin{align}
         & (d\beta)(gh,k,l)(d\beta)(g,h,kl) \\
         & \quad =
        \frac{\beta(k,l)\beta(h,kl)\beta(g,hkl)}
        {\beta(ghk,l)\beta(gh,k)\beta(g,h)} \\
         & \quad =
        (d\beta)(g,h,k)(d\beta)(g,hk,l)(d\beta)(h,k,l).
    \end{align}
\end{proof}

\begin{theorem}[Products and fusion gauges preserve three-cocycles]
    \label{thm:mpug_gauge_preserves_cocycle}
    \lean{TNLean.Algebra.ScalarThreeCochain.IsCocycle.mul,
        TNLean.Algebra.ScalarThreeCochain.IsCocycle.fusionGauge}
    \leanok
    \uses{def:mpug_three_cocycle, def:mpug_fusion_gauge}
    The pointwise product of two scalar $3$-cocycles is a scalar
    $3$-cocycle. Consequently, if $\omega$ satisfies
    (\ref{eq:mpug_three_cocycle}), then so does
    $\beta\mathbin{\triangleright}\omega$.
\end{theorem}

\begin{proof}\leanok
    \uses{def:mpug_three_cocycle, thm:mpug_coboundary_cocycle}
    Multiply the two instances of (\ref{eq:mpug_three_cocycle}) and reorder
    their factors in $\C^\times$. Apply this product identity to
    $(d\beta)\omega$, using Theorem~\ref{thm:mpug_coboundary_cocycle} for
    $d\beta$.
\end{proof}

\begin{theorem}[Normalization under a fusion gauge]
    \label{thm:mpug_gauge_preserves_normalization}
    \lean{TNLean.Algebra.ScalarThreeCochain.isNormalized_coboundary_iff,
        TNLean.Algebra.ScalarThreeCochain.isNormalized_coboundary,
        TNLean.Algebra.ScalarThreeCochain.IsNormalized.mul,
        TNLean.Algebra.ScalarThreeCochain.IsNormalized.fusionGauge}
    \leanok
    \uses{def:mpug_scalar_cochains, def:mpug_fusion_gauge}
    The $3$-coboundary $d\beta$ is normalized if and only if there is a
    constant $c\in\C^\times$ such that, for every $g\in G$,
    \begin{align}
        \beta(e,g) & =c=\beta(g,e).
        \label{eq:mpug_coboundary_normalized_iff}
    \end{align}
    In particular, a normalized scalar $2$-cochain has normalized
    $3$-coboundary. Hence $\beta\mathbin{\triangleright}\omega$ is normalized
    whenever both $\beta$ and $\omega$ are normalized.
\end{theorem}

\begin{proof}\leanok
    \uses{def:mpug_scalar_cochains, def:mpug_fusion_gauge}
    Set one argument at a time equal to $e$ in
    (\ref{eq:mpug_coboundary}). If $d\beta$ is normalized, then
    (\ref{eq:mpug_three_normalized}) for $d\beta$ at $(e,e,g)$ and at
    $(g,e,e)$ gives $\beta(e,g)=\beta(e,e)=\beta(g,e)$. Conversely,
    (\ref{eq:mpug_coboundary_normalized_iff}) gives
    \begin{align}
        (d\beta)(e,h,k) & =1, \\
        (d\beta)(g,e,k) & =1, \\
        (d\beta)(g,h,e) & =1.
        \label{eq:mpug_coboundary_normalized_converse}
    \end{align}
    Taking the common constant to be $1$ proves the normalized-$\beta$
    corollary. Since pointwise products preserve
    (\ref{eq:mpug_coboundary_normalized_converse}), the gauged cochain
    $\beta\mathbin{\triangleright}\omega$ of (\ref{eq:mpug_gauge_action}) is
    normalized whenever both $\beta$ and $\omega$ are normalized.
\end{proof}

\begin{definition}[Cohomologous scalar three-cochains]
    \label{def:mpug_cohomologous_to}
    \lean{TNLean.Algebra.ScalarThreeCochain.CohomologousTo}
    \leanok
    \uses{def:mpug_fusion_gauge}
    Two scalar $3$-cochains $\omega$ and $\eta$ are cohomologous when
    $\omega=\beta\mathbin{\triangleright}\eta$ for some scalar
    $2$-cochain $\beta$.
\end{definition}

\begin{theorem}[Cohomologous-to is an equivalence relation]
    \label{thm:mpug_cohomologous_to_equivalence}
    \lean{TNLean.Algebra.ScalarThreeCochain.CohomologousTo.refl,
        TNLean.Algebra.ScalarThreeCochain.CohomologousTo.symm,
        TNLean.Algebra.ScalarThreeCochain.CohomologousTo.trans,
        TNLean.Algebra.ScalarThreeCochain.CohomologousTo.equivalence}
    \leanok
    \uses{def:mpug_cohomologous_to, thm:mpug_gauge_action}
    Cohomologous-to is reflexive, symmetric, and transitive.
\end{theorem}

\begin{proof}\leanok
    \uses{thm:mpug_gauge_action}
    Reflexivity uses the constant cochain $1$. If $\omega=
        \beta\mathbin{\triangleright}\eta$, then
    \begin{align}
        \beta^{-1}\mathbin{\triangleright}\omega
         & =(\beta\beta^{-1})\mathbin{\triangleright}\eta
        =\eta,
    \end{align}
    which proves symmetry. Composition formula
    (\ref{eq:mpug_gauge_composition}) proves transitivity.
\end{proof}

\begin{definition}[Trivial scalar gauge class]
    \label{def:mpug_trivial_gauge_class}
    \lean{TNLean.Algebra.ScalarThreeCochain.IsTrivialGaugeClass}
    \leanok
    \uses{def:mpug_cohomologous_to}
    A scalar $3$-cochain has trivial gauge class when it is cohomologous to
    the constant cochain $1$.
\end{definition}

\begin{definition}[L-symbols and compatibility]
    \label{def:mpug_l_symbols}
    \lean{TNLean.Algebra.LSymbol,
        TNLean.Algebra.LSymbol.IsCompatible}
    \leanok
    \uses{def:mpug_scalar_cochains}
    Let $G$ act on a set $\mathsf X$. An $L$-symbol is a function
    $L\colon \mathsf X\times G^2\to\C^\times$, written $L^x_{g,h}$.
    It is compatible with a scalar $3$-cochain $\omega$ when, for every
    $x\in\mathsf X$ and $g,h,k\in G$,
    \begin{align}
        L^x_{g,hk}L^x_{h,k}
         & =\omega(g,h,k)L^{k x}_{g,h}L^x_{gh,k}.
        \label{eq:mpug_l_compatibility}
    \end{align}
    The $L$-symbols were introduced by
    \cite{GarreRubio2023MPOSym}; this is the scalar compatibility relation of
    \cite{FrancoRubio2025MPUGauging}. No transitivity or finiteness assumption
    is imposed on the action.
    % Source anchor: arXiv:2502.20257, eq:omega_and_Ls.
\end{definition}

\begin{definition}[Joint scalar gauge of L-symbols]
    \label{def:mpug_l_gauge}
    \lean{TNLean.Algebra.ActionTensorGauge,
        TNLean.Algebra.LSymbol.gauge}
    \leanok
    \uses{def:mpug_l_symbols, def:mpug_fusion_gauge}
    An action-tensor scalar gauge is a function
    $\gamma\colon G\times\mathsf X\to\C^\times$, written $\gamma^x_g$.
    Together with a fusion-tensor scalar gauge $\beta\colon G^2\to\C^\times$,
    it changes an $L$-symbol by
    \begin{align}
        (\beta,\gamma)\mathbin{\triangleright}L^x_{g,h}
         & =\frac{\gamma^x_{gh}\beta_{g,h}}
        {\gamma^{h x}_g\gamma^x_h}L^x_{g,h}.
        \label{eq:mpug_l_gauge}
    \end{align}
    This is the joint fusion-tensor and action-tensor gauge freedom in
    \cite{FrancoRubio2025MPUGauging}, written with the block label as a
    superscript throughout, as that source does for its block-dependent
    cochains.
    % Source anchor: arXiv:2502.20257, eq:Lsymbgauge (which spells the block
    %   label as a second subscript); the superscript spelling gamma^x_g is
    %   the paper's own at lines 5908 and 5988--5990.
\end{definition}

\begin{theorem}[Identity and composition of joint scalar gauges]
    \label{thm:mpug_l_gauge_composition}
    \lean{TNLean.Algebra.LSymbol.gauge_one,
        TNLean.Algebra.LSymbol.gauge_comp}
    \leanok
    \uses{def:mpug_l_gauge}
    The constant fusion and action gauges leave every $L$-symbol unchanged.
    Successive joint gauges compose by pointwise multiplication:
    \begin{align}
        (\delta,\varepsilon)\mathbin{\triangleright}
        \bigl((\beta,\gamma)\mathbin{\triangleright}L\bigr)
         & = (\beta\delta,\gamma\varepsilon)
        \mathbin{\triangleright}L.
        \label{eq:mpug_l_gauge_composition}
    \end{align}
\end{theorem}

\begin{proof}\leanok
    \uses{def:mpug_l_gauge}
    Substitute (\ref{eq:mpug_l_gauge}) twice and cancel the intermediate
    factors. The identity case follows by taking all four cochains equal to
    $1$.
\end{proof}

\begin{definition}[Normalized L-symbols and action-tensor gauges]
    \label{def:mpug_l_normalized}
    \lean{TNLean.Algebra.LSymbol.IsNormalized,
        TNLean.Algebra.ActionTensorGauge.IsNormalized}
    \leanok
    \uses{def:mpug_l_symbols, def:mpug_l_gauge}
    The $L$-symbol is normalized when, for every $x\in\mathsf X$ and $g\in G$,
    \begin{align}
        L^x_{g,e} & =1.
        \label{eq:mpug_l_normalized_right}
    \end{align}
    \begin{align}
        L^x_{e,g} & =1.
        \label{eq:mpug_l_normalized_left}
    \end{align}
    Equations (\ref{eq:mpug_l_normalized_right}) and
    (\ref{eq:mpug_l_normalized_left}) are the normalization of the
    $L$-symbols in \cite{FrancoRubio2025MPUGauging}. The action-tensor gauge
    is normalized when
    \begin{align}
        \gamma^x_e & =1
        \label{eq:mpug_action_gauge_normalized}
    \end{align}
    for every $x\in\mathsf X$. This last condition is not a displayed
    equation of the source: it is the scalar shadow of the standing
    convention that every action tensor involving the identity is trivial.
    % Source anchor: arXiv:2502.20257, eq:triv_Ls for the L-symbol axes;
    %   gamma^x_e = 1 is not a labelled source equation but the standing
    %   convention at lines 1936--1937 that any fusion or action tensor
    %   involving e is trivial, applied to the action-tensor gauge scalar of
    %   eq:scalar_act_ten.
\end{definition}

\begin{definition}[Restricted scalar gauges]
    \label{def:mpug_restricted_scalar_gauge}
    \lean{TNLean.Algebra.ScalarCocycle.IsInverseNormalized,
        TNLean.Algebra.RestrictedFusionGauge,
        TNLean.Algebra.RestrictedScalarGauge}
    \leanok
    \discussion{7700}
    \uses{def:mpug_fusion_gauge, def:mpug_l_gauge, def:mpug_l_normalized}
    A fusion scalar $\beta\colon G^2\to\C^\times$ is
    \emph{inverse-normalized} when, for every $g\in G$,
    \begin{align}
        \beta_{g,e}=\beta_{e,g}=\beta_{g,g^{-1}}
         & =1.
        \label{eq:mpug_restricted_fusion_gauge}
    \end{align}
    A restricted scalar gauge is a pair $(\beta,\gamma)$ in which $\beta$ is
    inverse-normalized and, for every $x\in\mathsf X$,
    \begin{align}
        \gamma^x_e & =1.
        \label{eq:mpug_restricted_action_gauge}
    \end{align}
    The conditions are the standing restriction imposed at
    \cite[lines~2050--2054]{FrancoRubio2025MPUGauging} after the earlier
    fusion and action gauges have been fixed.
    % Source anchor: arXiv:2502.20257, unlabelled display at lines 2050--2054,
    % immediately after eq:zetas_in_gauge.
\end{definition}

\begin{theorem}[Joint scalar gauges preserve L-symbol compatibility]
    \label{thm:mpug_l_gauge_preservation}
    \lean{TNLean.Algebra.LSymbol.IsCompatible.gauge,
        TNLean.Algebra.LSymbol.gauge_apply_right_one,
        TNLean.Algebra.LSymbol.gauge_apply_left_one,
        TNLean.Algebra.LSymbol.isNormalized_gauge_iff,
        TNLean.Algebra.LSymbol.IsNormalized.gauge}
    \leanok
    \uses{def:mpug_fusion_gauge, def:mpug_l_symbols,
        def:mpug_l_gauge, def:mpug_l_normalized}
    If $L$ is compatible with $\omega$, then
    $(\beta,\gamma)\mathbin{\triangleright}L$ is compatible with
    $\beta\mathbin{\triangleright}\omega$. On the identity axes, the gauge
    formula is exactly
    \begin{align}
        ((\beta,\gamma)\mathbin{\triangleright}L)^x_{g,e}
         & =\frac{\beta_{g,e}}{\gamma^x_e}L^x_{g,e}.
        \label{eq:mpug_l_gauge_right}
    \end{align}
    \begin{align}
        ((\beta,\gamma)\mathbin{\triangleright}L)^x_{e,g}
         & =\frac{\beta_{e,g}}{\gamma^{g x}_e}L^x_{e,g}.
        \label{eq:mpug_l_gauge_left}
    \end{align}
    Thus the gauged $L$-symbol is normalized exactly when
    $\beta_{g,e}L^x_{g,e}=\gamma^x_e$ and
    $\beta_{e,g}L^x_{e,g}=\gamma^{g x}_e$ for all $x$ and $g$. In particular,
    normalized $\beta$, $\gamma$, and $L$ give a normalized gauged $L$-symbol.
\end{theorem}

\begin{proof}\leanok
    \uses{def:mpug_fusion_gauge, def:mpug_l_symbols, def:mpug_l_gauge,
        def:mpug_l_normalized}
    Write $L'=(\beta,\gamma)\mathbin{\triangleright}L$. Substituting
    (\ref{eq:mpug_l_gauge}) into the left side of
    (\ref{eq:mpug_l_compatibility}) for $L'$ cancels the factor
    $\gamma^x_{hk}$ and leaves
    \begin{align}
        {L'}^x_{g,hk}{L'}^x_{h,k}
         & =\frac{\gamma^x_{ghk}}
        {\gamma^{(hk)x}_g\gamma^{k x}_h\gamma^x_k}\,
        \beta_{g,hk}\beta_{h,k}\,L^x_{g,hk}L^x_{h,k}.
        \label{eq:mpug_l_gauge_left_side}
    \end{align}
    The same substitution on the right side cancels
    $\gamma^{k x}_{gh}$ and leaves
    \begin{align}
        {L'}^{k x}_{g,h}{L'}^x_{gh,k}
         & =\frac{\gamma^x_{ghk}}
        {\gamma^{h(k x)}_g\gamma^{k x}_h\gamma^x_k}\,
        \beta_{g,h}\beta_{gh,k}\,L^{k x}_{g,h}L^x_{gh,k}.
        \label{eq:mpug_l_gauge_right_side}
    \end{align}
    The action law $h(k x)=(hk)x$ makes the two action-tensor prefactors of
    (\ref{eq:mpug_l_gauge_left_side}) and
    (\ref{eq:mpug_l_gauge_right_side}) equal. Applying
    (\ref{eq:mpug_l_compatibility}) for $L$ and $\omega$ to the right of
    (\ref{eq:mpug_l_gauge_left_side}) therefore gives
    \begin{align}
        {L'}^x_{g,hk}{L'}^x_{h,k}
         & =\frac{\beta_{g,hk}\beta_{h,k}}
        {\beta_{g,h}\beta_{gh,k}}\,\omega(g,h,k)\,
        {L'}^{k x}_{g,h}{L'}^x_{gh,k},
        \label{eq:mpug_l_gauge_compatibility}
    \end{align}
    whose scalar factor is
    $(\beta\mathbin{\triangleright}\omega)(g,h,k)$ by
    (\ref{eq:mpug_coboundary}) and
    (\ref{eq:mpug_gauge_action}). Setting either group argument equal
    to $e$ in (\ref{eq:mpug_l_gauge}) gives
    (\ref{eq:mpug_l_gauge_right}) and
    (\ref{eq:mpug_l_gauge_left}), from which the normalization statements
    follow.
\end{proof}

\begin{theorem}[Action of the restricted scalar gauge group]
    \label{thm:mpug_restricted_scalar_gauge_action}
    \lean{TNLean.Algebra.ScalarCocycle.isInverseNormalized_one,
        TNLean.Algebra.ScalarCocycle.IsInverseNormalized.mul,
        TNLean.Algebra.ScalarCocycle.IsInverseNormalized.inv,
        TNLean.Algebra.ActionTensorGauge.isNormalized_one,
        TNLean.Algebra.ActionTensorGauge.IsNormalized.mul,
        TNLean.Algebra.ActionTensorGauge.IsNormalized.inv,
        TNLean.Algebra.RestrictedFusionGauge.one,
        TNLean.Algebra.RestrictedFusionGauge.mul,
        TNLean.Algebra.RestrictedFusionGauge.inv,
        TNLean.Algebra.RestrictedScalarGauge.one,
        TNLean.Algebra.RestrictedScalarGauge.mul,
        TNLean.Algebra.RestrictedScalarGauge.inv,
        TNLean.Algebra.RestrictedFusionGauge.fusionGauge_one,
        TNLean.Algebra.RestrictedFusionGauge.isNormalized_fusionGauge,
        TNLean.Algebra.RestrictedFusionGauge.isCocycle_fusionGauge,
        TNLean.Algebra.RestrictedFusionGauge.fusionGauge_mul,
        TNLean.Algebra.RestrictedFusionGauge.fusionGauge_inv,
        TNLean.Algebra.RestrictedScalarGauge.gauge_one,
        TNLean.Algebra.RestrictedScalarGauge.isNormalized_gauge,
        TNLean.Algebra.RestrictedScalarGauge.isCompatible_gauge,
        TNLean.Algebra.RestrictedScalarGauge.gauge_mul,
        TNLean.Algebra.RestrictedScalarGauge.gauge_inv}
    \leanok
    \discussion{7700}
    \uses{def:mpug_restricted_scalar_gauge}
    Let $\kappa=(\beta,\gamma)$ be a restricted scalar gauge. If $\omega$ is
    a normalized scalar $3$-cocycle and $L$ is a normalized $L$-symbol
    compatible with $\omega$, then
    \begin{align}
        \omega' & =\beta\mathbin{\triangleright}\omega,     \\
        L'      & =(\beta,\gamma)\mathbin{\triangleright}L.
        \label{eq:mpug_restricted_gauge_action}
    \end{align}
    are again normalized, $\omega'$ is a scalar $3$-cocycle, and $L'$ is
    compatible with $\omega'$. The identity gauge acts trivially, successive
    restricted gauges compose by pointwise multiplication, and applying
    $\kappa^{-1}$ after $\kappa$ recovers $(\omega,L)$.
    These are scalar statements only.
    % The theorem does not choose fusion tensors, action tensors, or
    % tensor-level gauges satisfying the same normalization. Tensor-level
    % restricted gauge choices remain issue #7325.
\end{theorem}

\begin{proof}\leanok
    \uses{def:mpug_restricted_scalar_gauge,
        thm:mpug_gauge_preserves_cocycle,
        thm:mpug_gauge_preserves_normalization,
        thm:mpug_l_gauge_composition,
        thm:mpug_l_gauge_preservation}
    The three specified values of $\beta$ and the identity value of $\gamma$
    remain equal to $1$ under pointwise products and inverses. Apply the
    general preservation laws to~(\ref{eq:mpug_restricted_gauge_action}).
    The group laws and the composition formulas for the two gauge actions give
    the identity, composition, and inverse assertions.
\end{proof}

\begin{theorem}[Cocycles and coboundaries for trivial coefficients]
    \label{thm:mpug_h2_trivial_coefficients}
    \lean{TNLean.Algebra.ScalarCocycle.isCocycle_iff_isMulCocycle₂,
        TNLean.Algebra.ScalarCocycle.isCoboundary_iff_isMulCoboundary₂,
        TNLean.Algebra.ScalarCocycle.cohomologousTo_iff_isCoboundary_div,
        TNLean.Algebra.ScalarCocycle.isCoboundary_iff_mem_mathlibCoboundaries,
        TNLean.Algebra.ScalarCocycle.cohomologousTo_iff_H2π_eq}
    \leanok
    \uses{def:is_coboundary, def:cohomologous_to, def:is_cocycle}
    Let $G$ act trivially on $\C^\times$, and let
    $f(g,h)=\omega(g,h)$. The multiplicative cocycle equation
    \begin{align}
        f(gh,k)f(g,h)
         & =f(h,k)f(g,hk)
        \label{eq:mpug_h2_cocycle_comparison}
    \end{align}
    is equivalent to Definition~\ref{def:is_cocycle}. Likewise, the
    multiplicative coboundary equation
    \begin{align}
        \frac{\varphi(h)}{\varphi(gh)}\varphi(g)
         & =f(g,h)
        \label{eq:mpug_h2_coboundary_comparison}
    \end{align}
    is equivalent to Definition~\ref{def:is_coboundary}.

    If $\pi$ denotes the canonical map from multiplicative $2$-cocycles to
    degree-two group cohomology, then for any two cocycles $\omega_1$ and
    $\omega_2$,
    \begin{align}
        \omega_1\sim\omega_2
         & \quad\Longleftrightarrow\quad
        \pi(\omega_1)=\pi(\omega_2).
        \label{eq:mpug_h2_class_comparison}
    \end{align}
    % Source anchor: arXiv:2502.20257, lines 1927--1931.
\end{theorem}

\begin{proof}\leanok
    \uses{def:is_coboundary, def:cohomologous_to, def:is_cocycle}
    Since $\C^\times$ is commutative,
    (\ref{eq:mpug_h2_cocycle_comparison}) is obtained from the cocycle
    equation by interchanging the two factors on each side, and
    (\ref{eq:mpug_h2_coboundary_comparison}) is obtained from
    $\varphi(g)\varphi(h)\varphi(gh)^{-1}$ by reordering its factors.
    Moreover, $\omega_1\sim\omega_2$ holds exactly when
    $\omega_1\omega_2^{-1}$ is a coboundary. In additive notation this ratio
    is the difference of the corresponding cochains, so the quotient
    criterion for $\pi$ gives
    (\ref{eq:mpug_h2_class_comparison}).
\end{proof}

\begin{definition}[Comparison with low-degree group cohomology]
    \label{def:mpug_h2_group_cohomology_comparison}
    \lean{TNLean.Algebra.h2EquivGroupCohomology,
        TNLean.Algebra.ScalarCocycle.trivialMulDistribMulAction,
        TNLean.Algebra.scalarH2Representation,
        TNLean.Algebra.ScalarCocycle.toMathlibCocycle,
        TNLean.Algebra.ScalarCocycle.ofMathlibCocycle,
        TNLean.Algebra.ScalarCocycle.representativesEquivMathlib,
        TNLean.Algebra.h2ToGroupCohomology,
        TNLean.Algebra.h2ToGroupCohomology_injective,
        TNLean.Algebra.h2ToGroupCohomology_surjective}
    \leanok
    \uses{def:h2_cx, thm:mpug_h2_trivial_coefficients}
    Give $\C^\times$ the trivial $G$-action. Currying and uncurrying identify
    the cocycle representatives in Definition~\ref{def:h2_cx} with the
    multiplicative representatives in the inhomogeneous cochain complex.
    By (\ref{eq:mpug_h2_class_comparison}), this correspondence induces a
    canonical equivalence
    \begin{align}
        Z^2(G,\C^\times)/{\sim}
         & \simeq H^2\bigl(G,\C^\times_{\mathrm{triv}}\bigr).
        \label{eq:mpug_h2_equivalence}
    \end{align}
    The left-hand side is the quotient already defined in
    Definition~\ref{def:h2_cx}; no additional quotient is introduced.
    This is the $H^2(H,\C^\times)$ appearing in the pair $(H,\psi)$ in
    \cite[lines~1927--1931]{FrancoRubio2025MPUGauging}.
    % Reuse boundary: the right-hand side is Mathlib's low-degree H2 for
    % Rep.ofMulDistribMulAction with the explicit trivial action.
\end{definition}

\begin{definition}[Pullback of scalar cohomology]
    \label{def:mpug_h2_pullback}
    \lean{TNLean.Algebra.ScalarCocycle.pullback,
        TNLean.Algebra.ScalarCocycle.pullback_apply,
        TNLean.Algebra.H2.pullback,
        TNLean.Algebra.H2.pullback_mk}
    \leanok
    \uses{def:h2_cx}
    Let $f:K\to G$ be a group homomorphism. For a scalar $2$-cochain
    $\omega:G\times G\to\C^\times$, define
    $(f^*\omega)(a,b)=\omega(f(a),f(b))$.
    The induced pullback map
    $f^*:H^2(G,\C^\times)\to H^2(K,\C^\times)$ is defined by
    $f^*[\omega]=[f^*\omega]$.
\end{definition}

\begin{theorem}[Pullback preserves cocycles and cohomology]
    \label{thm:mpug_h2_pullback_preserves}
    \lean{TNLean.Algebra.ScalarCocycle.IsCocycle.pullback,
        TNLean.Algebra.ScalarCocycle.CohomologousTo.pullback}
    \leanok
    \uses{def:mpug_h2_pullback, def:cohomologous_to, def:is_cocycle}
    Let $f:K\to G$ be a group homomorphism. If $\omega$ is a scalar
    $2$-cocycle on $G$, then $f^*\omega$ is a scalar $2$-cocycle on $K$.
    If $\omega_1$ and $\omega_2$ are cohomologous scalar $2$-cochains on $G$,
    then $f^*\omega_1$ and $f^*\omega_2$ are cohomologous on $K$.
\end{theorem}

\begin{proof}\leanok
    \uses{def:mpug_h2_pullback, def:cohomologous_to, def:is_cocycle}
    For $a,b,c\in K$, substitute $f(a),f(b),f(c)$ in the cocycle equation
    for $\omega$ and use $f(ab)=f(a)f(b)$ and $f(bc)=f(b)f(c)$.
    If $\omega_1(g,h)=\varphi(g)\varphi(h)\varphi(gh)^{-1}\omega_2(g,h)$,
    set $\eta=\varphi\circ f$. Substituting $g=f(a)$ and $h=f(b)$ gives
    $(f^*\omega_1)(a,b)=\eta(a)\eta(b)\eta(ab)^{-1}(f^*\omega_2)(a,b)$,
    so $\eta$ witnesses the required cohomology relation.
\end{proof}

% Partial support for #7361 (child #7754), not the full prop:BI_psi criterion.
% No U(1) restriction naturality or tensor-level block-independence claim.
\begin{theorem}[Exact extension of a subgroup cocycle]
    \label{thm:mpug_h2_exact_extension}
    \lean{TNLean.Algebra.H2.mem_range_pullback_subtype_iff,
        TNLean.Algebra.ScalarCocycle.IsCocycle.coboundary_mul}
    \leanok
    \uses{def:mpug_h2_pullback, def:is_cocycle}
    Let $H\le G$ be any subgroup, let $\iota:H\hookrightarrow G$ be its
    inclusion, and let $\psi:H\times H\to\C^\times$ be a $2$-cocycle.
    Then $[\psi]$ belongs to the image of $\iota^*$ if and only if there is
    a $2$-cocycle $\Psi:G\times G\to\C^\times$ with
    $\Psi(a,b)=\psi(a,b)$ for every $a,b\in H$.
    No normalization of either cocycle is required.
    This is the representative adjustment implicit in the extension
    condition of \cite[lines~7042--7064]{FrancoRubio2025MPUGauging},
    here stated with $\C^\times$ coefficients.
    % Source: Papers/2502.20257/main.tex:7042--7064.
\end{theorem}

\begin{proof}\leanok
    \uses{def:mpug_h2_pullback, def:cohomologous_to, def:is_cocycle}
    Suppose $[\psi]=\iota^*[\omega]$. Choose a cochain
    $\varphi:H\to\C^\times$ such that
    $\psi(a,b)=\varphi(a)\varphi(b)\varphi(ab)^{-1}\omega(a,b)$.
    Extend $\varphi$ to a function $\Phi:G\to\C^\times$, taking value $1$
    outside $H$, and set
    $\Psi(g,h)=\Phi(g)\Phi(h)\Phi(gh)^{-1}\omega(g,h)$.
    In the cocycle equation for $\Psi$, cancellation leaves the common
    factor $\Phi(g)\Phi(h)\Phi(k)\Phi(ghk)^{-1}$ multiplied by the two sides
    of the cocycle equation for $\omega$. Hence $\Psi$ is a cocycle, and
    $\Phi|_H=\varphi$ gives $\Psi|_{H\times H}=\psi$ exactly.
    Conversely, any exact extension $\Psi$ satisfies
    $\iota^*[\Psi]=[\psi]$.
\end{proof}

\begin{theorem}[Circle and nonzero-complex coefficients in degree two]
    \label{thm:mpug_circle_complex_units_h2}
    \lean{TNLean.Algebra.Complex.div_toUnits_unitsPhase,
        TNLean.Algebra.ScalarCocycle.circlePhase_isMulCocycle₂,
        TNLean.Algebra.ScalarCocycle.positiveResidual_isCoboundary,
        TNLean.Algebra.ScalarCocycle.cohomologousTo_circlePhaseInclusion,
        TNLean.Algebra.ScalarCocycle.circlePhaseMathlibCocycle,
        TNLean.Algebra.ScalarCocycle.H2π_map_circlePhaseMathlibCocycle,
        TNLean.Algebra.circleH2ToComplexUnitsH2_injective,
        TNLean.Algebra.circleH2ToComplexUnitsH2_surjective,
        TNLean.Algebra.circleH2EquivComplexUnitsH2}
    \leanok
    \uses{def:mpug_h2_group_cohomology_comparison,
        thm:mpug_positive_generalized_cocycle}
    Let $G$ be finite, and give $U(1)$ and $\C^\times$ the trivial
    $G$-actions.  Inclusion of the unit circle induces an equivalence
    \begin{align}
        H^2(G,U(1))
         & \simeq H^2(G,\C^\times).
        \label{eq:mpug_circle_complex_units_h2}
    \end{align}
    More concretely, for a $\C^\times$-valued cocycle $\omega$, define its
    pointwise phase and radial residual by
    \begin{align}
        u(g,h)
         & =\frac{\omega(g,h)}{|\omega(g,h)|}\in U(1),
         &
        r(g,h)
         & =\frac{\omega(g,h)}{u(g,h)}=|\omega(g,h)|.
        \label{eq:mpug_circle_phase_residual}
    \end{align}
    Then $u$ is a cocycle, $r$ is a positive coboundary, and hence $\omega$
    is cohomologous to the inclusion of $u$.  The inclusion map is injective
    for every group because pointwise phase is its left inverse; finiteness is
    used only to prove surjectivity.

    This is an algebraic group-cohomology statement for a finite discrete
    group.  The footnote at
    \cite[line~1931]{FrancoRubio2025MPUGauging} instead invokes a comparison
    of sheaf cohomology with topological coefficient groups; the distinction
    and the direct finite-group proof are recorded in
    \cite{gap:fbc25_circle_complex_units_cohomology}.
\end{theorem}

\begin{proof}\leanok
    \uses{thm:mpug_positive_generalized_cocycle}
    The phase map $q(z)=z/|z|$ is multiplicative and satisfies
    $q(u)=u$ for $u\in U(1)$, so it preserves the cocycle equation and
    retracts inclusion.  The polar identity in
    (\ref{eq:mpug_circle_phase_residual}) identifies the residual $r$ with a
    positive cocycle.  Apply
    Theorem~\ref{thm:mpug_positive_generalized_cocycle} to the one-point block
    space to write $r=d\chi$.  Therefore $\omega\sim u$, which proves
    surjectivity of coefficient inclusion on degree-two classes.  Its left
    inverse induced by $q$ proves injectivity.
\end{proof}

\begin{definition}[Cocycle-twisted regular matrices]
    \label{def:mpug_twisted_regular_matrices}
    \lean{TNLean.Algebra.ScalarCocycle.phase,
        TNLean.Algebra.ScalarCocycle.twistedLeftRegular,
        TNLean.Algebra.ScalarCocycle.twistedRightRegular}
    \leanok
    \uses{def:scalar_cocycle, thm:mpug_circle_complex_units_h2}
    Let $G$ be a group and let $\omega$ be a $\C^\times$-valued scalar
    cochain.  Write
    \begin{align}
        \phi(g,h) & =\frac{\omega(g,h)}{|\omega(g,h)|}\in U(1).
    \end{align}
    On the standard basis of $\C^G$, define
    \begin{align}
        L_g^\phi
         & =\sum_{h\in G}\phi(g,h)^{-1}|gh\rangle\langle h|,
         &
        R_g^\phi
         & =\sum_{h\in G}\phi(h,g^{-1})^{-1}|hg^{-1}\rangle\langle h|.
        \label{eq:mpug_twisted_regular_matrices}
    \end{align}
    The source uses a $U(1)$-valued cocycle $\Omega$; taking $\Omega=\phi$
    gives exactly the leg and inverse conventions in
    \cite[lines~412--418]{FrancoRubio2025MPUGauging}.
\end{definition}

\begin{lemma}[Cocycle equation for the circle phase]
    \label{lem:mpug_circle_phase_cocycle}
    \lean{TNLean.Algebra.ScalarCocycle.phase_isCocycle}
    \leanok
    \uses{def:is_cocycle, thm:mpug_circle_complex_units_h2}
    If $\omega$ is a cocycle, then its circle phase satisfies
    \begin{align}
        \phi(g,h)\phi(gh,k)
         & =\phi(g,hk)\phi(h,k)
    \end{align}
    for all $g,h,k\in G$.
\end{lemma}

\begin{proof}\leanok
    \uses{thm:mpug_circle_complex_units_h2}
    The pointwise phase map $z\mapsto z/|z|$ is multiplicative, so it preserves
    the cocycle equation.
\end{proof}

\begin{theorem}[Action of the twisted regular matrices]
    \label{thm:mpug_twisted_regular_action}
    \lean{TNLean.Algebra.ScalarCocycle.twistedLeftRegular_mulVec_apply,
        TNLean.Algebra.ScalarCocycle.twistedRightRegular_mulVec_apply}
    \leanok
    \uses{def:mpug_twisted_regular_matrices}
    If $G$ is finite, then for $v\in\C^G$ and $g,h\in G$,
    \begin{align}
        (L_g^\phi v)(gh)
         & =\phi(g,h)^{-1}v(h),
         &
        (R_g^\phi v)(hg^{-1})
         & =\phi(h,g^{-1})^{-1}v(h).
    \end{align}
    % Source anchor: arXiv:2502.20257, lines 412--418.
\end{theorem}

\begin{proof}\leanok
    \uses{def:mpug_twisted_regular_matrices}
    In each matrix-vector product, exactly one summand is nonzero: the column
    indexed by $h$.
\end{proof}

\begin{theorem}[Multiplication laws for twisted regular matrices]
    \label{thm:mpug_twisted_regular_multiplication}
    \lean{TNLean.Algebra.ScalarCocycle.twistedLeftRegular_mul,
        TNLean.Algebra.ScalarCocycle.twistedRightRegular_mul}
    \leanok
    \uses{def:mpug_twisted_regular_matrices,
        lem:mpug_circle_phase_cocycle}
    If $G$ is finite and $\omega$ is a cocycle, then
    \begin{align}
        L_g^\phi L_h^\phi
         & =\phi(g,h)^{-1}L_{gh}^\phi,           \\
        R_g^\phi R_h^\phi
         & =\phi(h^{-1},g^{-1})^{-1}R_{gh}^\phi.
        \label{eq:mpug_twisted_regular_multiplication}
    \end{align}
    % Source anchor: arXiv:2502.20257, lines 406--418; the laws follow from
    % the displayed twisted regular matrices and the cocycle equation.
\end{theorem}

\begin{proof}\leanok
    \uses{lem:mpug_circle_phase_cocycle}
    Expand the matrix products entrywise.  The two scalar identities are the
    cocycle equation at $(g,h,k)$ and $(k,h^{-1},g^{-1})$, respectively.
\end{proof}

\begin{theorem}[Commutation of twisted left and right actions]
    \label{thm:mpug_twisted_regular_commutation}
    \lean{TNLean.Algebra.ScalarCocycle.twistedLeftRegular_commute}
    \leanok
    \uses{def:mpug_twisted_regular_matrices}
    If $G$ is finite and $\omega$ is a cocycle, then
    \begin{align}
        L_g^\phi R_h^\phi
         & =R_h^\phi L_g^\phi.
        \label{eq:mpug_twisted_regular_commutation}
    \end{align}
    % Source anchor: arXiv:2502.20257, lines 406--418 and 457; the latter
    % explicitly attributes commuting local symmetries to commutativity of the
    % left and right group actions.
\end{theorem}

\begin{proof}\leanok
    \uses{lem:mpug_circle_phase_cocycle}
    Expand both matrix products entrywise and use the cocycle equation at
    $(g,k,h^{-1})$.
\end{proof}

\begin{theorem}[Unitarity of twisted regular matrices]
    \label{thm:mpug_twisted_regular_unitary}
    \lean{TNLean.Algebra.ScalarCocycle.twistedLeftRegular_mem_unitaryGroup,
        TNLean.Algebra.ScalarCocycle.twistedRightRegular_mem_unitaryGroup}
    \leanok
    \uses{def:mpug_twisted_regular_matrices}
    If $G$ is finite, the matrices $L_g^\phi$ and $R_g^\phi$ are unitary for
    every $g\in G$.
    % Source anchor: arXiv:2502.20257, lines 406--418, where these matrices
    % form the two gauge legs of the local unitary representation.
\end{theorem}

\begin{proof}\leanok
    \uses{def:mpug_twisted_regular_matrices}
    Each matrix is a permutation matrix whose nonzero entries have modulus
    one.  Distinct columns therefore have disjoint support, while every column
    has squared norm one.
\end{proof}

\begin{theorem}[Onsite Gauss coefficient identity]
    \label{thm:mpug_onsite_gauss_phase}
    \lean{TNLean.Algebra.ScalarCocycle.onsiteGauss_phase_div}
    \leanok
    \uses{def:mpug_twisted_regular_matrices,
        lem:mpug_circle_phase_cocycle}
    Suppose that $\omega$ is a cocycle and impose the additional normalization
    \begin{align}
        \phi(t^{-1},t) & =1
    \end{align}
    for every $t\in G$.  Then
    \begin{align}
        \frac{\phi(ag^{-1},gb)}{\phi(a,b)}
         & =\frac{1}{\phi(a,g^{-1})\phi(g,b)}.
        \label{eq:mpug_onsite_gauss_phase}
    \end{align}
    The scalar convention $S_{t^{-1}}=S_t^{-1}$ appears in
    \cite[line~238]{FrancoRubio2025MPUGauging}, and the coefficient rewrite is
    at lines 3763--3766. The displayed normalization remains an explicit
    assumption here; it is not a consequence of cocyclicity.
\end{theorem}

\begin{proof}\leanok
    \uses{lem:mpug_circle_phase_cocycle}
    Use the cocycle equation first at $(a,g^{-1},gb)$ and then at
    $(g^{-1},g,b)$.  The normalization at $t=e$ also gives
    $\phi(e,b)=1$.
\end{proof}

\begin{definition}[Block scalar and block independence]
    \label{def:mpug_block_independence_factor}
    \lean{TNLean.Algebra.LSymbol.ell,
        TNLean.Algebra.LSymbol.IsBlockIndependentEll}
    \leanok
    \uses{def:mpug_l_symbols}
    For an $L$-symbol, define
    \begin{align}
        \ell^x_{a,b;r}
         & =\frac{L^x_{ar,r^{-1}b}}{L^x_{a,b}}.
        \label{eq:mpug_block_scalar}
    \end{align}
    The scalar $\ell$ is block-independent when
    \begin{align}
        \ell^x_{a,b;r}
         & =\ell^y_{a,b;r}
        \label{eq:mpug_block_independence}
    \end{align}
    for all $a,b,r\in G$ and $x,y\in\mathsf X$.
    % Source anchor: arXiv:2502.20257, app:block_indep, lines 6965--6974.
\end{definition}

\begin{definition}[Four scalar gauge formulations]
    \label{def:mpug_block_independence_gauge_forms}
    \lean{TNLean.Algebra.LSymbol.IsBlockConstant,
        TNLean.Algebra.LSymbol.IsTrivialEll,
        TNLean.Algebra.LSymbol.IsTrivial,
        TNLean.Algebra.LSymbol.HasTrivialGauge,
        TNLean.Algebra.LSymbol.HasBlockConstantGauge,
        TNLean.Algebra.LSymbol.HasTrivialEllGauge,
        TNLean.Algebra.LSymbol.IsBlockIndependent}
    \leanok
    \uses{def:mpug_l_gauge, def:mpug_block_independence_factor}
    These are the four scalar $L$-symbol properties used in the proof of the
    source proposition. Write
    $L'=(\beta,\gamma)\mathbin{\triangleright}L$ for a joint scalar gauge of a
    supplied $L$-symbol. The four properties are, respectively, the existence
    of $\beta$ and $\gamma$ such that
    \begin{align}
        (L')^x_{g,h}       & =1
                           &                     & \text{for all $x,g,h$},
        \label{eq:mpug_block_form_trivial}                                     \\
        (L')^x_{a,b}       & =(L')^y_{a,b}
                           &                     & \text{for all $x,y,a,b$},
        \label{eq:mpug_block_form_constant}                                    \\
        \ell[L']^x_{a,b;r} & =1
                           &                     & \text{for all $x,a,b,r$},
        \label{eq:mpug_block_form_trivial_ell}                                 \\
        \ell[L']^x_{a,b;r} & =\ell[L']^y_{a,b;r}
                           &                     & \text{for all $x,y,a,b,r$}.
        \label{eq:mpug_block_form_independent_ell}
    \end{align}
    % Source anchor: arXiv:2502.20257, prop:technical01, lines 6975--6989.
\end{definition}

\begin{theorem}[Corrected block factorization]
    \label{thm:mpug_block_factorization}
    \lean{TNLean.Algebra.LSymbol.IsCompatible.apply_right_one,
        TNLean.Algebra.LSymbol.blockFactor,
        TNLean.Algebra.LSymbol.factorization_of_isBlockIndependentEll}
    \leanok
    \uses{def:mpug_l_symbols, def:mpug_block_independence_factor}
    Suppose that $L$ is compatible with $\omega$ and that $\ell$ is
    block-independent. Fix $x_0\in\mathsf X$ and set
    \begin{align}
        A(g,h)
         & =\bigl(\omega(gh,e,e)\ell^{x_0}_{g,h;h}\bigr)^{-1}.
        \label{eq:mpug_block_factor}
    \end{align}
    Then, for every $x\in\mathsf X$ and $g,h\in G$,
    \begin{align}
        L^x_{g,h}
         & =A(g,h)L^x_{e,e}.
        \label{eq:mpug_corrected_block_factorization}
    \end{align}
    The reciprocal factor in this identity is recorded in
    \cite{gap:fbc25_block_independence_reciprocal_errors}.
    % Local correction:
    % docs/paper-gaps/fbc25_block_independence_reciprocal_errors.tex.
    % Source anchor: arXiv:2502.20257, eq:Lfact, lines 6994--7002.
\end{theorem}

\begin{proof}\leanok
    \uses{def:mpug_block_independence_factor}
    Definition~\ref{def:mpug_block_independence_factor} gives
    \begin{align}
        L^x_{g,h}
         & =\frac{L^x_{gh,e}}{\ell^x_{g,h;h}}.
        \label{eq:mpug_block_factorization_first}
    \end{align}
    Block independence replaces $x$ by $x_0$ in the denominator.
    Compatibility at $(gh,e,e)$ gives
    $L^x_{gh,e}=L^x_{e,e}/\omega(gh,e,e)$, which proves
    (\ref{eq:mpug_corrected_block_factorization}).
\end{proof}

\begin{theorem}[Relative block character]
    \label{thm:mpug_relative_block_character}
    \lean{TNLean.Algebra.LSymbol.IsCompatible.apply_left_one,
        TNLean.Algebra.LSymbol.relativeCharacter,
        TNLean.Algebra.LSymbol.relativeCharacter_mul_base,
        TNLean.Algebra.LSymbol.relativeCharacter_mul}
    \leanok
    \uses{thm:mpug_block_factorization}
    Under the hypotheses of Theorem~\ref{thm:mpug_block_factorization}, define
    \begin{align}
        \rho(g)
         & =\bigl(\ell^{x_0}_{e,g;g}\omega(g,e,e)
        \omega(e,e,g)\bigr)^{-1}.
        \label{eq:mpug_relative_block_character}
    \end{align}
    The reciprocal factor $(\ell^{x_0}_{e,g;g})^{-1}$ corrects the
    corresponding printed factor \cite{gap:fbc25_block_independence_reciprocal_errors}.
    Then
    \begin{align}
        L^{g x}_{e,e}
         & =\rho(g)L^x_{e,e},
        \label{eq:mpug_relative_block_transport} \\
        \rho(gh)
         & =\rho(g)\rho(h).
        \label{eq:mpug_relative_block_multiplicative}
    \end{align}
    Thus $\rho\colon G\to\C^\times$ is a one-dimensional character. No
    transitivity or finiteness assumption on the action is required.
    % Local correction:
    % docs/paper-gaps/fbc25_block_independence_reciprocal_errors.tex.
    % Source anchor: arXiv:2502.20257, eq:defrho, lines 7003--7018.
\end{theorem}

\begin{proof}\leanok
    \uses{thm:mpug_block_factorization}
    Compatibility at $(e,e,g)$ gives
    $L^x_{e,g}=\omega(e,e,g)L^{g x}_{e,e}$. Compare this with
    (\ref{eq:mpug_corrected_block_factorization}) at $(e,g)$ to obtain
    (\ref{eq:mpug_relative_block_transport}). Applying that identity first
    to $h$ and then to $g$, and comparing with its application to $gh$, proves
    (\ref{eq:mpug_relative_block_multiplicative}) by cancellation.
\end{proof}

\begin{theorem}[Explicit inverse gauge]
    \label{thm:mpug_block_independence_inverse_gauge}
    \lean{TNLean.Algebra.LSymbol.factorFusionGauge,
        TNLean.Algebra.LSymbol.factorActionGauge,
        TNLean.Algebra.LSymbol.gauge_one_eq_of_isBlockIndependentEll,
        TNLean.Algebra.LSymbol.gauge_inv_eq_one_of_isBlockIndependentEll,
        TNLean.Algebra.LSymbol.hasTrivialGauge_of_isBlockIndependentEll}
    \leanok
    \uses{def:mpug_l_gauge, thm:mpug_l_gauge_composition,
        thm:mpug_block_factorization, thm:mpug_relative_block_character}
    With $A$ and $\rho$ as above, define the forward gauges by
    \begin{align}
        \beta^+_{g,h}
         & =\frac{A(g,h)}{\rho(h)},
        \label{eq:mpug_forward_fusion_gauge} \\
        \gamma^{+,x}_g
         & =\bigl(L^{g x}_{e,e}\bigr)^{-1}.
        \label{eq:mpug_forward_action_gauge}
    \end{align}
    Their joint gauge multiplier equals $L$:
    \begin{align}
        (\beta^+,\gamma^+)\mathbin{\triangleright}1
         & =L.
        \label{eq:mpug_forward_gauge_multiplier}
    \end{align}
    Consequently the pointwise inverse gauges trivialize $L$:
    \begin{align}
        ((\beta^+)^{-1},(\gamma^+)^{-1})
        \mathbin{\triangleright}L
         & =1.
        \label{eq:mpug_inverse_gauge_trivialization}
    \end{align}
    The distinction between the forward multiplier and its pointwise inverse
    corrects the orientation of the displayed source calculation
    \cite{gap:fbc25_block_independence_reciprocal_errors}.
    If $\mathsf X$ is nonempty, a choice of $x_0$ therefore gives a
    trivializing joint scalar gauge for every compatible block-independent
    $L$-symbol.
    % Local correction:
    % docs/paper-gaps/fbc25_block_independence_reciprocal_errors.tex.
    % Source anchor: arXiv:2502.20257, eq:Lgauge and lines 7019--7027.
\end{theorem}

\begin{proof}\leanok
    \uses{thm:mpug_l_gauge_composition, thm:mpug_block_factorization,
        thm:mpug_relative_block_character}
    Substitute (\ref{eq:mpug_forward_fusion_gauge}) and
    (\ref{eq:mpug_forward_action_gauge}) into
    (\ref{eq:mpug_l_gauge}). Equation
    (\ref{eq:mpug_relative_block_transport}) reduces the multiplier to
    $A(g,h)L^x_{e,e}$, which equals $L^x_{g,h}$ by
    (\ref{eq:mpug_corrected_block_factorization}). Joint gauges compose by
    pointwise multiplication, so the inverse gauges give
    (\ref{eq:mpug_inverse_gauge_trivialization}).
\end{proof}

\begin{theorem}[Trivial and block-constant gauges]
    \label{thm:mpug_trivial_iff_block_constant}
    \lean{TNLean.Algebra.LSymbol.IsTrivial.isBlockConstant,
        TNLean.Algebra.LSymbol.IsBlockConstant.isBlockIndependentEll,
        TNLean.Algebra.LSymbol.HasTrivialGauge.hasBlockConstantGauge,
        TNLean.Algebra.LSymbol.HasBlockConstantGauge.isBlockIndependent,
        TNLean.Algebra.LSymbol.hasTrivialGauge_iff_hasBlockConstantGauge}
    \leanok
    \uses{def:mpug_block_independence_gauge_forms,
        thm:mpug_l_gauge_composition, thm:mpug_l_gauge_preservation,
        thm:mpug_block_independence_inverse_gauge}
    Let $L$ be a supplied $L$-symbol compatible with a scalar $3$-cochain,
    and suppose that $\mathsf X$ is nonempty. A joint scalar gauge trivializes
    $L$ if and only if a joint scalar gauge makes $L^x_{g,h}$ independent of
    $x$.
    % Source anchor: arXiv:2502.20257, prop:technical01, items 1 and 2.
\end{theorem}

\begin{proof}\leanok
    \uses{thm:mpug_l_gauge_composition, thm:mpug_l_gauge_preservation,
        thm:mpug_block_independence_inverse_gauge}
    A trivial $L$-symbol is block-constant. Conversely, a block-constant
    representative has block-independent $\ell$. Apply
    Theorem~\ref{thm:mpug_block_independence_inverse_gauge} to that
    representative and compose the two joint gauges.
\end{proof}

\begin{theorem}[Trivial gauges and trivial block scalars]
    \label{thm:mpug_trivial_iff_trivial_ell}
    \lean{TNLean.Algebra.LSymbol.IsTrivial.isTrivialEll,
        TNLean.Algebra.LSymbol.IsTrivialEll.isBlockIndependentEll,
        TNLean.Algebra.LSymbol.HasTrivialGauge.hasTrivialEllGauge,
        TNLean.Algebra.LSymbol.HasTrivialEllGauge.isBlockIndependent,
        TNLean.Algebra.LSymbol.hasTrivialGauge_iff_hasTrivialEllGauge}
    \leanok
    \uses{def:mpug_block_independence_gauge_forms,
        thm:mpug_l_gauge_composition, thm:mpug_l_gauge_preservation,
        thm:mpug_block_independence_inverse_gauge}
    Let $L$ be a supplied $L$-symbol compatible with a scalar $3$-cochain,
    and suppose that $\mathsf X$ is nonempty. A joint scalar gauge trivializes
    $L$ if and only if a joint scalar gauge makes every
    $\ell^x_{a,b;r}$ equal to $1$.
    % Source anchor: arXiv:2502.20257, prop:technical01, items 1 and 3.
\end{theorem}

\begin{proof}\leanok
    \uses{thm:mpug_l_gauge_composition, thm:mpug_l_gauge_preservation,
        thm:mpug_block_independence_inverse_gauge}
    A trivial $L$-symbol has $\ell=1$. Conversely, the identity $\ell=1$
    implies block independence. Apply
    Theorem~\ref{thm:mpug_block_independence_inverse_gauge} and compose the
    gauges.
\end{proof}

\begin{theorem}[Scalar gauge triviality and block independence]
    \label{thm:mpug_trivial_iff_block_independent}
    \lean{TNLean.Algebra.LSymbol.IsBlockIndependent.hasTrivialGauge,
        TNLean.Algebra.LSymbol.hasTrivialGauge_iff_isBlockIndependent}
    \leanok
    \uses{def:mpug_block_independence_gauge_forms,
        thm:mpug_l_gauge_composition, thm:mpug_l_gauge_preservation,
        thm:mpug_block_independence_inverse_gauge}
    Let $L$ be a supplied $L$-symbol compatible with a scalar $3$-cochain,
    and suppose that $\mathsf X$ is nonempty. A joint scalar gauge trivializes
    $L$ if and only if a joint scalar gauge makes $\ell^x_{a,b;r}$ independent
    of $x$.
    % Source anchor: arXiv:2502.20257, prop:technical01, items 1 and 4.
\end{theorem}

\begin{proof}\leanok
    \uses{thm:mpug_l_gauge_composition, thm:mpug_l_gauge_preservation,
        thm:mpug_block_independence_inverse_gauge}
    Gauge triviality implies block constancy and hence block independence.
    For the converse, apply
    Theorem~\ref{thm:mpug_block_independence_inverse_gauge} to the
    block-independent representative and compose its trivializing gauge with
    the original gauge.
\end{proof}

\begin{theorem}[Four scalar gauge formulations of block independence]
    \label{thm:mpug_block_independence_four_forms}
    \lean{TNLean.Algebra.LSymbol.hasTrivialGauge_tfae}
    \leanok
    \uses{def:mpug_block_independence_gauge_forms,
        thm:mpug_trivial_iff_block_constant,
        thm:mpug_trivial_iff_trivial_ell,
        thm:mpug_trivial_iff_block_independent}
    Let $L$ be a supplied $L$-symbol compatible with a scalar $3$-cochain,
    and suppose that $\mathsf X$ is nonempty. The four scalar gauge
    properties (\ref{eq:mpug_block_form_trivial}),
    (\ref{eq:mpug_block_form_constant}),
    (\ref{eq:mpug_block_form_trivial_ell}), and
    (\ref{eq:mpug_block_form_independent_ell}) of $L$ are equivalent. This is
    the scalar core of the block-independence proposition of
    \cite{FrancoRubio2025MPUGauging}, whose fourth condition is stated for an
    MPS--MPU pair.
    % Source anchor: arXiv:2502.20257, prop:technical01, lines 6975--6989.
\end{theorem}

\begin{proof}\leanok
    \uses{thm:mpug_trivial_iff_block_constant,
        thm:mpug_trivial_iff_trivial_ell,
        thm:mpug_trivial_iff_block_independent}
    Theorems~\ref{thm:mpug_trivial_iff_block_constant},
    \ref{thm:mpug_trivial_iff_trivial_ell}, and
    \ref{thm:mpug_trivial_iff_block_independent} identify each of the last
    three properties with the first.
\end{proof}

The source's fourth condition is a statement about an MPS--MPU pair. The
results above begin instead with a supplied compatible scalar $L$-symbol and
prove the algebraic equivalence among its four scalar gauge properties. A
separate tensor-level statement must construct the associated $L$-symbol and
show that the tensor gauge freedoms induce the joint scalar gauge
(\ref{eq:mpug_l_gauge}).
% Scope boundary: the tensor-facing bridge for arXiv:2502.20257,
% prop:technical01 is not contained in the scalar L-symbol declarations.

\begin{theorem}[Stabilizer transition elements]
    \label{thm:mpug_stabilizer_transition}
    \lean{TNLean.Algebra.StabilizerRepresentatives.transitionElement_mul,
        TNLean.Algebra.StabilizerRepresentatives,
        TNLean.Algebra.StabilizerRepresentatives.ofPretransitive,
        TNLean.Algebra.StabilizerRepresentatives.transitionElement,
        TNLean.Algebra.StabilizerRepresentatives.coe_transitionElement,
        TNLean.Algebra.StabilizerRepresentatives.transitionElement_stabilizer}
    \leanok
    Suppose that $G$ acts transitively on $\mathsf X$, and fix
    $x_0\in\mathsf X$. There are representatives $k_x\in G$ satisfying
    $k_xx_0=x$ and $k_{x_0}=e$. With
    $H=\operatorname{Stab}_G(x_0)$, define
    \begin{align}
        t_k(g,x)
         & =k_{g x}^{-1}gk_x.
        \label{eq:mpug_stabilizer_transition}
    \end{align}
    Then $t_k(g,x)\in H$ and
    \begin{align}
        t_k(g_1g_2,x)
         & =t_k(g_1,g_2x)t_k(g_2,x).
        \label{eq:mpug_stabilizer_transition_product}
    \end{align}
    Consequently, the source factors are
    \begin{align}
        h_2
         & =t_k(g_2,x)=k_{g_2x}^{-1}g_2k_x,
        \label{eq:mpug_stabilizer_h2}                  \\
        h_1
         & =t_k(g_1,g_2x)=k_{g_1g_2x}^{-1}g_1k_{g_2x}.
        \label{eq:mpug_stabilizer_h1}
    \end{align}
    The last representative in $h_1$ corrects the printed $k_{g_1x}$
    \cite{gap:fbc25_stabilizer_transition_index_typo}. No finiteness or
    normality assumption is required.
    % Local correction:
    % docs/paper-gaps/fbc25_stabilizer_transition_index_typo.tex.
    % Source anchor: arXiv:2502.20257, lemma:h1h2 and eq:h1h2,
    % lines 7028--7042.
\end{theorem}

\begin{proof}\leanok
    Transitivity supplies one element carrying $x_0$ to each $x$; choose the
    identity at $x=x_0$. Acting with (\ref{eq:mpug_stabilizer_transition}) on
    $x_0$ gives
    $k_{gx}^{-1}g(k_xx_0)=k_{gx}^{-1}(gx)=x_0$.
    For (\ref{eq:mpug_stabilizer_transition_product}), substitute
    (\ref{eq:mpug_stabilizer_transition}) and cancel
    $k_{g_2x}k_{g_2x}^{-1}$. The two formulas for $h_1$ and $h_2$ are the
    corresponding specializations.
\end{proof}

\begin{theorem}[L-symbol induced by a stabilizer cocycle]
    \label{thm:mpug_stabilizer_cocycle_l_symbol}
    \lean{TNLean.Algebra.StabilizerRepresentatives.inducedLSymbol_isCompatible,
        TNLean.Algebra.StabilizerRepresentatives.inducedLSymbol,
        TNLean.Algebra.StabilizerRepresentatives.inducedLSymbol_base,
        TNLean.Algebra.StabilizerRepresentatives.inducedActionGauge,
        TNLean.Algebra.StabilizerRepresentatives.inducedLSymbol_eq_gauge,
        TNLean.Algebra.StabilizerRepresentatives.cohomologousTo_iff_exists_actionGauge_inducedLSymbol_eq}
    \leanok
    \uses{def:mpug_l_symbols, def:mpug_l_gauge, def:is_cocycle,
        def:cohomologous_to, thm:mpug_stabilizer_transition}
    Let $H=\operatorname{Stab}_G(x_0)$ and let
    $\psi\colon H\times H\to\C^\times$ be a scalar $2$-cocycle. Define
    \begin{align}
        L^x_{g,h}
         & =\psi\bigl(t_k(g,hx),t_k(h,x)\bigr).
        \label{eq:mpug_induced_stabilizer_l_symbol}
    \end{align}
    Formula (\ref{eq:mpug_induced_stabilizer_l_symbol}) is the $\omega=1$
    scalar specialization of \cite[Eq.~(20)]{GarreRubio2023MPOSym}.
    The $L$-symbol satisfies (\ref{eq:mpug_l_compatibility}) with $\omega=1$,
    and its base-point restriction recovers $\psi$:
    \begin{align}
        L^{x_0}_{a,b} & =\psi(a,b)
        \label{eq:mpug_induced_stabilizer_base}
    \end{align}
    for all $a,b\in H$.

    More precisely, two stabilizer cocycles are cohomologous exactly when
    their induced $L$-symbols differ by an action-tensor gauge and the trivial
    fusion-tensor gauge:
    \begin{align}
        \psi_1\sim\psi_2
         & \quad\Longleftrightarrow\quad
        \exists\gamma,\quad
        L[\psi_1]=(1,\gamma)\mathbin{\triangleright}L[\psi_2].
        \label{eq:mpug_induced_cohomology_iff}
    \end{align}
    If the cohomology relation is witnessed by
    $\varphi\colon H\to\C^\times$, with
    \begin{align}
        \psi_1(a,b)
         & =\varphi(a)\varphi(b)\varphi(ab)^{-1}\psi_2(a,b),
        \label{eq:mpug_stabilizer_cocycle_gauge}
    \end{align}
    then the forward action-tensor gauge is
    \begin{align}
        \gamma_{g,x} & =\varphi\bigl(t_k(g,x)\bigr)^{-1}.
        \label{eq:mpug_induced_action_gauge}
    \end{align}
    Thus the induced action-gauge class depends only on the cohomology class
    of $\psi$.

    The representation-level lemma of
    \cite[lines~7028--7042]{FrancoRubio2025MPUGauging} also asserts that a
    tensor-derived $L$-symbol can be put in this form. The result here is only
    its scalar part and does not make that representation-theoretic claim.
    % Source anchor: arXiv:2203.12563, Eq. (20).
\end{theorem}

\begin{proof}\leanok
    \uses{def:mpug_l_gauge, thm:mpug_stabilizer_transition}
    Put
    $a=t_k(g,hkx)$, $b=t_k(h,kx)$, and $c=t_k(k,x)$.
    The product law (\ref{eq:mpug_stabilizer_transition_product}) gives
    $t_k(gh,kx)=ab$ and $t_k(hk,x)=bc$. Compatibility is therefore exactly
    \begin{align}
        \psi(a,b)\psi(ab,c)
         & =\psi(a,bc)\psi(b,c),
    \end{align}
    which is the cocycle equation. If $p,q\in H$, then
    $t_k(p,x_0)=p$ and $t_k(q,x_0)=q$ because $k_{x_0}=e$, proving
    (\ref{eq:mpug_induced_stabilizer_base}). For the forward gauge
    calculation, set $a=t_k(g,hx)$ and $b=t_k(h,x)$, so that
    $t_k(gh,x)=ab$. Substituting (\ref{eq:mpug_induced_action_gauge}) into
    (\ref{eq:mpug_l_gauge}) gives the multiplier
    \begin{align}
        \frac{\gamma_{gh,x}}{\gamma_{g,hx}\gamma_{h,x}}
         & =\varphi(ab)^{-1}
        \bigl(\varphi(a)^{-1}\varphi(b)^{-1}\bigr)^{-1}
        =\varphi(a)\varphi(b)\varphi(ab)^{-1}.
    \end{align}
    Formula (\ref{eq:mpug_stabilizer_cocycle_gauge}) now identifies the two
    induced $L$-symbols. Conversely, suppose
    $L[\psi_1]=(1,\gamma)\mathbin{\triangleright}L[\psi_2]$ and define
    $\varphi(a)=\gamma_{a,x_0}^{-1}$ for $a\in H$. Evaluating the gauge
    equality at $x_0$ and $a,b\in H$ gives
    \begin{align}
        \psi_1(a,b)
         & =\frac{\gamma_{ab,x_0}}
        {\gamma_{a,x_0}\gamma_{b,x_0}}\psi_2(a,b)
        =\varphi(a)\varphi(b)\varphi(ab)^{-1}\psi_2(a,b).
    \end{align}
    Hence $\psi_1\sim\psi_2$, which proves the reverse implication in
    (\ref{eq:mpug_induced_cohomology_iff}).
\end{proof}

\begin{definition}[Generalized coboundary and cocycle]
    \label{def:mpug_generalized_cocycle}
    \lean{TNLean.Algebra.GeneralizedThreeCochain,
        TNLean.Algebra.GeneralizedFourCochain,
        TNLean.Algebra.ActionTensorGauge.coboundary,
        TNLean.Algebra.LSymbol.coboundary,
        TNLean.Algebra.GeneralizedThreeCochain.coboundary,
        TNLean.Algebra.LSymbol.IsGeneralizedCocycle}
    \leanok
    \uses{def:mpug_l_symbols, def:mpug_l_gauge}
    Let $\chi\colon G\times\mathsf X\to\C^\times$ be block-dependent. Its
    generalized coboundary is
    \begin{align}
        (d\chi)^x_{g,h}
         & =\frac{\chi^{h x}_g\chi^x_h}{\chi^x_{gh}}.
        \label{eq:mpug_generalized_coboundary}
    \end{align}
    For a block-dependent scalar $2$-cochain
    $f\colon\mathsf X\times G^2\to\C^\times$, set
    \begin{align}
        (df)^x_{g,h,k}
         & =\frac{f^x_{g,hk}f^x_{h,k}}
        {f^{k x}_{g,h}f^x_{gh,k}}.
        \label{eq:mpug_generalized_two_coboundary}
    \end{align}
    The final group index in the denominator is $k$; this local correction
    of the printed formula is recorded in
    \cite{gap:fbc25_two_cochain_coboundary_index_typo}. For a block-dependent
    scalar $3$-cochain
    $\omega\colon\mathsf X\times G^3\to\C^\times$, the corresponding next
    coboundary is
    \begin{align}
        (d\omega)^x_{g,h,k,l}
         & =\frac{\omega^{l x}_{g,h,k}\omega^x_{g,hk,l}
            \omega^x_{h,k,l}}
        {\omega^x_{gh,k,l}\omega^x_{g,h,kl}}.
        \label{eq:mpug_generalized_three_coboundary}
    \end{align}
    A function $\Omega\colon\mathsf X\times G^2\to\C^\times$ is a
    generalized $2$-cocycle when
    \begin{align}
        \Omega^x_{g,hk}\Omega^x_{h,k}
         & =\Omega^{k x}_{g,h}\Omega^x_{gh,k}.
        \label{eq:mpug_generalized_cocycle}
    \end{align}
    The source gives the first two coboundaries and the equation $d\Omega=1$
    in \cite[lines~5914--5922 and Lemma~\texttt{lemma:G\_cocycle}]
    {FrancoRubio2025MPUGauging}. Equation
    \eqref{eq:mpug_generalized_three_coboundary} spells out the
    three-to-four-cochain instance needed to read its assertion $d(df)=1$.
    No transitivity or finiteness assumption is imposed on $\mathsf X$.

    % Local fix: docs/paper-gaps/fbc25_two_cochain_coboundary_index_typo.tex.
    % Source anchor: arXiv:2502.20257, lines 5914--5922 and lemma:G_cocycle.
\end{definition}

\begin{theorem}[Generalized coboundaries square to one]
    \label{thm:mpug_generalized_coboundary_squared}
    \lean{TNLean.Algebra.ActionTensorGauge.isGeneralizedCocycle_coboundary,
        TNLean.Algebra.GeneralizedThreeCochain.coboundary_coboundary}
    \leanok
    \uses{def:mpug_generalized_cocycle}
    The generalized coboundary of every block-dependent scalar $1$-cochain
    is a generalized $2$-cocycle. For every block-dependent scalar
    $2$-cochain $f$,
    \begin{align}
        d(df) & =1.
        \label{eq:mpug_generalized_coboundary_squared}
    \end{align}
    This is the identity stated in
    \cite[line~5922]{FrancoRubio2025MPUGauging}, with the index correction in
    \cite{gap:fbc25_two_cochain_coboundary_index_typo}.
\end{theorem}

\begin{proof}\leanok
    \uses{def:mpug_generalized_cocycle}
    Substitute the coboundary formulas. The action law
    $k(lx)=(kl)x$ and associativity identify the shifted block labels and
    group products. Every value of the original cochain then occurs once in
    the numerator and once in the denominator, so all factors cancel.
\end{proof}

\begin{definition}[Positive block-dependent cochains]
    \label{def:mpug_positive_generalized_cochains}
    \lean{TNLean.Algebra.PositiveUnits.embedding,
        TNLean.Algebra.PositiveUnits.IsPositive,
        TNLean.Algebra.ActionTensorGauge.IsPositiveValued,
        TNLean.Algebra.LSymbol.IsPositiveValued}
    \leanok
    \uses{def:mpug_l_symbols, def:mpug_l_gauge}
    Let $\iota\colon\mathbb R_{>0}\to\C^\times$ be the natural
    multiplicative embedding. A block-dependent scalar $1$-cochain $\chi$ is
    positive when
    \begin{align}
        \chi^x_g & \in\iota(\mathbb R_{>0})
        \label{eq:mpug_positive_one_cochain}
    \end{align}
    for every $x\in\mathsf X$ and $g\in G$. A block-dependent scalar
    $2$-cochain $\Omega$ is positive when
    \begin{align}
        \Omega^x_{g,h} & \in\iota(\mathbb R_{>0})
        \label{eq:mpug_positive_two_cochain}
    \end{align}
    for every $x\in\mathsf X$ and $g,h\in G$. Thus positivity means
    membership in the embedded subgroup of strictly positive real numbers,
    not positivity of the real part.
\end{definition}

\begin{definition}[Hat operation on block-dependent cochains]
    \label{def:mpug_generalized_hat}
    \lean{TNLean.Algebra.ActionTensorGauge.hat,
        TNLean.Algebra.LSymbol.hat}
    \leanok
    \uses{def:mpug_l_symbols, def:mpug_l_gauge}
    For block-dependent scalar cochains, define
    \begin{align}
        \widehat\chi^{\,x}_g
         & =\chi^{\,g x}_{g^{-1}},
        \label{eq:mpug_generalized_hat_one}
    \end{align}
    and
    \begin{align}
        \widehat\Omega^{\,x}_{g,h}
         & =\Omega^{\,(gh)x}_{h^{-1},g^{-1}}.
        \label{eq:mpug_generalized_hat_two}
    \end{align}
    These are the one-index and two-index cases of the hat operation in
    \cite[line~5912]{FrancoRubio2025MPUGauging}.
\end{definition}

\begin{theorem}[Hat involution and generalized coboundaries]
    \label{thm:mpug_generalized_hat_coboundary}
    \lean{TNLean.Algebra.ActionTensorGauge.hat_hat,
        TNLean.Algebra.LSymbol.hat_hat,
        TNLean.Algebra.LSymbol.hat_coboundary}
    \leanok
    \uses{def:mpug_generalized_cocycle, def:mpug_generalized_hat}
    The block-dependent hat operation is an involution:
    \begin{align}
        \widehat{\widehat\chi} & =\chi,
        \label{eq:mpug_generalized_hat_hat_one}
    \end{align}
    and
    \begin{align}
        \widehat{\widehat\Omega} & =\Omega.
        \label{eq:mpug_generalized_hat_hat_two}
    \end{align}
    It commutes with the generalized coboundary:
    \begin{align}
        \widehat{d\chi} & =d\widehat\chi.
        \label{eq:mpug_generalized_hat_coboundary}
    \end{align}
\end{theorem}

\begin{proof}\leanok
    \uses{def:mpug_generalized_cocycle, def:mpug_generalized_hat}
    Applying (\ref{eq:mpug_generalized_hat_one}) twice and cancelling each
    group element with its inverse gives
    (\ref{eq:mpug_generalized_hat_hat_one}), and applying
    (\ref{eq:mpug_generalized_hat_two}) twice gives
    (\ref{eq:mpug_generalized_hat_hat_two}). For
    (\ref{eq:mpug_generalized_hat_coboundary}), the action law
    $g(h x)=(gh)x$ and $g^{-1}\bigl((gh)x\bigr)=h x$ give
    \begin{align}
        (d\widehat\chi)^x_{g,h}
         & =\frac{\widehat\chi^{\,h x}_g\widehat\chi^{\,x}_h}
        {\widehat\chi^{\,x}_{gh}}
        =\frac{\chi^{(gh)x}_{g^{-1}}\chi^{h x}_{h^{-1}}}
        {\chi^{(gh)x}_{(gh)^{-1}}}
        =(d\chi)^{(gh)x}_{h^{-1},g^{-1}}
        =\widehat{d\chi}^{\,x}_{g,h}.
    \end{align}
    The two sides therefore carry the same three values of $\chi$, reordered
    in $\C^\times$.
\end{proof}

\begin{definition}[Regular-action matrices]
    \label{def:mpug_generalized_regular_matrices}
    \lean{TNLean.Algebra.LSymbol.regularMatrix}
    \leanok
    \uses{def:mpug_l_symbols}
    Let $\Omega\colon\mathsf X\times G^2\to\C^\times$ be an $L$-symbol.
    On the group algebra $\C G$, define $X^x_g$ by
    \begin{align}
        X^x_g\ket{k}
         & =\Omega^{k^{-1}x}_{g,k}\ket{gk}.
        \label{eq:mpug_generalized_regular_matrix}
    \end{align}
    This is the regular-action matrix in
    \cite[Lemma~\texttt{lemma:G\_cocycle}]{FrancoRubio2025MPUGauging}.
    % Source anchor: arXiv:2502.20257, lines 5938--5942.
\end{definition}

\begin{theorem}[Generalized projective multiplication law]
    \label{thm:mpug_generalized_regular_multiplication}
    \lean{TNLean.Algebra.LSymbol.regularMatrix_mul}
    \leanok
    \uses{def:mpug_generalized_cocycle,
        def:mpug_generalized_regular_matrices}
    Let $G$ be finite. If $\Omega$
    satisfies~(\ref{eq:mpug_generalized_cocycle}), then
    \begin{align}
        X^{h x}_gX^x_h
         & =\Omega^x_{g,h}X^x_{gh}.
        \label{eq:mpug_generalized_regular_multiplication}
    \end{align}
    % Source anchor: arXiv:2502.20257, lemma:G_cocycle, lines 5931--5943,
    %   which states this multiplication law "for G being a finite group".
    %   The blueprint statement previously omitted that hypothesis, which the
    %   formal statement carries.
\end{theorem}

\begin{proof}\leanok
    \uses{def:mpug_generalized_cocycle,
        def:mpug_generalized_regular_matrices}
    Apply the left side to $\ket{k}$. Its coefficient is
    \begin{align}
        \Omega^{(hk)^{-1}h x}_{g,hk}\Omega^{k^{-1}x}_{h,k}
         & =\Omega^{k^{-1}x}_{g,hk}\Omega^{k^{-1}x}_{h,k} \\
         & =\Omega^x_{g,h}\Omega^{k^{-1}x}_{gh,k},
    \end{align}
    where the last equality is
    (\ref{eq:mpug_generalized_cocycle}) at the block $k^{-1}x$.
    Formula~(\ref{eq:mpug_generalized_regular_matrix}) now gives
    (\ref{eq:mpug_generalized_regular_multiplication}).
\end{proof}

\begin{theorem}[Finite-group power trivialization]
    \label{thm:mpug_generalized_cocycle_power}
    \lean{TNLean.Algebra.LSymbol.regularMatrix_eq_permutation_mul_diagonal,
        TNLean.Algebra.LSymbol.regularDeterminant_ne_zero,
        TNLean.Algebra.LSymbol.determinantCochain,
        TNLean.Algebra.LSymbol.pow_card_eq_coboundary_apply,
        TNLean.Algebra.LSymbol.pow_card_eq_coboundary}
    \leanok
    \uses{def:mpug_generalized_cocycle,
        def:mpug_generalized_regular_matrices}
    Let $G$ be finite and let $\Omega$ be a generalized $2$-cocycle. Then
    $\Omega^{|G|}$ is a generalized coboundary. More precisely, with
    \begin{align}
        \chi^x_g
         & =\det X^x_g,
        \label{eq:mpug_generalized_determinant_cochain}
    \end{align}
    one has
    \begin{align}
        \left(\Omega^x_{g,h}\right)^{|G|}
         & =\frac{\det X^{h x}_g\det X^x_h}{\det X^x_{gh}}
        =(d\chi)^x_{g,h}.
        \label{eq:mpug_generalized_power_coboundary}
    \end{align}
    This is
    \cite[Lemma~\texttt{lemma:G\_cocycle}]{FrancoRubio2025MPUGauging}. The
    second numerator determinant in
    (\ref{eq:mpug_generalized_power_coboundary}) carries the index $h$; this
    local correction of the printed formula is recorded in
    \cite{gap:fbc25_generalized_cocycle_determinant_typo}.

    % Local fix: docs/paper-gaps/fbc25_generalized_cocycle_determinant_typo.tex.
    % Source anchor: arXiv:2502.20257, lemma:G_cocycle, lines 5931--5948.
\end{theorem}

\begin{proof}\leanok
    \uses{def:mpug_generalized_regular_matrices,
        thm:mpug_generalized_regular_multiplication}
    Let $P_g$ be the regular-action permutation matrix defined by
    $P_g\ket{k}=\ket{gk}$, and set
    $D^x_g=\operatorname{diag}(\Omega^{k^{-1}x}_{g,k})_{k\in G}$.
    Then
    \begin{align}
        X^x_g
         & =P_gD^x_g.
    \end{align}
    Since $\det P_g\in\{1,-1\}$ and every diagonal entry of $D^x_g$
    is nonzero, $\det X^x_g\ne0$. Thus
    (\ref{eq:mpug_generalized_determinant_cochain}) takes values in
    $\C^\times$. Taking determinants in
    (\ref{eq:mpug_generalized_regular_multiplication}), using
    multiplicativity of the determinant on the left and
    $\det(cX)=c^{|G|}\det X$ for the $|G|\times|G|$ matrix $X^x_{gh}$ on the
    right, gives
    \begin{align}
        \det X^{h x}_g\det X^x_h
         & =\left(\Omega^x_{g,h}\right)^{|G|}\det X^x_{gh}.
        \label{eq:mpug_generalized_determinant_product}
    \end{align}
    Dividing (\ref{eq:mpug_generalized_determinant_product}) by
    $\det X^x_{gh}$ gives
    (\ref{eq:mpug_generalized_power_coboundary}).
\end{proof}

\begin{theorem}[Positive generalized cocycles are coboundaries]
    \label{thm:mpug_positive_generalized_cocycle}
    \lean{TNLean.Algebra.PositiveUnits.positiveRpow,
        TNLean.Algebra.PositiveUnits.positiveRpow_inv_natCast_pow,
        TNLean.Algebra.ActionTensorGauge.positiveRpow,
        TNLean.Algebra.LSymbol.positivePrimitive,
        TNLean.Algebra.LSymbol.isPositiveValued_positivePrimitive,
        TNLean.Algebra.LSymbol.exists_positive_coboundary}
    \leanok
    \uses{def:mpug_generalized_cocycle,
        def:mpug_positive_generalized_cochains,
        thm:mpug_generalized_cocycle_power}
    Let $G$ be finite and let $G$ act on a set $\mathsf X$. If $\Omega$ is a
    positive generalized $2$-cocycle, then there is a positive
    block-dependent scalar $1$-cochain $\chi_0$ such that
    \begin{align}
        \Omega^x_{g,h}
         & =(d\chi_0)^x_{g,h}
        =\frac{\chi^{\,h x}_{0,g}\chi^x_{0,h}}
        {\chi^x_{0,gh}}
        \label{eq:mpug_positive_coboundary}
    \end{align}
    for all $x\in\mathsf X$ and $g,h\in G$.
    This is the first assertion of
    \cite[Corollary at lines~5950--5960]{FrancoRubio2025MPUGauging}.
    The source writes $\chi_g$ here, whereas
    Theorem~\ref{thm:mpug_generalized_cocycle_power} constructs $\chi^x_g$
    through (\ref{eq:mpug_generalized_determinant_cochain}); the block label
    $x$ is therefore retained, as recorded in
    \cite{gap:fbc25_positive_generalized_cocycle_block_index}.
    % Local correction:
    % docs/paper-gaps/fbc25_positive_generalized_cocycle_block_index.tex.
    % Source anchor: arXiv:2502.20257, cor:positive_trivial, lines 5950--5956.
\end{theorem}

\begin{proof}\leanok
    \uses{def:mpug_positive_generalized_cochains,
        thm:mpug_generalized_cocycle_power}
    Write $\delta^x_g=\det X^x_g$ for the determinant cochain from
    Theorem~\ref{thm:mpug_generalized_cocycle_power}, and set
    \begin{align}
        \chi^x_{0,g}
         & =\lvert\delta^x_g\rvert^{1/\lvert G\rvert}.
        \label{eq:mpug_positive_determinant_root}
    \end{align}
    The regular-action permutation can contribute a sign to $\delta^x_g$.
    Its norm is therefore taken before the positive $\lvert G\rvert$-th root,
    as recorded in
    \cite{gap:fbc25_positive_generalized_cocycle_determinant_sign}.
    Since $\Omega$ is positive and the complex norm is multiplicative, taking
    norms in (\ref{eq:mpug_generalized_power_coboundary}) gives
    \begin{align}
        \Omega^{\lvert G\rvert} & =d\lvert\delta\rvert.
        \label{eq:mpug_positive_norm_power}
    \end{align}
    Positive real powers are multiplicative, so
    (\ref{eq:mpug_positive_determinant_root}) and
    (\ref{eq:mpug_positive_norm_power}) give
    \begin{align}
        (d\chi_0)^{\lvert G\rvert}
         & =d\lvert\delta\rvert=\Omega^{\lvert G\rvert},
        \label{eq:mpug_positive_root_power}
    \end{align}
    and the $\lvert G\rvert$-th power is injective on $\mathbb R_{>0}$. Hence
    $d\chi_0=\Omega$.
    % Consolidated clarification:
    % docs/paper-gaps/fbc25_positive_generalized_cocycle_determinant_sign.tex.
\end{proof}

\begin{theorem}[Hat-compatible positive primitive]
    \label{thm:mpug_positive_generalized_cocycle_hat}
    \lean{TNLean.Algebra.ActionTensorGauge.coboundary_div,
        TNLean.Algebra.ActionTensorGauge.coboundary_positiveRpow,
        TNLean.Algebra.LSymbol.exists_positive_coboundary_hat_mul_eq_one}
    \leanok
    \uses{def:mpug_generalized_hat,
        def:mpug_positive_generalized_cochains,
        thm:mpug_generalized_hat_coboundary,
        thm:mpug_positive_generalized_cocycle}
    Under the hypotheses of
    Theorem~\ref{thm:mpug_positive_generalized_cocycle}, suppose additionally
    that
    \begin{align}
        \Omega^x_{g,h}\widehat\Omega^{\,x}_{g,h}
         & =1
        \label{eq:mpug_positive_hat_hypothesis}
    \end{align}
    for every $x\in\mathsf X$ and $g,h\in G$. Then the positive cochain can
    be chosen so that
    \begin{align}
        \Omega & =d\chi,
        \label{eq:mpug_positive_hat_coboundary}
    \end{align}
    and
    \begin{align}
        \chi^x_g\widehat\chi^{\,x}_g
         & =1.
        \label{eq:mpug_positive_hat_witness_identity}
    \end{align}
    This is the second assertion of
    \cite[Corollary at lines~5950--5960]{FrancoRubio2025MPUGauging}.
    % Source anchor: arXiv:2502.20257, cor:positive_trivial, lines 5955--5960.
\end{theorem}

\begin{proof}\leanok
    \uses{thm:mpug_generalized_hat_coboundary,
        thm:mpug_positive_generalized_cocycle}
    Let $\chi_0$ satisfy $d\chi_0=\Omega$ as in
    Theorem~\ref{thm:mpug_positive_generalized_cocycle}. Define
    \begin{align}
        \chi^x_g
         & =\sqrt{\frac{\chi^x_{0,g}}
            {\widehat\chi^{\,x}_{0,g}}}
        =\sqrt{\frac{\chi^x_{0,g}}
            {\chi^{\,g x}_{0,g^{-1}}}}.
        \label{eq:mpug_positive_hat_witness}
    \end{align}
    All values are positive, so the positive square roots are unambiguous.
    Using (\ref{eq:mpug_generalized_hat_coboundary}) and the hypothesis
    (\ref{eq:mpug_positive_hat_hypothesis}) gives
    \begin{align}
        d\chi
         & =\sqrt{d\chi_0/d\widehat\chi_0}
        =\sqrt{\Omega/\widehat\Omega}
        =\sqrt{\Omega^2}
        =\Omega.
    \end{align}
    The same substitution in
    (\ref{eq:mpug_positive_hat_witness}) gives
    \begin{align}
        \widehat\chi
         & =\sqrt{\widehat\chi_0/\chi_0}
        =\chi^{-1},
    \end{align}
    which proves
    (\ref{eq:mpug_positive_hat_witness_identity}).
\end{proof}

\begin{definition}[Inversion cochains]
    \label{def:mpug_inversion_cochains}
    \lean{TNLean.Algebra.ScalarOneCochain,
        TNLean.Algebra.ScalarOneCochain.coboundary,
        TNLean.Algebra.ScalarCocycle.hat,
        TNLean.Algebra.ScalarThreeCochain.hat,
        TNLean.Algebra.ScalarThreeCochain.Xi,
        TNLean.Algebra.ScalarThreeCochain.sigma}
    \leanok
    \uses{def:mpug_scalar_cochains}
    These conventions and inversion cochains follow
    \cite[lines~5910--5916 and 5967--5973]{FrancoRubio2025MPUGauging}.
    For scalar cochains with one, two, and three arguments, set
    \begin{align}
        (d\sigma)(g,h) & =\frac{\sigma(g)\sigma(h)}{\sigma(gh)}.
        \label{eq:mpug_one_coboundary}
    \end{align}
    The hat reverses the arguments and replaces each group element by its
    inverse:
    \begin{align}
        \widehat\beta(g,h) & =\beta(h^{-1},g^{-1}).
        \label{eq:mpug_two_hat}
    \end{align}
    \begin{align}
        \widehat\omega(g,h,k) & =\omega(k^{-1},h^{-1},g^{-1}).
        \label{eq:mpug_three_hat}
    \end{align}
    Define the inversion cochains by
    \begin{align}
        \Xi(g,h) & =
        \frac{\omega(h^{-1},g^{-1},g)}
        {\omega(h^{-1}g^{-1},g,h)}.
        \label{eq:mpug_intro_Xi}
    \end{align}
    \begin{align}
        \sigma(g) & =\omega(g,g^{-1},g).
        \label{eq:mpug_intro_sigma}
    \end{align}
    % Source anchors: arXiv:2502.20257, lines 5910--5916 and eq:introXi.
\end{definition}

\begin{theorem}[The hat operation is an involution]
    \label{thm:mpug_scalar_hat_involution}
    \lean{TNLean.Algebra.ScalarCocycle.hat_hat,
        TNLean.Algebra.ScalarThreeCochain.hat_hat}
    \leanok
    \uses{def:mpug_inversion_cochains}
    Let $G$ be an arbitrary group. Every scalar $2$-cochain $\beta$ and every
    scalar $3$-cochain $\omega$ satisfy
    \begin{align}
        \widehat{\widehat\beta} & =\beta,
        \label{eq:mpug_two_hat_hat}
    \end{align}
    \begin{align}
        \widehat{\widehat\omega} & =\omega.
        \label{eq:mpug_three_hat_hat}
    \end{align}
    % Source anchor: arXiv:2502.20257, line 5912, where the hat operation is
    %   introduced as an involution for cochains of any arity.
\end{theorem}

\begin{proof}\leanok
    \uses{def:mpug_inversion_cochains}
    Reversing the arguments of (\ref{eq:mpug_two_hat}) and inverting each of
    them twice restores the original arguments, and likewise for
    (\ref{eq:mpug_three_hat}).
\end{proof}

\begin{theorem}[Two inversion-cochain formulas]
    \label{thm:mpug_inversion_cochain_formulas}
    \lean{TNLean.Algebra.ScalarThreeCochain.Xi_eq_second_formula,
        TNLean.Algebra.ScalarThreeCochain.sigma_inv}
    \leanok
    \uses{def:mpug_scalar_cochains, def:mpug_three_cocycle,
        def:mpug_inversion_cochains}
    These formulas are part of the cocycle calculation in
    \cite[lines~5967--5979]{FrancoRubio2025MPUGauging}.
    Let $G$ be an arbitrary group and let $\omega$ be a normalized scalar
    $3$-cocycle. Then
    \begin{align}
        \Xi(g,h) & =
        \frac{\omega(h^{-1},g^{-1},g)}
        {\omega(h^{-1}g^{-1},g,h)}
        =
        \frac{\omega(h^{-1},g^{-1},gh)}
        {\omega(g^{-1},g,h)}.
        \label{eq:mpug_Xi_two_formulas}
    \end{align}
    The inversion one-cochain satisfies
    \begin{align}
        \sigma(g^{-1}) & =\sigma(g)^{-1}.
        \label{eq:mpug_sigma_inverse}
    \end{align}
    % Source anchor: arXiv:2502.20257, eq:introXi and lemma:omegaXi.
\end{theorem}

\begin{proof}\leanok
    \uses{def:mpug_scalar_cochains, def:mpug_three_cocycle,
        def:mpug_inversion_cochains}
    Apply (\ref{eq:mpug_three_cocycle}) to $(h^{-1},g^{-1},g,h)$ and use
    (\ref{eq:mpug_three_normalized}); this gives the second quotient in
    (\ref{eq:mpug_Xi_two_formulas}). Applying
    (\ref{eq:mpug_three_cocycle}) to
    $(g,g^{-1},g,g^{-1})$ gives
    $1=\sigma(g)\sigma(g^{-1})$, hence
    (\ref{eq:mpug_sigma_inverse}).
\end{proof}

\begin{theorem}[Inversion identities for a normalized three-cocycle]
    \label{thm:mpug_omega_Xi}
    \lean{TNLean.Algebra.ScalarThreeCochain.hat_eq_mul_coboundary_Xi,
        TNLean.Algebra.ScalarThreeCochain.Xi_eq_hat_mul_coboundary_sigma}
    \leanok
    \uses{def:mpug_scalar_cochains, def:mpug_three_cocycle,
        def:mpug_fusion_gauge, def:mpug_inversion_cochains}
    These inversion identities and their cocycle substitutions are given in
    \cite[lines~5967--5979]{FrancoRubio2025MPUGauging}.
    Let $\omega$ be a normalized scalar $3$-cocycle. Then
    \begin{align}
        \widehat\omega & =\omega\,d\Xi.
        \label{eq:mpug_hat_omega}
    \end{align}
    The inversion two-cochain also satisfies
    \begin{align}
        \Xi & =\widehat\Xi\,d\sigma.
        \label{eq:mpug_Xi_hat_sigma}
    \end{align}
    % Source anchor: arXiv:2502.20257, lemma:omegaXi, lines 5967--5979.
\end{theorem}

\begin{proof}\leanok
    \uses{def:mpug_three_cocycle, def:mpug_fusion_gauge,
        def:mpug_inversion_cochains, thm:mpug_inversion_cochain_formulas}
    For (\ref{eq:mpug_hat_omega}), multiply the cocycle identities at
    \begin{align}
         & (k^{-1},h^{-1},g^{-1},g),\quad
        (k^{-1},h^{-1}g^{-1},g,h),\quad
        (k^{-1}h^{-1}g^{-1},g,h,k),
    \end{align}
    and cancel the normalized factors. For (\ref{eq:mpug_Xi_hat_sigma}), use
    the second expression in (\ref{eq:mpug_Xi_two_formulas}) for
    $\widehat\Xi$, then combine the identities at
    \begin{align}
         & (h^{-1}g^{-1},g,h,h^{-1}g^{-1}),\quad
        (h^{-1},h,h^{-1},g^{-1}),\quad
        (h^{-1},g^{-1},g,g^{-1}),\quad
        (g,g^{-1},g,g^{-1}).
    \end{align}
    The last identity at $g$, $h$, and $gh$ supplies
    (\ref{eq:mpug_sigma_inverse}) in the three required places.
\end{proof}

Write $\Z_2=\{e,s\}$, where $s^2=e$. The following statements concern only
the scalar data and their gauge classes.

\begin{theorem}[Normalized $L$-symbol ratio]
    \label{thm:mpug_sigma_l_ratio}
    \lean{TNLean.Algebra.LSymbol.sigma_eq_div}
    \leanok
    \uses{def:mpug_l_symbols, def:mpug_l_normalized,
        def:mpug_inversion_cochains}
    Let $G$ act on a set $\mathsf X$. Suppose that a normalized $L$-symbol
    $L$ is compatible with a scalar $3$-cochain $\omega$. Then, for every
    $x\in\mathsf X$ and $g\in G$,
    \begin{align}
        \sigma(g)
         & =\frac{L^x_{g^{-1},g}}{L^{g x}_{g,g^{-1}}}.
        \label{eq:mpug_sigma_l_ratio}
    \end{align}
    This is the scalar identity in
    \cite[lines~5327--5329]{FrancoRubio2025MPUGauging}.
    % Source anchor: arXiv:2502.20257, lines 5327--5329.
\end{theorem}

\begin{proof}\leanok
    \uses{def:mpug_l_symbols, def:mpug_l_normalized,
        def:mpug_inversion_cochains}
    Apply the compatibility relation
    (\ref{eq:mpug_l_compatibility}) to $(g,g^{-1},g)$. The two factors with an
    identity argument are $1$ by normalization, leaving
    (\ref{eq:mpug_sigma_l_ratio}).
\end{proof}

\begin{theorem}[Trivial scalar gauge class for $\Z_2$]
    \label{thm:mpug_z2_trivial_scalar_class}
    \lean{TNLean.Algebra.ScalarThreeCochain.isTrivialGaugeClass_iff_sigma_generator_eq_one}
    \leanok
    \uses{def:mpug_scalar_cochains, def:mpug_fusion_gauge,
        def:mpug_trivial_gauge_class, def:mpug_inversion_cochains}
    Let $\omega$ be a normalized scalar $3$-cochain on $\Z_2$. Its gauge
    class is trivial if and only if
    \begin{align}
        \sigma(s) & =1.
        \label{eq:mpug_z2_trivial_sigma}
    \end{align}
    This strengthens the scalar step in
    \cite[lines~5339--5341]{FrancoRubio2025MPUGauging}: only normalization,
    rather than the cocycle identity, is needed for this equivalence.
\end{theorem}

\begin{proof}\leanok
    \uses{def:mpug_fusion_gauge,
        thm:mpug_gauge_preserves_normalization,
        thm:mpug_cohomologous_to_equivalence}
    Suppose first that $\omega=d\beta$ for an arbitrary scalar $2$-cochain
    $\beta$. Although $\beta$ need not be normalized, the normalization of
    $d\beta$ implies that there is a constant $c\in\C^\times$ such that
    $\beta(e,h)=\beta(h,e)=c$ for every $h\in\Z_2$. Since $s^2=e$,
    \begin{align}
        (d\beta)(s,s,s)
         & =\frac{\beta(s,e)\beta(s,s)}
        {\beta(s,s)\beta(e,s)}=\frac{c}{c}=1.
        \label{eq:mpug_z2_coboundary_generator}
    \end{align}
    Thus every scalar $2$-cochain whose coboundary is normalized contributes
    the factor $1$ at $(s,s,s)$. In particular, an arbitrary gauge between
    normalized scalar $3$-cochains cannot change this value, and
    $\sigma(s)=1$. Conversely, normalization determines every
    value of $\omega$ except $\omega(s,s,s)=\sigma(s)$. Hence
    (\ref{eq:mpug_z2_trivial_sigma}) makes $\omega$ the constant cochain $1$.
\end{proof}

\begin{theorem}[Sign dichotomy for the $\Z_2$ inversion scalar]
    \label{thm:mpug_z2_sigma_sign}
    \lean{TNLean.Algebra.ScalarThreeCochain.sigma_generator_eq_one_or_eq_neg_one}
    \leanok
    \uses{def:mpug_three_cocycle, def:mpug_inversion_cochains}
    Let $\omega$ be a normalized scalar $3$-cocycle on $\Z_2$. Then
    \begin{align}
        \sigma(s) & \in\{1,-1\}.
        \label{eq:mpug_z2_sigma_sign}
    \end{align}
\end{theorem}

\begin{proof}\leanok
    \uses{thm:mpug_inversion_cochain_formulas}
    Since $s^{-1}=s$, formula (\ref{eq:mpug_sigma_inverse}) gives
    $\sigma(s)=\sigma(s)^{-1}$. Hence $\sigma(s)^2=1$, whose two solutions in
    $\C^\times$ are $1$ and $-1$.
\end{proof}

For an element $g$ whose square is the identity, the source condition
$\xi_g(x)\xi_{g^{-1}}(g\mathbin{\cdot}x)=-1$ reduces to
$\xi_g(x)\xi_g(g\mathbin{\cdot}x)=-1$ because $g^{-1}=g$.

\begin{definition}[Anomalous sign for an involution]
    \label{def:mpug_anomalous_involution_sign}
    \lean{TNLean.Algebra.IsComplexUnitSign,
        TNLean.Algebra.HasAnomalousInvolutionSign}
    \leanok
    Let $f\colon\mathsf X\to\mathsf X$. An anomalous sign for $f$ is a
    function $\xi\colon\mathsf X\to\C^\times$ such that
    \begin{align}
        \xi(x) & \in\{1,-1\},
        \label{eq:mpug_anomalous_sign_range}
    \end{align}
    and
    \begin{align}
        \xi(x)\xi(f(x)) & =-1
        \label{eq:mpug_anomalous_sign_pairing}
    \end{align}
    for every $x\in\mathsf X$.
    % Source anchor: arXiv:2502.20257, prop:def_tens_anomalous,
    % eq:defects_Z2b, lines 2979--2981.
\end{definition}

\begin{theorem}[Anomalous signs and fixed points]
    \label{thm:mpug_anomalous_involution_sign_iff}
    \lean{TNLean.Algebra.hasAnomalousInvolutionSign_of_fixedPointFree,
        TNLean.Algebra.not_hasAnomalousInvolutionSign_of_fixedPoint,
        TNLean.Algebra.hasAnomalousInvolutionSign_iff_fixedPointFree,
        TNLean.Algebra.involutive_smul_of_mul_self_eq_one,
        TNLean.Algebra.hasAnomalousInvolutionSign_smul_iff}
    \leanok
    \uses{def:mpug_anomalous_involution_sign}
    Let $f\colon\mathsf X\to\mathsf X$ be an involution. An anomalous sign
    exists if and only if $f$ has no fixed point. In particular, let a group
    act on $\mathsf X$ from the left, and suppose that $g^2=e$. The action
    $x\mapsto g\mathbin{\cdot}x$ admits an anomalous sign if and only if
    $g\mathbin{\cdot}x\neq x$ for every $x\in\mathsf X$.
    This is the sign condition used for anomalous defect tensors in
    \cite[lines~2979--2981]{FrancoRubio2025MPUGauging}.
    % Source anchor: arXiv:2502.20257, prop:def_tens_anomalous,
    % eq:defects_Z2b, lines 2979--2981.
\end{theorem}

\begin{proof}\leanok
    \uses{def:mpug_anomalous_involution_sign}
    Choose a well-order on $\mathsf X$. In each two-element orbit
    $\{x,f(x)\}$, assign $1$ to the smaller element and $-1$ to the larger.
    Involutivity exchanges the two elements, so their signs multiply to
    $-1$. Conversely, at a fixed point the product is the square of either
    $1$ or $-1$, hence equals $1$ rather than $-1$. For the action statement,
    $g\mathbin{\cdot}(g\mathbin{\cdot}x)=(g^2)\mathbin{\cdot}x=x$.
\end{proof}

\begin{theorem}[Nontrivial scalar gauge class for $\Z_2$]
    \label{thm:mpug_z2_nontrivial_scalar_class}
    \lean{TNLean.Algebra.ScalarThreeCochain.not_isTrivialGaugeClass_iff_sigma_generator_eq_neg_one}
    \leanok
    \uses{def:mpug_three_cocycle, def:mpug_trivial_gauge_class,
        def:mpug_inversion_cochains}
    Let $\omega$ be a normalized scalar $3$-cocycle on $\Z_2$. Its gauge
    class is nontrivial if and only if
    \begin{align}
        \sigma(s) & =-1.
        \label{eq:mpug_z2_nontrivial_sigma}
    \end{align}
    This is the scalar cohomology classification used in the proof at
    \cite[lines~5339--5341]{FrancoRubio2025MPUGauging}.
    % Source anchor: arXiv:2502.20257, lines 5339--5341.
\end{theorem}

\begin{proof}\leanok
    \uses{thm:mpug_z2_trivial_scalar_class, thm:mpug_z2_sigma_sign}
    By Theorem~\ref{thm:mpug_z2_trivial_scalar_class}, a nontrivial class has
    $\sigma(s)\neq1$. The dichotomy
    (\ref{eq:mpug_z2_sigma_sign}) then forces $\sigma(s)=-1$, and the converse
    follows from the same two statements.
\end{proof}

These results classify only the normalized scalar $3$-cocycle by $\sigma(s)$.
They do not establish the tensor identity
$T_sT_s^*=\sigma(s)\mathbf 1$, nor the proposition identifying the MPU anomaly
with the nontrivial time-reversal SPT class. Both statements require the tensor
argument in \cite[lines~5207--5338]{FrancoRubio2025MPUGauging}, beyond the
scalar calculation above.

\section{The GHZ--blocked-cluster tensor}
\label{sec:mpug_ghz_blocked_cluster}

\begin{definition}[GHZ--blocked-cluster tensor and its two blocks]
    \label{def:mpug_ghz_blocked_cluster_tensor}
    \lean{MPSTensor.clusterBlocked,
        MPSTensor.ghzSectorTensor,
        MPSTensor.ghzSectorTensor_apply_same,
        MPSTensor.z4z2GHZClusterTensor,
        MPSTensor.z4z2GHZClusterBlock,
        MPSTensor.z4z2GHZClusterBlock_apply,
        MPSTensor.z4z2GHZClusterDirectSum}
    \leanok
    \discussion{7699}
    \uses{def:ghz_tensor, def:cluster_tensor,
        def:mps_independent_tensor_product}
    Let $C=\{C^q\}_{q=0}^3$ be the once-blocked cluster-state tensor. The tensor
    obtained by taking its independent tensor product with the GHZ tensor has
    physical dimension $8$ and bond dimension $4$:
    \begin{align}
        A & =A^{\mathrm{GHZ}}\boxtimes C.
        \label{eq:mpug_ghz_blocked_cluster_tensor}
    \end{align}
    For $x\in\{0,1\}$, let $E_x$ be the one-dimensional GHZ sector supported
    at the physical value $x$. The two bond-dimension-$2$ tensors
    \begin{align}
        A_x & =E_x\boxtimes C
    \end{align}
    satisfy, for $p,x\in\{0,1\}$ and $q\in\{0,1,2,3\}$,
    \begin{align}
        A_x^{(p,q)} & =\delta_{p,x}C^q.
        \label{eq:mpug_ghz_blocked_cluster_blocks}
    \end{align}
    Their canonical direct sum is denoted $A_0\oplus A_1$. This is the tensor
    and block formula in
    \cite[lines~4427--4479]{FrancoRubio2025MPUGauging}.
    % Source anchor: arXiv:2502.20257, example:z4z2, eq:z4z2MPS and
    % lines 4427--4479.
\end{definition}

\begin{theorem}[Injective and locally orthogonal GHZ--cluster blocks]
    \label{thm:mpug_ghz_blocked_cluster_blocks}
    \lean{MPSTensor.ghzSectorTensor_isInjective,
        MPSTensor.z4z2GHZClusterBlock_isInjective,
        MPSTensor.z4z2GHZClusterBlocks_isBNTLocallyOrthogonal,
        MPSTensor.z4z2GHZClusterDirectSum_eq_tensor}
    \leanok
    \discussion{7699}
    \uses{def:mpug_ghz_blocked_cluster_tensor, def:injective,
        def:is_bnt_locally_orthogonal}
    Each tensor $A_x$ is injective. For $x\neq y$ and every
    $X\in\operatorname{Mat}_2(\C)$, their one-site mixed transfer map satisfies
    \begin{align}
        \mathcal E_{x,y}(X)
         & =\sum_{p=0}^1\sum_{q=0}^3
        A_x^{(p,q)}X\bigl(A_y^{(p,q)}\bigr)^\dagger
        =0.
        \label{eq:mpug_ghz_blocked_cluster_local_orthogonality}
    \end{align}
    Thus the two blocks are locally orthogonal, and
    \begin{align}
        A_0\oplus A_1 & =A
        \label{eq:mpug_ghz_blocked_cluster_direct_sum}
    \end{align}
    after the canonical identification of the direct-sum and product bond
    coordinates. The local orthogonality is the observation at
    \cite[line~4481]{FrancoRubio2025MPUGauging}.
    % The theorem does not construct the Z4 x Z2 action, its virtual action
    % tensors, its L-symbols, or its gauging data. The group action and gauging
    % of this example remain issue #7352.
\end{theorem}

\begin{proof}\leanok
    \uses{def:mpug_ghz_blocked_cluster_tensor, def:injective,
        def:is_bnt_locally_orthogonal, thm:mps_word_span_tensor_product,
        thm:cluster_normal}
    The one-dimensional sector $E_x$ is injective, and injectivity is preserved
    by independent tensor products, so $A_x=E_x\boxtimes C$ is injective.
    If $x\neq y$, formula~(\ref{eq:mpug_ghz_blocked_cluster_blocks}) makes
    every summand in
    (\ref{eq:mpug_ghz_blocked_cluster_local_orthogonality}) vanish: a physical
    value $p$ cannot equal both $x$ and $y$. For $p,x,y\in\{0,1\}$,
    $q\in\{0,1,2,3\}$, and $a,b\in\{0,1\}$, the direct-sum and product tensors
    have the same entry in their canonical bond coordinates:
    \begin{align}
        (A_0\oplus A_1)^{(p,q)}_{(x,a),(y,b)}
         & =A^{(p,q)}_{(x,a),(y,b)}
        =\delta_{x,y}\delta_{p,x}C^q_{a,b}.
    \end{align}
    Equality of all entries proves
    (\ref{eq:mpug_ghz_blocked_cluster_direct_sum}).
\end{proof}

\section{The gauge-invariant subspace of the CZX circuit tuple}
\label{sec:mpug_czx_gauge_invariant}

The CZX model of \cite[lines~4503--5183]{FrancoRubio2025MPUGauging} is the
anomalous $\Z_2=\{e,g\}$ matrix product unitary acting on the GHZ state.
After blocking, one site carries two matter qubits, and the fusion operators on
two neighboring blocked sites are $\lambda^R_{e,e}=\lambda^R_{e,g}=\Id$,
$\lambda^R_{g,e}=w_R$, together with the modified fusion operator
$\tilde\lambda^R_{g,g}=\lambda^R_{g,g}Z_1$ introduced at the end of that
passage. The question of \cite[lines~5198--5204]{FrancoRubio2025MPUGauging}
asks whether the common $+1$ eigenspace of the resulting Gauss projectors is
bounded or grows with the system size. This section answers it for the
displayed circuit tuple $(\Id,\Id,w,\tilde\lambda)$: the dimension is $2^N$ on
$N\geq3$ blocked sites. Nothing here identifies the displayed tuple
with a member of the full physical completion class of the CZX defect data;
that identification requires the three defect domains and maps not supplied
by the displayed matrices.

\subsection{Monomial operators and their common fixed space}

\begin{definition}[Monomial matrix]
    \label{def:mpug_monomial_matrix}
    \lean{Matrix.monomial,
        Matrix.monomial_apply,
        Matrix.monomial_apply_perm,
        Matrix.monomial_apply_of_ne}
    \leanok
    Let $\sigma$ be a permutation of a finite index set $S$ and let
    $\varphi\colon S\to\C$. The monomial matrix $M(\sigma,\varphi)$ is the
    operator on $\C^S$ with
    \begin{align}
        M(\sigma,\varphi)\ket{s} & =\varphi(s)\ket{\sigma(s)}.
        \label{eq:mpug_monomial_matrix}
    \end{align}
    Its entry in row $\sigma(s)$ and column $s$ is $\varphi(s)$, and every
    other entry vanishes.
\end{definition}

\begin{theorem}[Calculus of monomial matrices]
    \label{thm:mpug_monomial_matrix_laws}
    \lean{Matrix.monomial_mulVec,
        Matrix.monomial_mulVec_eq_self_iff,
        Matrix.monomial_mul_monomial,
        Matrix.conjTranspose_monomial,
        Matrix.smul_monomial,
        Matrix.monomial_one,
        Matrix.reindex_monomial,
        Matrix.monomial_mem_unitaryGroup}
    \leanok
    \uses{def:mpug_monomial_matrix}
    A vector $v\in\C^S$ is fixed by $M(\sigma,\varphi)$ if and only if
    $v_{\sigma(s)}=\varphi(s)v_s$ for every $s\in S$. Products, adjoints,
    scalar multiples, and simultaneous relabellings of monomial matrices are
    monomial:
    \begin{align}
        M(\sigma,\varphi)M(\tau,\psi)
         & =M\bigl(\sigma\tau,\ s\mapsto\varphi(\tau(s))\psi(s)\bigr),
        \label{eq:mpug_monomial_mul}                                               \\
        M(\sigma,\varphi)^\dagger
         & =M\bigl(\sigma^{-1},\ t\mapsto\overline{\varphi(\sigma^{-1}(t))}\bigr).
        \label{eq:mpug_monomial_adjoint}
    \end{align}
    The identity is $M(\id,1)$, and $M(\sigma,\varphi)$ is unitary when
    $|\varphi(s)|=1$ for every $s$.
\end{theorem}

\begin{proof}\leanok
    \uses{def:mpug_monomial_matrix}
    The action of $M(\sigma,\varphi)$ on a vector is
    $(M(\sigma,\varphi)v)_t=\varphi(\sigma^{-1}(t))v_{\sigma^{-1}(t)}$, which
    gives the fixed-vector criterion after the substitution $t=\sigma(s)$.
    The remaining identities are entrywise: in a product only the column
    $\tau(s)$ of the first factor meets the row $\tau(s)$ of the second, the
    adjoint transposes the single nonzero entry of each column, and a
    relabelling conjugates the permutation. For unitarity,
    (\ref{eq:mpug_monomial_mul}) and (\ref{eq:mpug_monomial_adjoint}) give
    $M(\sigma,\varphi)^\dagger M(\sigma,\varphi)=M(\id,|\varphi|^2)=\Id$.
\end{proof}

\begin{theorem}[Flat phases on the hypercube have a potential]
    \label{thm:mpug_hypercube_potential}
    \lean{TNLean.Algebra.hypercube_induction,
        TNLean.Algebra.exists_potential_of_flat}
    \leanok
    Let $M$ be a commutative group and let
    $\varphi_j\colon\mathbb F_2^r\to M$ for $0\leq j<r$ be phases attached
    to the edges $(\gamma,\gamma+e_j)$ of the hypercube. Assume that
    \begin{align}
        \varphi_j(\gamma)\varphi_j(\gamma+e_j)
         & =1,
        \label{eq:mpug_hypercube_involutive}       \\
        \varphi_j(\gamma)\varphi_k(\gamma+e_j)
         & =\varphi_k(\gamma)\varphi_j(\gamma+e_k)
        \label{eq:mpug_hypercube_flat}
    \end{align}
    for all $j,k$ and $\gamma$. Then there is $\Phi\colon\mathbb F_2^r\to M$
    with $\Phi(0)=1$ and $\Phi(\gamma+e_j)=\varphi_j(\gamma)\Phi(\gamma)$ for
    every $j$ and $\gamma$. In addition, a property of vectors of
    $\mathbb F_2^r$ that holds at $0$ and is preserved by every coordinate
    flip holds everywhere.
\end{theorem}

\begin{proof}\leanok
    Connectivity: every $\gamma$ is the sum of the $e_j$ over its support,
    so induction on the support carries the property from $0$ to $\gamma$.

    The potential is constructed by induction on $r$. Write
    $\gamma=(a,\gamma')$ with $a\in\mathbb F_2$ and $\gamma'\in\mathbb F_2^{r-1}$.
    The phases $\gamma'\mapsto\varphi_{j+1}(0,\gamma')$ for $0\leq j<r-1$
    satisfy (\ref{eq:mpug_hypercube_involutive}) and
    (\ref{eq:mpug_hypercube_flat}) on $\mathbb F_2^{r-1}$, so the induction
    hypothesis gives a potential $\Phi'$ for them. Define
    \begin{align}
        \Phi(0,\gamma') & =\Phi'(\gamma'),
        \qquad
        \Phi(1,\gamma')=\varphi_0(0,\gamma')\Phi'(\gamma').
        \label{eq:mpug_hypercube_potential_step}
    \end{align}
    The edge relation in direction $0$ at $(0,\gamma')$ is the definition,
    and at $(1,\gamma')$ it follows from
    (\ref{eq:mpug_hypercube_involutive}). The edge relation in a direction
    $j+1$ at $(0,\gamma')$ is the relation for $\Phi'$, and at $(1,\gamma')$
    it is the relation for $\Phi'$ combined with
    (\ref{eq:mpug_hypercube_flat}) for the square spanned by the directions
    $0$ and $j+1$ at $(0,\gamma')$.
\end{proof}

\begin{definition}[Trivial holonomy on a fiber]
    \label{def:mpug_trivial_holonomy}
    \lean{TNLean.Algebra.fiberFlip,
        TNLean.Algebra.fiberFlip_apply,
        TNLean.Algebra.IsTrivialHolonomy}
    \leanok
    Let $X$ be a finite set and consider the basis $X\times\mathbb F_2^r$.
    For $0\leq j<r$ let $\sigma_j$ be the coordinate flip
    $(x,\gamma)\mapsto(x,\gamma+e_j)$ and let
    $\varphi_j\colon X\times\mathbb F_2^r\to\C$. The phases have trivial
    holonomy on the fiber over $x\in X$ when
    \begin{align}
        \varphi_j(x,\gamma)\varphi_k(x,\gamma+e_j)
         & =\varphi_k(x,\gamma)\varphi_j(x,\gamma+e_k)
        \label{eq:mpug_trivial_holonomy}
    \end{align}
    for all $j,k$ and $\gamma\in\mathbb F_2^r$.
\end{definition}

\begin{theorem}[Monomial fixed-space criterion]
    \label{thm:mpug_monomial_fixed_space_criterion}
    \lean{TNLean.Algebra.reindex_mulVec,
        TNLean.Algebra.commonFixedSubmodule_reindex,
        TNLean.Algebra.finrank_commonFixedSubmodule_reindex,
        TNLean.Algebra.add_self_eq_zero_pi_zmod_two,
        TNLean.Algebra.eq_zero_of_mem_commonFixedSubmodule_monomial_fiberFlip,
        TNLean.Algebra.finrank_commonFixedSubmodule_monomial_fiberFlip}
    \leanok
    \uses{def:mpug_monomial_matrix, def:mpug_trivial_holonomy,
        def:mpug_common_fixed_subspace}
    In the setting of Definition~\ref{def:mpug_trivial_holonomy}, let
    $T_j=M(\sigma_j,\varphi_j)$ and assume that every $T_j$ is an involution,
    that is, $\varphi_j(s)\varphi_j(\sigma_j(s))=1$ for every basis element
    $s$. Then
    \begin{align}
        \dim\mathcal V(T)
         & =\#\{x\in X:\text{the phases have trivial holonomy on the fiber over }x\}.
        \label{eq:mpug_monomial_fixed_space_count}
    \end{align}
    The dimension of a common fixed subspace is also unchanged by a
    simultaneous relabelling of the basis.
\end{theorem}

\begin{proof}\leanok
    \uses{thm:mpug_monomial_matrix_laws, thm:mpug_hypercube_potential,
        thm:mpug_mem_common_fixed_subspace}
    By Theorem~\ref{thm:mpug_monomial_matrix_laws}, a vector $v$ lies in
    $\mathcal V(T)$ exactly when
    $v_{(x,\gamma+e_j)}=\varphi_j(x,\gamma)v_{(x,\gamma)}$ for all $j$, $x$,
    and $\gamma$. Let $X_0$ be the set of fibers with trivial holonomy and
    consider the linear map $E\colon\mathcal V(T)\to\C^{X_0}$,
    $v\mapsto(v_{(x,0)})_{x\in X_0}$.

    $E$ is injective. If $v\in\mathcal V(T)$ vanishes at one point of a
    fiber, the transport relation and the connectivity statement of
    Theorem~\ref{thm:mpug_hypercube_potential} make it vanish on the whole
    fiber. On a fiber with trivial holonomy, $v$ therefore vanishes as soon
    as $v_{(x,0)}=0$. On a fiber where (\ref{eq:mpug_trivial_holonomy})
    fails at some $(j,k,\gamma_0)$, the two paths from $(x,\gamma_0)$ to
    $(x,\gamma_0+e_j+e_k)$ give
    \begin{align}
        v_{(x,\gamma_0+e_j+e_k)} & =\varphi_j(x,\gamma_0)\varphi_k(x,\gamma_0+e_j)\,v_{(x,\gamma_0)}, \\
        v_{(x,\gamma_0+e_j+e_k)} & =\varphi_k(x,\gamma_0)\varphi_j(x,\gamma_0+e_k)\,v_{(x,\gamma_0)}.
    \end{align}
    The two phase products differ, so $v_{(x,\gamma_0)}=0$ and again $v$
    vanishes on the fiber.

    $E$ is surjective. Given $w\in\C^{X_0}$, take on each fiber with trivial
    holonomy the potential $\Phi_x$ of
    Theorem~\ref{thm:mpug_hypercube_potential} for the invertible phases
    $\gamma\mapsto\varphi_j(x,\gamma)$, and set
    $v_{(x,\gamma)}=w_x\Phi_x(\gamma)$ there and $v=0$ on the other fibers.
    The potential relation is the transport relation, so $v\in\mathcal V(T)$,
    and $\Phi_x(0)=1$ gives $Ev=w$.

    Hence $\dim\mathcal V(T)=\dim\C^{X_0}=\#X_0$. For the relabelling
    statement, a relabelled matrix acts by relabelling the vector, acting,
    and relabelling back, so the common fixed subspace is the image of the
    original one under the coordinate relabelling.
\end{proof}

\begin{theorem}[Local Gauss operator of monomial matter unitaries]
    \label{thm:mpug_gauss_operator_monomial}
    \lean{TNLean.Algebra.gaussMonomialPerm,
        TNLean.Algebra.gaussMonomialPerm_apply,
        TNLean.Algebra.gaussMonomialPhase,
        TNLean.Algebra.gaussOperator_eq_monomial}
    \leanok
    \uses{def:mpug_gauss_operator, def:mpug_monomial_matrix}
    Suppose that every matter unitary of the tuple $R$ is monomial,
    $R_{a,b}=M(\sigma_{a,b},\varphi_{a,b})$. Then the local Gauss operator
    of Definition~\ref{def:mpug_gauss_operator} is the monomial operator
    \begin{align}
        \mathcal G_g\ket{i,(a,b)}
         & =\overline{\varphi_{T_g(a,b)}\bigl(\sigma_{T_g(a,b)}^{-1}\sigma_{a,b}(i)\bigr)}\,
        \varphi_{a,b}(i)\,
        \ket{\sigma_{T_g(a,b)}^{-1}\sigma_{a,b}(i),\ T_g(a,b)}.
        \label{eq:mpug_gauss_operator_monomial}
    \end{align}
\end{theorem}

\begin{proof}\leanok
    \uses{thm:mpug_monomial_matrix_laws, def:mpug_gauss_operator}
    The block of $\mathcal G_g$ from $(a,b)$ to $T_g(a,b)$ is
    $R_{T_g(a,b)}^\dagger R_{a,b}$, which is monomial by
    (\ref{eq:mpug_monomial_mul}) and (\ref{eq:mpug_monomial_adjoint}); all
    other blocks vanish.
\end{proof}

\begin{theorem}[Periodic placement of monomial operators]
    \label{thm:mpug_embed_monomial}
    \lean{MPOTensor.windowPerm,
        MPOTensor.windowPerm_apply,
        MPOTensor.embedLocalOperator_monomial,
        MPOTensor.extractWindow_two,
        MPOTensor.replaceWindow_two_apply}
    \leanok
    \uses{def:mpug_monomial_matrix}
    Let $M(\sigma,\varphi)$ be a monomial operator on $L$ consecutive sites.
    Its placement on the periodic window beginning at site $i$ of an $N$-site
    chain is the monomial operator whose permutation replaces the window of a
    configuration $s$ by $\sigma$ applied to that window and whose phase is
    $\varphi$ evaluated on the window of $s$. For $L=2$ and $N\geq2$ the
    window at $j$ is $(s_j,s_{j+1})$, and replacing it changes exactly the
    sites $j$ and $j+1$.
\end{theorem}

\begin{proof}\leanok
    \uses{def:mpug_monomial_matrix}
    The placed operator has entry $(t,s)$ equal to the window entry
    $M(\sigma,\varphi)_{t_W,s_W}$ when $t$ and $s$ agree outside the window
    and $0$ otherwise. Both conditions together say that $t$ is $s$ with its
    window replaced by $\sigma(s_W)$, and the entry is then
    $\varphi(s_W)$. The two-site statements are the cyclic offset
    computations $(j+N-j)\bmod N=0$ and $(j+1+N-j)\bmod N=1$.
\end{proof}

\subsection{The displayed CZX circuit tuple}

\begin{definition}[The displayed CZX circuit tuple]
    \label{def:mpug_czx_circuit_tuple}
    \lean{MPOTensor.CZX.bitFlip,
        MPOTensor.CZX.pauliX,
        MPOTensor.CZX.pauliZ,
        MPOTensor.CZX.controlledZ,
        MPOTensor.CZX.w,
        MPOTensor.CZX.lambda,
        MPOTensor.CZX.tildeLambda,
        MPOTensor.CZX.siteBits,
        MPOTensor.CZX.localBits,
        MPOTensor.CZX.localBits_apply_zero,
        MPOTensor.CZX.localBits_apply_one,
        MPOTensor.CZX.localBits_apply_two,
        MPOTensor.CZX.localBits_apply_three,
        MPOTensor.CZX.matterMatrix,
        MPOTensor.CZX.gen,
        MPOTensor.CZX.circuitTuple,
        MPOTensor.CZX.circuitTuple_one,
        MPOTensor.CZX.circuitTuple_gen_one,
        MPOTensor.CZX.circuitTuple_gen_gen}
    \leanok
    \uses{def:mpug_monomial_matrix}
    Number the four matter qubits of two neighboring blocked sites by
    $0,1,2,3$, so that $(x_0,x_1)$ are the qubits of the first site and
    $(x_2,x_3)$ those of the second, and write
    $\ket{x}=\ket{x_0}\otimes\ket{x_1}\otimes\ket{x_2}\otimes\ket{x_3}$ for
    $x\in\mathbb F_2^4$. The elementary gates act by
    \begin{align}
        X_r\ket{x}              & =\ket{x+e_r},          \\
        Z_r\ket{x}              & =(-1)^{x_r}\ket{x},    \\
        \mathrm{CZ}_{rs}\ket{x} & =(-1)^{x_rx_s}\ket{x},
    \end{align}
    and operator products act from right to left. Define
    \begin{align}
        w             & =(X_0X_1)\,\mathrm{CZ}_{01}\mathrm{CZ}_{12}\,(Z_2Z_3),
        \label{eq:mpug_czx_w}                                                       \\
        \lambda       & =-i\,(X_0X_1X_2Z_3)\,\mathrm{CZ}_{01}\mathrm{CZ}_{12}\,X_2,
        \label{eq:mpug_czx_lambda}                                                  \\
        \tilde\lambda & =\lambda Z_0.
        \label{eq:mpug_czx_tilde_lambda}
    \end{align}
    These are the fusion operators $\lambda^R_{g,e}=w_R$ and
    $\lambda^R_{g,g}$ and the modified fusion operator
    $\tilde\lambda^R_{g,g}=\lambda^R_{g,g}Z_1$ of
    \cite[lines~4503--5183]{FrancoRubio2025MPUGauging} in this qubit
    numbering. The displayed circuit tuple, indexed by the ordered pair of
    gauge labels in $\Z_2$ with $e\leftrightarrow0$ and $g\leftrightarrow1$,
    is
    \begin{align}
        R^{\mathrm{CZX}}
         & =\bigl(U_{0,0},U_{0,1},U_{1,0},U_{1,1}\bigr)
        =\bigl(\Id,\Id,w,\tilde\lambda\bigr).
        \label{eq:mpug_czx_tuple}
    \end{align}
    Each site index in $\{0,1,2,3\}$ is read as the pair of matter bits
    $(m^{(1)},m^{(2)})$ with $m^{(1)}$ the leading binary digit, and the
    four-qubit operators are transported to the two-site matter space by
    this identification.
    % Source anchor: arXiv:2502.20257, lines 4503--5183 (CZX model and modified
    % fusion operator).
\end{definition}

\begin{theorem}[Phase tables of the CZX fusion operators]
    \label{thm:mpug_czx_phase_tables}
    \lean{MPOTensor.CZX.barFlip,
        MPOTensor.CZX.barFlip_apply,
        MPOTensor.CZX.barFlip_mul_barFlip,
        MPOTensor.CZX.barFlip_symm,
        MPOTensor.CZX.eExponent,
        MPOTensor.CZX.hExponent,
        MPOTensor.CZX.fExponent,
        MPOTensor.CZX.neg_one_pow_val_add,
        MPOTensor.CZX.w_eq,
        MPOTensor.CZX.lambda_eq,
        MPOTensor.CZX.tildeLambda_eq,
        MPOTensor.CZX.fExponent_barFlip,
        MPOTensor.CZX.eExponent_barFlip,
        MPOTensor.CZX.matterBarFlip,
        MPOTensor.CZX.matterBarFlip_symm,
        MPOTensor.CZX.localBits_matterBarFlip,
        MPOTensor.CZX.matterMatrix_monomial,
        MPOTensor.CZX.w_mem_unitaryGroup,
        MPOTensor.CZX.tildeLambda_mem_unitaryGroup}
    \leanok
    \uses{def:mpug_czx_circuit_tuple}
    Write $\overline x=(x_0+1,x_1+1,x_2,x_3)$ and
    \begin{align}
        e(x) & =x_0x_1+x_1x_2+x_2+x_3,     \\
        h(x) & =x_0x_1+x_1x_2+x_1+x_3,     \\
        f(x) & =x_0x_1+x_1x_2+x_0+x_1+x_3.
    \end{align}
    Then
    \begin{align}
        w\ket{x}             & =(-1)^{e(x)}\ket{\overline x},
        \label{eq:mpug_czx_w_action}                            \\
        \lambda\ket{x}       & =-i(-1)^{h(x)}\ket{\overline x},
        \label{eq:mpug_czx_lambda_action}                       \\
        \tilde\lambda\ket{x} & =-i(-1)^{f(x)}\ket{\overline x},
        \label{eq:mpug_czx_tilde_lambda_action}
    \end{align}
    the exponents satisfy $f(\overline x)=e(x)+1$ and
    $e(\overline x)=f(x)+1$. After transport to the two-site matter coordinates,
    $w$ and $\tilde\lambda$ are unitary.
\end{theorem}

\begin{proof}\leanok
    \uses{thm:mpug_monomial_matrix_laws, def:mpug_czx_circuit_tuple}
    Each gate is monomial, so (\ref{eq:mpug_monomial_mul}) expresses each
    product as a monomial operator. Reading (\ref{eq:mpug_czx_w}) from right
    to left, the diagonal gates contribute the phase
    $(-1)^{x_0x_1+x_1x_2+x_2+x_3}$ and the two bit flips send $x$ to
    $\overline x$. In (\ref{eq:mpug_czx_lambda}) the first $X_2$ replaces
    $x_2$ by $x_2+1$ in the phase of $\mathrm{CZ}_{12}$, giving
    $x_1(x_2+1)=x_1x_2+x_1$, and the second $X_2$ restores $x_2$, so the
    permutation is again $\overline{\,\cdot\,}$. The factor $Z_0$ in
    (\ref{eq:mpug_czx_tilde_lambda}) adds $x_0$ to the exponent. The
    identities between $e$ and $f$ are checked on the sixteen bit strings.
    After relabelling the computational basis as two four-level matter sites,
    unitarity follows from Theorem~\ref{thm:mpug_monomial_matrix_laws} because
    all phases are unimodular.
\end{proof}

\begin{theorem}[Local monomial action of the CZX Gauss operator]
    \label{thm:mpug_czx_local_monomial_action}
    \lean{MPOTensor.CZX.iPow,
        MPOTensor.CZX.iPow_add,
        MPOTensor.CZX.iPow_natCast,
        MPOTensor.CZX.iPow_zero,
        MPOTensor.CZX.iPow_two,
        MPOTensor.CZX.iPow_three,
        MPOTensor.CZX.iPow_two_mul_val,
        MPOTensor.CZX.iPow_two_mul_cast,
        MPOTensor.CZX.uExponent,
        MPOTensor.CZX.phaseExponent,
        MPOTensor.CZX.phaseExponent_zero_zero,
        MPOTensor.CZX.phaseExponent_zero_one,
        MPOTensor.CZX.phaseExponent_one_zero,
        MPOTensor.CZX.phaseExponent_one_one,
        MPOTensor.CZX.localPerm,
        MPOTensor.CZX.localPhase,
        MPOTensor.CZX.gen_inv,
        MPOTensor.CZX.gen_mul_gen,
        MPOTensor.CZX.toAdd_gen_ne_zero,
        MPOTensor.CZX.gaussOperator_circuitTuple_gen}
    \leanok
    \uses{def:mpug_czx_circuit_tuple, def:mpug_gauss_operator}
    Let $a,a'\in\mathbb F_2$ be the two gauge bits and write
    $s=(x_0,x_1,a,x_2,x_3,a')$ and
    $\sigma(s)=(x_0+1,x_1+1,a+1,x_2,x_3,a'+1)$. The nontrivial local Gauss
    operator obtained from (\ref{eq:mpug_czx_tuple}) and
    (\ref{eq:mpug_gauss_operator}) is the monomial operator
    \begin{align}
        G\ket{s} & =i^{q(s)}\ket{\sigma(s)}
        \label{eq:mpug_czx_monomial_action}
    \end{align}
    with, for $u(x)=2(x_0x_1+x_1x_2+x_3)\bmod4$, the phase table
    \begin{align}
        q(x_0,x_1,0,x_2,x_3,0) & =3+u(x)+2x_2,      \\
        q(x_0,x_1,0,x_2,x_3,1) & =2+u(x)+2x_0+2x_1, \\
        q(x_0,x_1,1,x_2,x_3,0) & =u(x)+2x_2,        \\
        q(x_0,x_1,1,x_2,x_3,1) & =3+u(x)+2x_0+2x_1
        \label{eq:mpug_czx_phase_table}
    \end{align}
    modulo $4$.
\end{theorem}

\begin{proof}\leanok
    \uses{thm:mpug_gauss_operator_monomial, thm:mpug_czx_phase_tables}
    By Theorem~\ref{thm:mpug_gauss_operator_monomial} applied to the monomial
    unitaries of Theorem~\ref{thm:mpug_czx_phase_tables}, the operator
    $\mathcal G_g$ is monomial. The permutation part is
    $\overline{\,\cdot\,}$ on the matter bits in all four gauge sectors,
    since $\overline{\,\cdot\,}$ is an involution, and $T_g$ flips both
    gauge bits. The phases are read off
    (\ref{eq:mpug_czx_w_action})--(\ref{eq:mpug_czx_tilde_lambda_action})
    and the adjoint formula (\ref{eq:mpug_monomial_adjoint}) together with
    $f(\overline x)=e(x)+1$ and $e(\overline x)=f(x)+1$: for example the
    first row is $-i(-1)^{e(x)}=i^{3+2e(x)}=i^{3+u(x)+2x_2}$, and the other
    three rows follow in the same way. The four identities between the phase
    table and $e$, $f$ are checked on the sixteen bit strings.
\end{proof}

\subsection{Orbits, holonomy, and the exact dimension}

\begin{definition}[Chain configurations, bond flips, and labels]
    \label{def:mpug_czx_chain_labels}
    \lean{MPOTensor.CZX.Site,
        MPOTensor.CZX.siteDecode,
        MPOTensor.CZX.siteDecode_equivFin,
        MPOTensor.CZX.chainDecode,
        MPOTensor.CZX.chainDecode_apply,
        MPOTensor.CZX.flipPattern,
        MPOTensor.CZX.chainFlip,
        MPOTensor.CZX.chainFlip_apply,
        MPOTensor.CZX.succ_ne_self,
        MPOTensor.CZX.succ_succ_ne_self,
        MPOTensor.CZX.flipPattern_apply_self,
        MPOTensor.CZX.flipPattern_apply_succ,
        MPOTensor.CZX.flipPattern_apply_of_ne,
        MPOTensor.CZX.bondExponent,
        MPOTensor.CZX.bLabel,
        MPOTensor.CZX.pLabel}
    \leanok
    \uses{thm:mpug_czx_local_monomial_action}
    A computational-basis configuration of the periodic chain of $N$ blocked
    sites is
    \begin{align}
        s & =\bigl(m^{(1)}_k,m^{(2)}_k,g_k\bigr)_{k\in\Z/N\Z}\in S_N:=(\mathbb F_2^3)^N,
        \label{eq:mpug_czx_configuration}
    \end{align}
    read from the enumerated site indices of the placed Gauss projector. The
    bond-$j$ flip $\sigma_j$ adds $1$ to $m^{(1)}_j$, $m^{(2)}_j$, $g_j$, and
    $g_{j+1}$. The bond-$j$ phase exponent $q_j(s)$ is the table
    (\ref{eq:mpug_czx_phase_table}) evaluated at
    $x=(m^{(1)}_j,m^{(2)}_j,m^{(1)}_{j+1},m^{(2)}_{j+1})$ and
    $(a,a')=(g_j,g_{j+1})$. The affine labels are
    \begin{align}
        p_k(s) & =m^{(1)}_k+m^{(2)}_k,
        \label{eq:mpug_czx_p_label}            \\
        b_k(s) & =g_k+m^{(1)}_k+m^{(1)}_{k-1}.
        \label{eq:mpug_czx_b_label}
    \end{align}
\end{definition}

\begin{theorem}[Permutation-orbit classification]
    \label{thm:mpug_czx_orbit_classification}
    \lean{MPOTensor.CZX.fiberEquiv,
        MPOTensor.CZX.fiberEquiv_apply,
        MPOTensor.CZX.pLabel_fiberEquiv,
        MPOTensor.CZX.bLabel_fiberEquiv,
        MPOTensor.CZX.chainFlip_fiberEquiv}
    \leanok
    \uses{def:mpug_czx_chain_labels}
    For $N\geq2$, the map
    \begin{align}
        \bigl((p,b),\gamma\bigr)
         & \longmapsto
        \bigl(m^{(1)}_k,m^{(2)}_k,g_k\bigr)_k
        =\bigl(\gamma_k,\ p_k+\gamma_k,\ b_k+\gamma_k+\gamma_{k-1}\bigr)_k
        \label{eq:mpug_czx_fiber_coordinates}
    \end{align}
    is a bijection $\mathbb F_2^N\times\mathbb F_2^N\times\mathbb F_2^N\to S_N$
    whose inverse reads off $p(s)$, $b(s)$, and $\gamma=m^{(1)}$. In these
    coordinates $\sigma_j$ is $\gamma\mapsto\gamma+e_j$ and leaves $(p,b)$
    unchanged.
\end{theorem}

\begin{proof}\leanok
    \uses{def:mpug_czx_chain_labels}
    Both compositions of (\ref{eq:mpug_czx_fiber_coordinates}) with the
    label map are the identity, by the identity $a+a=0$ in $\mathbb F_2$.
    The flip $\sigma_j$ adds $((1,1),1)$ at site $j$ and $((0,0),1)$ at
    site $j+1$; in fiber coordinates this is $\gamma_j\mapsto\gamma_j+1$
    because $k-1=j$ exactly when $k=j+1$.
\end{proof}

\begin{theorem}[Involution and neighboring holonomy]
    \label{thm:mpug_czx_holonomy}
    \lean{MPOTensor.CZX.bondExponent_add_bondExponent_chainFlip,
        MPOTensor.CZX.bondExponent_holonomy,
        MPOTensor.CZX.bondExponent_chainFlip_of_disjoint,
        MPOTensor.CZX.fiberPhase,
        MPOTensor.CZX.fiberPhase_mul_fiberPhase_flip,
        MPOTensor.CZX.fiberPhase_flip_of_disjoint,
        MPOTensor.CZX.fiberPhase_holonomy,
        MPOTensor.CZX.fiberPhase_ne_zero}
    \leanok
    \uses{def:mpug_czx_chain_labels}
    For every $N\geq3$, every bond $j$, and every configuration $s$,
    \begin{align}
        q_j(s)+q_j(\sigma_js)
         & =0,
        \label{eq:mpug_czx_involution}                \\
        q_j(s)+q_{j+1}(\sigma_js)
         & =q_{j+1}(s)+q_j(\sigma_{j+1}s)+2b_{j+1}(s)
        \label{eq:mpug_czx_holonomy}
    \end{align}
    modulo $4$, and $q_j(\sigma_ks)=q_j(s)$ whenever
    $k\notin\{j-1,j,j+1\}$. For the phases
    $\varphi_j(s)=i^{q_j(s)}$, (\ref{eq:mpug_czx_involution}) gives
    $\varphi_j(s)\varphi_j(\sigma_js)=1$, while
    (\ref{eq:mpug_czx_holonomy}) gives the neighboring square relation with
    factor $(-1)^{b_{j+1}(s)}$.
\end{theorem}

\begin{proof}\leanok
    \uses{def:mpug_czx_chain_labels}
    The exponent $q_j$ depends only on the six bits at the sites $j$ and
    $j+1$, and $\sigma_k$ changes only the sites $k$ and $k+1$; this gives
    the disjoint-bond statement. For (\ref{eq:mpug_czx_involution}) the
    flip $\sigma_j$ changes $x_0$, $x_1$, $a$, and $a'$, and the identity is
    verified on the $64$ values of the six bits. For
    (\ref{eq:mpug_czx_holonomy}) the sites $j$, $j+1$, $j+2$ are distinct
    because $N\geq3$; $\sigma_j$ changes the left gauge bit seen by the bond
    $(j+1,j+2)$, and $\sigma_{j+1}$ changes $x_2$, $x_3$, and $a'$ seen by
    the bond $(j,j+1)$, so the identity is verified on the $512$ values of
    the nine bits at the three sites.
\end{proof}

\begin{theorem}[Fixed spaces on the orbit fibers]
    \label{thm:mpug_czx_good_fibers}
    \lean{MPOTensor.CZX.isTrivialHolonomy_fiberPhase_iff}
    \leanok
    \uses{def:mpug_trivial_holonomy, thm:mpug_czx_orbit_classification}
    Let $N\geq3$. In the fiber coordinates
    (\ref{eq:mpug_czx_fiber_coordinates}), the phases $i^{q_j}$ have trivial
    holonomy on the fiber over $(p,b)$ if and only if $b=0$.
\end{theorem}

\begin{proof}\leanok
    \uses{thm:mpug_czx_holonomy, thm:mpug_czx_orbit_classification}
    If $b=0$, then for $k=j+1$ or $j=k+1$ the square identity is
    (\ref{eq:mpug_czx_holonomy}) with $b_{j+1}=0$, for $k=j$ it is trivial,
    and otherwise the bonds are disjoint and both sides agree. Conversely,
    trivial holonomy at the square spanned by the bonds $k-1$ and $k$ at
    $\gamma=0$ gives $(-1)^{b_k}=1$ by (\ref{eq:mpug_czx_holonomy}), since the
    phases are nonzero, hence $b_k=0$ for every $k$.
\end{proof}

\begin{theorem}[The placed CZX Gauss operators]
    \label{thm:mpug_czx_placed_operators}
    \lean{MPOTensor.CZX.windowCoordinates,
        MPOTensor.CZX.placedGaussOperator,
        MPOTensor.CZX.sum_multiplicative_zmod_two,
        MPOTensor.CZX.gaussProjector_circuitTuple,
        MPOTensor.CZX.placedGaussProjector_circuitTuple,
        MPOTensor.CZX.inv_two_smul_add_mulVec_eq_iff,
        MPOTensor.CZX.commonFixedSubmodule_placedGaussProjector_circuitTuple,
        MPOTensor.CZX.placedGaussOperator_eq_monomial,
        MPOTensor.CZX.siteBits_eq_localBits_zero,
        MPOTensor.CZX.siteBits_eq_localBits_one,
        MPOTensor.CZX.barFlip_apply_zero,
        MPOTensor.CZX.barFlip_apply_one,
        MPOTensor.CZX.barFlip_apply_two,
        MPOTensor.CZX.barFlip_apply_three,
        MPOTensor.CZX.toAdd_mul_gen_inv,
        MPOTensor.CZX.toAdd_gen_mul,
        MPOTensor.CZX.gaussLegAction_fst,
        MPOTensor.CZX.gaussLegAction_snd,
        MPOTensor.CZX.chainDecode_window,
        MPOTensor.CZX.chainDecode_windowPerm,
        MPOTensor.CZX.localPhase_window,
        MPOTensor.CZX.fiberCoordinates,
        MPOTensor.CZX.placedGaussOperator_eq_reindex}
    \leanok
    \uses{def:mpug_placed_gauss_projector, def:mpug_czx_circuit_tuple,
        def:mpug_czx_chain_labels, def:mpug_common_fixed_subspace}
    Let $N\geq2$ and let $G_j$ be the local Gauss operator
    (\ref{eq:mpug_czx_monomial_action}) placed on the bond $(j,j+1)$. For
    $\Z_2$ the placed Gauss projector of
    Definition~\ref{def:mpug_placed_gauss_projector} is
    \begin{align}
        \widetilde P_j & =\tfrac12\bigl(\Id+G_j\bigr),
        \label{eq:mpug_czx_placed_projector}
    \end{align}
    so the common fixed subspace $\mathcal V_N(R^{\mathrm{CZX}})$ of the
    $\widetilde P_j$ is the common fixed subspace of the $G_j$. In the
    coordinates (\ref{eq:mpug_czx_configuration}), $G_j$ is the monomial
    operator $G_j\ket{s}=i^{q_j(s)}\ket{\sigma_js}$, and in the fiber
    coordinates (\ref{eq:mpug_czx_fiber_coordinates}) it is the monomial
    operator with permutation $\gamma\mapsto\gamma+e_j$ and phase
    $i^{q_j}$.
\end{theorem}

\begin{proof}\leanok
    \uses{def:mpug_gauss_projector, thm:mpug_gauss_representation,
        thm:mpug_czx_local_monomial_action, thm:mpug_embed_monomial,
        thm:mpug_czx_orbit_classification, thm:mpug_mem_common_fixed_subspace}
    The average over $\Z_2=\{e,g\}$ is $\frac12(\mathcal G_e+\mathcal G_g)$
    with $\mathcal G_e=\Id$, and relabelling and placement are linear and
    unital, which gives (\ref{eq:mpug_czx_placed_projector}). A vector $v$
    satisfies $\frac12(v+G_jv)=v$ exactly when $G_jv=v$.
    Theorem~\ref{thm:mpug_czx_local_monomial_action} and
    Theorem~\ref{thm:mpug_embed_monomial} express $G_j$ as a monomial
    operator; decoding the two-site window at $j$ into the bits of the sites
    $j$ and $j+1$ identifies its permutation with $\sigma_j$ and its phase
    with $i^{q_j(s)}$. Finally
    Theorem~\ref{thm:mpug_czx_orbit_classification} transports both to the
    fiber coordinates.
\end{proof}

\begin{theorem}[Exact dimension for the displayed CZX circuit tuple]
    \label{thm:mpug_czx_dimension}
    \lean{MPOTensor.CZX.finrank_commonFixedSubmodule_placedGaussProjector_circuitTuple}
    \leanok
    \uses{def:mpug_placed_gauss_projector, def:mpug_czx_circuit_tuple,
        def:mpug_common_fixed_subspace}
    On a periodic chain of $N\geq3$ blocked sites,
    \begin{align}
        \dim\mathcal V_N(R^{\mathrm{CZX}}) & =2^N.
        \label{eq:mpug_czx_dimension}
    \end{align}
    This answers the growth question of
    \cite[lines~5198--5204]{FrancoRubio2025MPUGauging} for the displayed
    circuit tuple; it does not assert that the tuple belongs to the full
    physical completion class of the CZX defect data, and it says nothing
    about the minimum or maximum over completions.
\end{theorem}

\begin{proof}\leanok
    \uses{thm:mpug_czx_placed_operators,
        thm:mpug_monomial_fixed_space_criterion,
        thm:mpug_czx_holonomy, thm:mpug_czx_good_fibers}
    By Theorem~\ref{thm:mpug_czx_placed_operators},
    $\mathcal V_N(R^{\mathrm{CZX}})$ is the common fixed subspace of monomial
    operators which, in fiber coordinates, have permutations
    $\gamma\mapsto\gamma+e_j$ on the fibers over
    $(p,b)\in\mathbb F_2^N\times\mathbb F_2^N$, and by
    (\ref{eq:mpug_czx_involution}) these operators are involutions. The
    monomial fixed-space criterion
    (\ref{eq:mpug_monomial_fixed_space_count}) counts the fibers with trivial
    holonomy, which by Theorem~\ref{thm:mpug_czx_good_fibers} are exactly
    the fibers with $b=0$. There is one such $b$ and there are $2^N$ values
    of $p$, so
    \begin{align}
        \dim\mathcal V_N(R^{\mathrm{CZX}})
         & =\sum_{p\in\mathbb F_2^N}1=2^N.
    \end{align}
\end{proof}
